% 日本語 LuaLaTeX。2026-11-25 / v1.71
% エンジン、出力先、SyncTeX は repo ルートの .latexmkrc で指定する。
\RequirePackage{fix-cm}
\documentclass[a4paper,11pt]{ltjsarticle}
\usepackage{amsmath,amssymb,amsthm,mathtools,mathrsfs}
% rsfs の小さな添字も、代替サイズではなく設計フォントの拡縮で組む。
\DeclareFontFamily{U}{rsfs}{\skewchar\font127}
\DeclareFontShape{U}{rsfs}{m}{n}{<-6> rsfs5 <6-8> rsfs7 <8-> rsfs10}{}
\usepackage[margin=25mm]{geometry}
\usepackage{booktabs,enumitem}
\usepackage[unicode,pdfusetitle,colorlinks=true,
 linkcolor=blue,citecolor=blue,urlcolor=blue]{hyperref}
\numberwithin{equation}{section}
\theoremstyle{plain}
\newtheorem{theorem}{定理}[section]
\newtheorem{proposition}[theorem]{命題}
\newtheorem{lemma}[theorem]{補題}
\newtheorem{corollary}[theorem]{系}
\theoremstyle{definition}
\newtheorem{definition}[theorem]{定義}
\newtheorem{example}[theorem]{例}
\newtheorem{problem}[theorem]{問題}
\theoremstyle{remark}
\newtheorem{remark}[theorem]{注意}
\renewcommand{\proofname}{証明}
\newcommand{\kk}{k}
\newcommand{\OO}{\mathcal O}
\newcommand{\cC}{\mathcal C}
\newcommand{\PP}{\mathbb P}
\newcommand{\HH}{\mathcal H_2}
\newcommand{\HHstack}{\mathfrak H_2}
\DeclareMathOperator{\Hom}{Hom}
\DeclareMathOperator{\End}{End}
\DeclareMathOperator{\Ext}{Ext}
\DeclareMathOperator{\coh}{coh}
\DeclareMathOperator{\qgr}{qgr}
\DeclareMathOperator{\gldim}{gldim}
\DeclareMathOperator{\Aut}{Aut}
\DeclareMathOperator{\Pic}{Pic}
\DeclareMathOperator{\Sym}{Sym}
\DeclareMathOperator{\Gr}{Gr}
\DeclareMathOperator{\GL}{GL}
\DeclareMathOperator{\PGL}{PGL}
\DeclareMathOperator{\rank}{rank}
\DeclareMathOperator{\ch}{ch}
\DeclareMathOperator{\td}{td}
\DeclareMathOperator{\length}{length}
\title{構造対象を指定した非可換二次曲面：\\平方根と二次 AS 正則表示の忘却}
\author{}
\date{2026年11月25日\quad v1.71}

\begin{document}
\maketitle
\begin{abstract}
構造対象 $\OO$ のみを指定した強い非可換二次曲面
$(\cC,\OO)$ を pointed exact 同値で分類する。
$\OO$ は Gabriel--Krull 商で rank-one torsion-free exceptional とし、
反 Serre 関手の exact ample な平方根と、その軌道環の
型 $(1,4,6,4,1)$ の二次 AS 正則代数による正則正規二次商表示を要求する。
平方根・その同型・ambient は標識として固定しない。

体上の全対象は、linear データと、被約・非被約を含む
Van den Bergh 許容四つ組から構成される cubic 領域、
および中央二重平方表示を持つ noncubic 領域に分かれる。
前者では四生成・七二次関係に対する二つの可逆行列の存在によって、
選んだ平方根を忘れた pointed exact 同値を区別する。
後者は $\PP^1$ 上の二つの被約二点集合の非順序対の
同時射影同値で分類できる。全退化型を含み、
未知の AS 条件や圏同値の解決を分類条件には残さない。

\end{abstract}
\tableofcontents

\section{Introduction}
\label{sec:introduction}

滑らかな二次曲面 $Q=\PP^1\times\PP^1$ の偏極
$H=\OO_Q(1,1)$ は $H^{\otimes2}=\omega_Q^{-1}$ を満たす。
従って $\coh Q$ の Serre 関手 $S$ に対し
$s=-\otimes H$ は $s^2\simeq S^{-1}[2]$ という平方根である。
本稿ではこの性質と、非可換射影三次元空間内の二次超曲面としての表示を用い、
構造対象付き圏 $(\cC,\OO)$ の分類を目標とする。
以下の $\HH$ の添字 $2$ は超曲面の次数を表す。
可換例の反標準次数は $(-K_Q)^2=8$ であり、次数二の del Pezzo 曲面とは異なる。

三次超曲面の切断環では、二次関係だけから ambient を回復できる場合がある。
しかし二次超曲面では曲面を切り出す式自体が二次関係に含まれる。
切断環 $B$ の二次関係は七次元、ambient の二次関係は六次元となるため、
通常の quadratic envelope をそのまま用いることはできない。
適切な強い条件は、正則正規元 $g\in A_2$ による表示 $B\simeq A/(g)$ の存在である。
この修正は ambient の選択を新たな marking として加えるものではない。

\subsection{分類の中心結果}
\label{sec:intro-classification}

本稿の主定理は定理~\ref{thm:complete-h2-classification} である。
許容データから対象を構成し、任意の強い構造対象付き二次曲面から
cubic/noncubic のいずれかのデータを回復する。
その際、選択した Serre 平方根、helix foundation や ambient を
最終的な同値関係に含めず、四生成環の二次関係と根対配置によって
pointed exact 同値を判定する。
特に cubic 側では、同値が反 Serre 関手を移すことから、偶数・奇数辺に
一つずつの行列を選んで全ての関係を検査できる
（系~\ref{cor:two-matrix-pointed-equivalence}）。
例~\ref{ex:two-matrix-segre} は二行列を等しいと要求すると見落とす同値を与える。
一般の twisting system に関しては、Sierra
\cite[Theorem 1.2, Corollary 1.3]{Sierra2011} が添字付き環の同型と
Zhang twist の存在を同値とした。本稿で必要な段階は、標識だけを保つ
tails 同値から全ての自由シフトに対応する軌道対象を回復し、
その添字付き同型を二周期に選ぶことである。
Tran--Vancliff \cite[Theorem 2.7]{TranVancliff2024} は二次環上で、
指定した線形写像の冪と関係空間への作用に関する条件から twisting system を構成する。
既知の cubic $\mathbb Z$-algebra の分類は
Van den Bergh \cite{VdB2011} を入力とするが、
構造対象と平方根の忘却の扱いは別に証明する。
Normal helix、適合二周期の Serre 平方関係、AS4 被覆、
最初の軌道対象の内在的認識は、この分類を支える主要な道具である。

\subsection{相対モジュライ別稿との分担}

局所有限型基底の内在的相対 $\mathfrak H_2$ の有限型smooth Artin性、
cubic/noncubic 族の相対効果性、framed 実現、
被約noncubic境界の群商・慣性・good/coarse moduli と形式局所商は
別稿 \cite{QuadricModuli2026} で扱う。
体上の分類だけから非被約基底の相対結論は従わない。
本稿の分類の定理と補助結果は後続の各節にまとめた。
別稿の相対結果は分類主定理の証明には用いず、背景・後続問題として引用する。
別稿Bでは無標識の分類射の標的に対する対角の障害も扱う。

\subsection{先行研究と本稿の範囲}

分類する際に何を固定するかで、同じ cubic 曲面から異なる問題が生じる。
Van den Bergh \cite[arXiv v3, Theorem 5.5.10]{VdB2011} の対象は添字付きの
三次元 cubic AS 正則 $\mathbb Z$-代数であり、linear または許容四つ組と
geometric quintuple の間の対応は同型を射とする。偶数部分の hull についても、
次数 $-4$ の自己同型の類を指定して初めて一意性が成立し、標準的な類には
hull が存在する \cite[arXiv v3, Theorem 6.1]{VdB2011}。
通常の次数付き cubic 代数では Itaba--Matsuno--Saito
\cite[Theorems 3.3, 4.2, 4.3]{IMS2025} が、標数零の代数閉体上で
三種の可約点スキームについて定義関係を列挙し、次数付き環同型と
次数付き加群圏の Morita 同値をそれぞれ判定している。
本稿の系~\ref{cor:two-matrix-pointed-equivalence} が検査するのは、
適合平方根で作った四生成の切断環から、その平方根を忘れた
$(\cC,\OO)$ の同値を回復するための二次関係である。

Artin--Zhang \cite{AZ1994} の ample pair の再構成を用いる。
Mori--Ueyama \cite[Definition 2.5, Theorem 2.6(2)]{MU2021} は、
Hom 有限な Abel 圏に ample pair $(\OO,s)$ を指定すると、
非負軌道環の tails が $\OO$ と $s$ を含む triple を再現することを示した。
また Mori--Nyman \cite[Theorem 6.13]{MN2025} では、四次元二次 AS 正則環の
正則中心二次商が domain で、その tails が滑らかであり商が standard なら、
次元三、Gorenstein parameter 四の AS 正則 $\mathbb Z$-代数の tails が得られる。
中心切断の ruling と standard 性の関係は \cite{SVdB2013,MU2021}、
正規切断の Zhang twist に関する判定は \cite{MUW2026} による。
定理~\ref{thm:normal-helix} は強い $\HH$ の正則正規切断から出発し、
二つの初期 MCM 型から構成した窓の乗法を調べる。この段階で
cubic 表示を持つ範囲と、その外の中央二重平方型を判別できる。
退化曲線を含む一般のcubic偶数環のdomain性には
Presotto \cite[Theorem 2.8; Appendix A]{Presotto2018} を用いるが、
二周期の平方と反Serreの比較は本稿で個別に証明する。

Okawa--Ueda \cite[Theorem 1.1]{OU2014} の
geometric quintupleのGIT商$\PP(2,4,4,6)$、
conifold potentialのGIT商$\PP(1,2,3,4)$ と
次数四の射は既知の重要なモジュライ記述である。
GITの軌道閉包・S-equivalenceを本稿のpointed exact同値と
同一視しない。可換quadricの四つの適合根と既存の次数四射の関係も
未解決である。
Okawa による Kitamura との共同研究の報告
\cite[Definition 2, Theorems 1, 3]{OkawaKitamura2022OWR} では、
cubic $\mathbb Z$-代数の無標識 $\operatorname{qgr}$ 同値を
$D_\infty$ の軌道で判定する。そこでいう line bundle は階数一で
torsion-free、自己 Euler 数一の対象であり、$\mathbb Z^2$ で記述される。
本稿では同値に指定対象 $\OO$ の保存を課す。これを忘れる際の
ファイバーは、滑らかな非共鳴楕円型について別稿 \cite{QuadricModuli2026}
で計算する。

Belmans \cite[Proposition 31, Corollary 32, Table 2]{Belmans2019} は
既存のcubic AS正則$\mathbb Z$-algebraに由来する
非可換二次曲面のHochschild高次消滅と退化型の数値を計算した。
これは本稿のcubic領域の重要な先行研究である一方、
cubic表示を持たない本稿のnoncubic対象の
全分類・pointed同値判定を与えるものではない。

局所有限型基底上の相対Koszul・Noetherian貼合せ、
Artin変形・Serre根、canonical tailsと
noncubic境界の相対groupoidは別稿
\cite{QuadricModuli2026} にまとめる。
そこではLowen--Van den Bergh \cite{LowenVanDenBergh2006}、
De Deken--Lowen \cite{DeDekenLowen2011}、
Polishchuk--Positselski \cite{PolishchukPositselski2005}の
既知定理と相対効果性を区別する。
八・十・十二本矢の表示とavatarは別稿
\cite{QuiverAvatarModuli2026} の対象とする。
本稿は体上の分類に必要な証明を保持し、
B固有の相対主定理を分類の証明に循環して用いない。

\section{構造対象と存在条件}
\label{sec:definition}

以下、$k$ は標数零の代数閉体とする。
圏は本質的に小さい $k$-線形圏とし、Abel 圏が noetherian であるとは
その全対象が noetherian であることをいう。
次数付き環については右加群を用いる。

\begin{definition}[Gabriel--Krull rank one torsion-free]
\label{def:gk-rank}
noetherian Abel 圏 $\cC$ の有限長対象からなる Serre 部分圏を $\cC_0$、
$\cC/\cC_0$ で有限長となる対象の逆像を $\cC_1$ とし、
$q:\cC\to\cC/\cC_1$ を商関手とする。
$\OO\in\cC$ が Gabriel--Krull rank one torsion-free であるとは、
\[
 q\OO\text{ が非零単純であり、}\qquad
 \Hom(T,\OO)=0\quad(T\in\cC_1)
\]
であることをいう。これは ``GK dimension が $1$'' という条件ではない。
さらに $\OO$ が exceptional であるとは
$\End(\OO)=k$ および $\Ext^i(\OO,\OO)=0$（$i>0$）をいう。
\end{definition}

exact 自己同値は有限長性と Serre 商の有限長性を保つので、
$\cC_0,\cC_1$ を保つ。従ってこの rank と torsion-free 条件は exact 同値で不変である。

\begin{definition}[Artin--Zhang ample pair]
\label{def:ample}
$s$ を $\cC$ の exact 自己同値とする。$(\OO,s)$ が ample であるとは、
次の二条件を満たすことをいう。
\begin{enumerate}[label=(A\arabic*)]
\item 任意の $M\in\cC$ に対し、ある $n_1,\ldots,n_r\ge0$ と全射
      $\bigoplus_{j=1}^r s^{-n_j}\OO\twoheadrightarrow M$ が存在する。
\item 任意の全射 $M\twoheadrightarrow N$ に対し、十分大きい全ての $n$ で
      $\Hom(s^{-n}\OO,M)\to\Hom(s^{-n}\OO,N)$ が全射である。
\end{enumerate}
切断環を
\[
 B(\cC,\OO,s)=\bigoplus_{n\ge0}\Hom(\OO,s^n\OO),
 \qquad a*b=s^n(a)\circ b\quad(a\in B_m,\ b\in B_n)
\]
と定める。この積は右加群の再構成と整合する規約である。
\end{definition}

\begin{definition}[型 $(1,4,6,4,1)$]
\label{def:as14641}
$A$ は局所有限、連結、次数一生成、関係イデアルが二次で生成され、左右 noetherian とする。
$A$ がここでいう型 $(1,4,6,4,1)$ とは、Koszul であり、
global dimension と Gorenstein parameter がともに $4$ の AS 正則環で、
自明加群の極小分解が
\[
 0\to A(-4)\to A(-3)^4\to A(-2)^6\to A(-1)^4\to A\to k\to0
\]
となることをいう。AS 条件は左右で要求する。
特に $\Ext_A^i(k,A)$ は $i=4$ 以外で零、$i=4$ で $k(4)$ であり、
$H_A(t)=(1-t)^{-4}$ である。
元 $g$ が正則正規であるとは、左右の零因子でなく、$gA=Ag$ となることをいう。
\end{definition}

\begin{definition}[$\OO$-marked 非可換二次曲面の強いクラス]
\label{def:h2}
$\HH$ の対象は対 $(\cC,\OO)$ であり、次を要求する。
\begin{enumerate}[label=(Q\arabic*)]
\item $\cC$ は $k$-線形 noetherian Abel 圏で、$D^b(\cC)$ は Hom 有限かつ
      Serre 関手 $S$ を持つ。
\item $\OO$ は定義~\ref{def:gk-rank} の rank one torsion-free exceptional 対象である。
\item exact 自己同値 $s$ が存在し、その導来自己同値について
      $s^2\simeq S^{-1}[2]$ かつ $(\OO,s)$ が Artin--Zhang ample である。
\item \emph{同じ $s$} に対し、型 $(1,4,6,4,1)$ の環 $A$ と
      正則正規元 $g\in A_2$ が存在し、全次数で
      $B(\cC,\OO,s)\simeq A/(g)$ という次数付き環同型がある。
\end{enumerate}
$s$、根の同型 $s^2\simeq S^{-1}[2]$、$A$、$g$ および商表示は指定構造ではない。
対象間の同値は $k$-線形 exact 同値 $F:\cC\xrightarrow{\sim}\cC'$ と
$F(\OO)\simeq\OO'$ によるものとする。
\end{definition}

\begin{remark}
条件 (Q4) は一次生成性と Hilbert 関数を含意する。
これらを平方根と豊富性のみの帰結とはしない。
また $g$ の中心性は定義に含めない。中心性を課した表示だけを先に調べる場合も、
それが全対象を覆うかは別途証明する必要がある。
単なる $\qgr A/(g)\simeq\cC$ だけでなく、指定 $\OO$ と選んだ $s$ に対する
\emph{全切断環} の表示を要求していることにも注意する。
\end{remark}

平方根を保持した補助対象 $(\cC,\OO,s)$、さらに ambient 表示を保持した補助対象から、
最終対象 $(\cC,\OO)$ への忘却を考える。
平方根付きの同値では $Fs\simeq s'F$ も要求する。
根同型を固定した場合の整合性や慣性は、本稿では分類しない。
この区別により、表示の重複を異なる二次曲面として数えることを避ける。

% !TEX root = ../nc-quadric-surface.tex
\section{全対象の分類定理}
\label{sec:classification}

この節では分類のデータ、許容条件、同値関係を先に述べる。
構成と被覆の証明は後続の各節に置く。
平方根と ambient は構成に使うが、最終的な同値関係では保持しない。

\subsection{cubic 側の幾何データと四生成環}

\begin{definition}[許容四つ組]
\label{def:classification-admissible}
$P_{\mathrm{cub}}$ は一つの記号 $\mathrm{lin}$ と、次の四つ組
$U=(C,L_0,L_1,L_2)$ 全体との直和とする。
$C$ は射影曲線、$L_i$ は直線束であり、次を満たす。
\begin{enumerate}[label=(D\arabic*)]
\item $h^0(C,L_i)=2$ で、$L_i$ は大域切断で生成される。
      隣接する組 $(L_0,L_1)$ と $(L_1,L_2)$ による射はそれぞれ
      $C$ を $\PP^1\times\PP^1$ の二次数 $(2,2)$ の Cartier 因子として閉埋め込む。
\item 各既約成分 $C_j\subset C_{\mathrm{red}}$ 上で
      $\deg(L_0|_{C_j})=\deg(L_2|_{C_j})$ である。
\item $L_0\not\simeq L_2$ である。
\end{enumerate}
これは \cite[arXiv v3, Proposition 5.1.2]{VdB2011} の許容四つ組の条件である。
$C$ の整性、被約性、滑らかさは要求しない。
記号 $\mathrm{lin}$ は $(\coh(\PP^1\times\PP^1),\OO)$ を表す。
\end{definition}

四つ組から $L_{i+3}=L_i^{-1}\otimes L_{i+1}\otimes L_{i+2}$ により
両側の helix を作り、
\[
 V_i=H^0(C,L_i),\qquad
 R_i=\ker\bigl(V_i\otimes V_{i+1}\otimes V_{i+2}
       \longrightarrow H^0(C,L_i\otimes L_{i+1}\otimes L_{i+2})\bigr)
\]
と置く。$\dim V_i=\dim R_i=2$ である。
これらを隣接空間と三次関係とする $\mathbb Z$-algebra を $D(U)$ とする。
定理~\ref{thm:all-admissible-inverse} の構成に従い、
反 Serre の四周期 $\theta$ と $\beta^2=\theta$ を対角型まで満たす
二周期 $\beta:D(U)\to D(U)(2)$ を一つ選ぶ。
この選択の存在は全ての許容データで証明されている。
具体的には曲線の自己同型 $\sigma$ と束の同型
$\rho_i:\sigma^*L_i\xrightarrow\sim L_{i+2}$ を取り、
$\beta_i=H^0(\rho_i)\circ\sigma^*$ とする。

$V=V_0\otimes V_1$ とし、$\beta_{01}=\beta_0\otimes\beta_1:
V\to V_2\otimes V_3$ と書く。有限次元の空間
\begin{equation}
 K_U=(1_V\otimes\beta_{01}^{-1})
       (R_0\otimes V_3+V_0\otimes R_1)\subset V\otimes V,
 \qquad B_U=T(V)/(K_U)
 \label{eq:classification-cubic-kernel}
\end{equation}
を作る。$\dim V=4$、$\dim K_U=7$ である。
ここで $R_0\otimes V_3$ と $V_0\otimes R_1$ の交わりは
四次テンソルの張る一次元空間 $W_0$ である。
三次関係を偶数頂点間でまとめれば、二次環 $B_U$ は
\[
 (B_U)_n=D(U)_{0,2n},\qquad a*b=a\beta^m(b)
 \quad(a\in(B_U)_m)
\]
と一致する。これは曲線 $C$ の切断環ではない。
全次数での二次表示は、偶数長の道に現れる $R_j$ の開始添字が偶数なら
右に一文字、奇数なら左に一文字を加え、上の二種類の四文字関係にまとめることで従う。
特に次数四の全切断乗法の核を課すと曲線への追加商となるので、
式~\eqref{eq:classification-cubic-kernel} の七次元空間を用いる。
linear データでは、$B_{\mathrm{lin}}$ を通常の Segre 環
$k[x_{00},x_{01},x_{10},x_{11}]/(x_{00}x_{11}-x_{01}x_{10})$ とする。

\subsection{cubic データ上の有限同値判定}

各 $V$ の基底を取り、上で得た二次関係を
$K,K'\subset k^4\otimes k^4$ とする。
$G=\GL_4(k)$ と、方程式で与えられる閉部分集合
\[
 E(K,K')=\{(g,h)\in G\times G\mid(g\otimes h)K=K'\}
\]
を使い、次の反復を定める。
\begin{equation}
 \begin{split}
 C_0&=G,\\
 C_{n+1}&=\operatorname{pr}_1(E\cap(G\times C_n))
       \cap\operatorname{pr}_2(E\cap(C_n\times G)).
 \end{split}
 \label{eq:classification-finite-relation}
\end{equation}
命題~\ref{prop:first-orbit-finite-equivalence} により、この下降列は有限段で安定する。
安定集合が非空のとき $U\sim_{\mathrm{cub}}U'$ と定める。
行列式を可逆とした座標環で、射影の消去と根基イデアルの比較だけを用いる判定である。
標数零の抽象的な代数閉体上では有限の代数的判定を意味し、
機械上で実行するには係数体が有効に与えられていることを前提とする。
基底、束の表示、二周期 $\beta$ の選び方によらないことは、次の定理で示す。

\subsection{noncubic データと主定理}

$P_{\mathrm{nc}}$ を、$\PP^1$ 上の被約な二点集合 $Z_1,Z_2$ の非順序対の全体とする。
同値関係 $\sim_{\mathrm{nc}}$ は $\PGL_2$ の同時作用である。
それぞれの零点が $Z_i$ となる非退化二次形式 $f_i(u,v)$ を選び、
\[
 C_{Z_1,Z_2}=k[u,v]\langle c,d\rangle/
       (c^2-f_1(u,v),\ d^2-f_2(u,v))
\]
と置く。全ての生成元の次数は一で、$u,v$ は中央とする。
この表示自体の次数シフトが適合根であるとは要求しない。
適合根と AS 正則被覆は定理~\ref{thm:ore-anticommuting-central} で構成する。

\begin{theorem}[全 $\HH$ の pointed exact 分類]
\label{thm:complete-h2-classification}
標数零の代数閉体 $k$ 上で、次の対応は全単射である。
\[
 (P_{\mathrm{cub}}/{\sim_{\mathrm{cub}}})
 \ \sqcup\ (P_{\mathrm{nc}}/{\sim_{\mathrm{nc}}})
 \ \xrightarrow{\ \sim\ }\ 
 \HH/{\text{pointed exact equivalence}}.
\]
cubic 側では $U$ を $(\qgr B_U,\pi B_U)$ へ、
noncubic 側では $\{Z_1,Z_2\}$ を
$(\qgr C_{Z_1,Z_2},\pi C_{Z_1,Z_2})$ へ送る。
両領域の像は交わらず、前者は pointed cubic 表示を持つ対象の全体、
後者はそのような表示を持たない対象の全体である。

noncubic 側の商は、二対が一致する一類、一点を共有する一類と、
四点が異なる場合の $\lambda\in k\setminus\{0,1\}$ を
$\lambda\sim\lambda^{-1}$ で割った一パラメータ族からなる。
この領域で non-standard な適合表示を持つのは、二対が一致する一類だけである。
\end{theorem}

\begin{proof}
まず全許容四つ組の $D(U)$ は cubic AS 正則 $\mathbb Z$-algebra である
\cite[arXiv v3, Theorems 5.5.7--5.5.10]{VdB2011}。
定理~\ref{thm:all-admissible-inverse} は、全曲線型で適合二周期を構成し、
even part の全切断環としての表示、型 $(1,4,6,4,1)$ の被覆、
正則正規元、Gabriel--Krull 条件、exceptionality、AZ ampleness を確認する。
従って cubic 側の写像は定義される。
linear データは通常の二次曲面であり、$-\otimes\OO(1,1)$ を使えばよい。
noncubic 側の構成と同値判定は定理~\ref{thm:noncubic-classification} による。

任意の $(\cC,\OO)\in\HH$ を取る。
定理~\ref{thm:cubic-coverage} の階数判定により、pointed cubic 表示を持つか、
二つの MCM 型の階数がともに三であるかに分かれる。
前者では \cite[arXiv v3, Theorem 5.5.10]{VdB2011} の分類により、
linear データまたは定義~\ref{def:classification-admissible} の四つ組から得られる。
後者では定理~\ref{thm:noncubic-classification} により中央二重平方環へ帰着する。
これで全被覆と二領域の非交差を得た。

cubic 側の各 $B_U$ は、構成した同じ pointed 圏の適合切断環である。
命題~\ref{prop:first-orbit-finite-equivalence} は二つの適合切断環について、
反復~\eqref{eq:classification-finite-relation} の非空性が
pointed exact 同値と必要十分であることを示す。
従って $\sim_{\mathrm{cub}}$ は同値関係であり、写像は単射である。
一つの四つ組から異なる基底・束同型・適合二周期を選んでも、
いずれも $(\qgr D(U),\pi e_0D(U))$ と pointed 同値なので、同じ判定で同一視される。
具体的には $\check B_{ij}=B_{j-i}$ から $D(U)_{2i,2j}$ への写像
$a\mapsto\beta^i(a)$ が積と整合するので、偶数頂点の圏と添字零の対象を保つ。
未知の自己同値群や全ての平方根の列挙を、判定の前提として残してはいない。
noncubic 側の最後の具体的記述も定理~\ref{thm:noncubic-classification} による。
\end{proof}

\begin{remark}[分類の範囲]
主定理は全対象の同値類の分類である。
cubic 側の同値を幾何的な少数の正規形へさらに簡約すること、
全自己同値、平方根と ambient の各忘却ファイバー、
族の代数性と相対モジュライは別の問題である。
Sklyanin pencil の局所的な候補数の評価に現れる追加の単射性条件は、
式~\eqref{eq:classification-finite-relation} では使わず、
主定理の許容領域を狭めない。
既知入力と本稿の直接証明は区別しており、後者の包括的な新規性は未確認である。
\end{remark}


\section{内在的条件からの基本的帰結}
\label{sec:categorical-consequences}

\begin{proposition}[大域次元]
\label{prop:gldim-two}
条件 (Q1), (Q3) の平方根が存在し、$\cC\ne0$ ならば $\gldim\cC=2$ である。
\end{proposition}
\begin{proof}
$S\simeq s^{-2}[2]$ なので、Serre 双対性から
\[
 \Ext^i(M,N)\simeq\Hom_{D^b(\cC)}(s^2N,M[2-i])^* .
\]
$i>2$ では heart の対象間の負次数 Hom であるため零である。
非零 $N$ に対し $\Ext^2(s^2N,N)\simeq\End(s^2N)^*\ne0$ なので
大域次元は正確に $2$ である。
\end{proof}

\begin{lemma}[非零切断は単射]
\label{lem:sections-monic}
$\OO$ が rank one torsion-free で $s$ が exact 自己同値ならば、
任意の整数 $n$ と非零射 $f:\OO\to s^n\OO$ に対し $f$ は単射である。
従って非負次数の切断環 $B(\cC,\OO,s)$ は domain である。
\end{lemma}
\begin{proof}
$s$ は $\cC_1$ を保つため $q\OO$ と $q(s^n\OO)$ は非零単純であり、
$s^n\OO$ も torsion-free である。
$qf=0$ ならば $\operatorname{im}f\in\cC_1$ となり、
$\operatorname{im}f\hookrightarrow s^n\OO$ は零となるので $f=0$ に矛盾する。
よって $qf$ は単純対象間の非零射、従って同型であり、$\ker f\in\cC_1$ となる。
$\ker f\hookrightarrow\OO$ と torsion-free 性から $\ker f=0$ を得る。
切断環の非零斉次元の積は単射の合成であり零にならない。
一般の非零元については最高次数の斉次成分を取れば domain 性が従う。
\end{proof}

\begin{proposition}[切断環と再構成]
\label{prop:reconstruction}
$\HH$ の対象と条件を満たす平方根 $s$ に対して、切断環 $B$ は二次 domain であり、
\[
 H_B(t)=\frac{1-t^2}{(1-t)^4},\qquad \dim B_n=(n+1)^2\quad(n\ge0).
\]
負次数では $\Hom(\OO,s^{-m}\OO)=0$（$m>0$）であり、
\[
 (\cC,\OO,s)\simeq(\qgr B,\pi B,(1))
\]
という algebraic triple の同値がある。
\end{proposition}
\begin{proof}
Hilbert 級数と二次性は二次の正則元による商の完全列から従う
（命題~\ref{prop:quadric-hilbert}）。domain 性は補題~\ref{lem:sections-monic} による。
もし非零射 $\OO\to s^{-m}\OO$ があればそれは単射であり、
$s^{n+m}$ を作用させてから $\Hom(\OO,-)$ を取ると
$B_{n+m}\hookrightarrow B_n$ を得る。$n\ge0$ では
$(n+m+1)^2>(n+1)^2$ となり矛盾する。
最後は Hom 有限性と Artin--Zhang ampleness に対する再構成
\cite[Theorem 2.6(2)]{MU2021} の適用である。
ここでは $B$ が noetherian なので coherent module と finitely generated module の
区別は生じず、tails は通常の $\qgr$ となる。
\end{proof}

\begin{lemma}[一次生成と偶数部分列の ampleness]
\label{lem:ample-square-from-generation}
$\mathcal C$ を Hom 有限な $k$ 線型 Abel 圏，$E\in\mathcal C$，
$s$ を完全自己同値とする．$(E,s)$ が Artin--Zhang ample であり，
右規約の orbit algebra
$B=\bigoplus_{n\ge0}\operatorname{Hom}(E,s^nE)$ が一次生成であるとする．
このとき評価射
\[
 \operatorname{Hom}(s^{-1}E,E)\otimes_k s^{-1}E\longrightarrow E
 \tag{E}\label{eq:degree-one-evaluation}
\]
は全射であり，$(E,s^2)$ も Artin--Zhang ample である．
\end{lemma}
\begin{proof}
AZ の生成条件を $sE$ に適用してから $s^{-1}$ を作用させると，
$E$ は $s^{-m_i}E$ の有限直和の商になり，すべての $m_i$ を正に取れる．
これは生成条件の記述で指数 $0$ を許していても成り立つ．
一次生成より，$m>0$ に対する任意の射 $s^{-m}E\to E$ は
\eqref{eq:degree-one-evaluation}を経由する．
実際，対応する $B_m$ の元を次数一の元の積の線型結合として書き，
積 $a*b=s^{\deg b}(a)\circ b$ を $s^{-m}$ で移すと，
最後の射は常に $s^{-1}E\to E$ となる．
従って上記の全射は評価射を経由し，評価射自身も全射である．

任意の対象を生成する負の $s$-shift の直和のうち，
奇数指数の項 $s^{-m}E$ は，評価射を $s^{-m}$ で移すことにより
$s^{-(m+1)}E$ の有限直和で置き換えられる．
このため偶数指数だけで生成条件を満たす．
必要なら生成条件を先に $s^N F$ に適用してから $s^{-N}$ を作用させ，
すべての指数を任意の下限より大きくできる．
全射に対する Hom の最終的全射性は，元の AZ 条件を偶数部分列へ制限すればよい．
\end{proof}

\begin{corollary}
$\HH$ の全対象に対し、$(\OO,s^2)=(\OO,S^{-1}[2])$ も
Artin--Zhang ample である。
\end{corollary}
\begin{proof}
条件 (Q4) により切断環が一次生成なので、直前の補題を適用できる。
\end{proof}

\begin{remark}[Gabriel--Krull 次元]
条件 (Q2) と (A1) により、$\cC/\cC_1$ の全対象は有限長である。
実際、任意の対象は有限個の $s^{-n}\OO$ の直和の商であり、その像は単純対象の
有限直和の商となる。$q\OO\ne0$ なので Gabriel--Krull のこの規約で
$\cC$ は次元二である。さらに $B_1\ne0$ から $q\OO\simeq q(s\OO)$ となり、
商圏の単純対象は $q\OO$ の一種類に限られる。
従って rank-one 条件を次数付き環の GK dimension と取り違える必要はない。
\end{remark}

\begin{remark}[Nakayama twist の点検]
\label{rem:nakayama}
例えば \cite[Section 5]{MU2021} の中心二次切断の仮定の下では、
右加群規約で canonical bimodule は $M\mapsto M_\nu(-2)$ と書ける。
従って $S^{-1}[2]$ は $M\mapsto M_{\nu^{-1}}(2)$ である。
通常の次数シフト $s=(1)$ が平方根になるには、誘導される
$M\mapsto M_{\nu^{-1}}$ が tails 上で自然に恒等であることが必要である。
この性質を環自己同型 $\nu=\mathrm{id}$ と同一視しない。
また単に一部の対象が $\nu$ で不変であるだけでは、関手の自然同型は得られない。
一般の正規切断へ公式を使う場合は双対化加群の仮定と規約を個別に照合する。
\end{remark}

% !TeX root = ../nc-quadric-classification.tex
% 2026-11-21 / v0.1
\section{二次包絡の障害と二次超曲面表示}
\label{sec:quadric-envelope-obstruction}

以下、$k$ は代数閉体とし、代数は $k$ 上の局所有限な連結次数付き代数とする。
次数一で生成される代数 $B$ に対し、$V=B_1$ と置き、通常の二次包絡を
\[
 \operatorname{QEnv}(B)
 :=T(V)/\bigl(\ker(V\otimes V\longrightarrow B_2)\bigr)
\]
と定める。括弧は両側イデアルの生成を表す。
これは $B$ の二次関係をすべて課し、三次以上の独立な関係を忘れる構成である。
したがって $B$ 自身が二次代数ならば
$\operatorname{QEnv}(B)=B$ である。

型 $(1,4,6,4,1)$ は定義~\ref{def:as14641}の意味で用いる。

\begin{proposition}[次数二の正則元による商の数値]
\label{prop:quadric-hilbert}
$A$ を上記の型の代数、$g\in A_2$ を非零の両側正則正規元とし、
$B=A/(g)$ と置く。このとき $B$ は二次代数であり、
\[
 H_B(t)=\frac{1-t^2}{(1-t)^4}
       =\frac{1+t}{(1-t)^3},\qquad
 \dim_k B_n=(n+1)^2\quad(n\geq0).
\]
特に $\dim B_1=4$、$\dim B_2=9$ であり、
$\dim\ker(B_1\otimes B_1\to B_2)=7$ である。
\end{proposition}

\begin{proof}
正規性から $(g)=gA=Ag$ であり、左から $g$ を掛ける写像によって
右加群の完全列
\[
 0\longrightarrow A(-2)\xrightarrow{\,g\cdot\,}A
 \longrightarrow B\longrightarrow0
\]
を得る。したがって Hilbert 級数は主張の通りである。
係数は $\binom{n+3}{3}-\binom{n+1}{3}=(n+1)^2$ である
（第二項は $n<2$ なら零とする）。
$A$ が二次代数で $g$ が次数二なので、$B$ も二次代数である。
最後の次元は $4^2-9=7$ から従う。
\end{proof}

\begin{corollary}[通常の二次包絡を用いる定義の障害]
\label{cor:quadric-envelope-fails}
命題~\ref{prop:quadric-hilbert} の商 $B$ について
$\operatorname{QEnv}(B)=B$ であり、これは型 $(1,4,6,4,1)$ の
AS 正則代数ではない。
\end{corollary}

\begin{proof}
$B$ は二次代数なので二次包絡と一致する。
その次数二成分の次元は $9$ であるが、指定された型の AS 正則代数では
$10$ である。
\end{proof}

とくに $Q=\mathbb P^1\times\mathbb P^1$、
$\mathcal O=\mathcal O_Q$、$s=-\otimes\mathcal O_Q(1,1)$ に対しては、
$s^2\simeq S^{-1}[2]$ であり、$(\mathcal O,s)$ は Artin--Zhang ample である。
しかしその切断代数は
\[
 B=\bigoplus_{n\geq0}H^0(Q,\mathcal O_Q(n,n))
   \simeq k[x,y,u,v]/(xy-uv)
\]
なので、通常の二次包絡の AS 正則性を要求するとこの基本例を除外してしまう。
従って quadric における ambient 条件は、二次包絡の正則性から、
\emph{次数二の正則正規元を核に持つ AS 正則な二次被覆の存在}へ置き換える。
$s$ も被覆も、$\mathcal O$-marked object の指定データには含めず、
存在条件の証人として扱う。

\section{固定された切断代数の ambient 被覆}
\label{sec:ambient-incidence}

ここでは $B$ を二次代数とし、$V=B_1$、
$K=\ker(V\otimes V\to B_2)$ と置く。
さらに $\dim V=4$ および
$H_B(t)=(1-t^2)/(1-t)^4$ を仮定するので、$\dim K=7$ である。

\begin{lemma}[正規性と Hilbert 級数による正則性の判定]
\label{lem:normal-hilbert-regular}
$A$ は局所有限な連結次数付き代数、$0\ne g\in A_d$ は正規元とし、
$B=A/(g)$ とする。もし
$H_B(t)=(1-t^d)H_A(t)$ ならば、$g$ は左右の零因子でない。
\end{lemma}

\begin{proof}
右加群の完全列
\[
 0\longrightarrow\ker(g\cdot:A(-d)\to A)
 \longrightarrow A(-d)\xrightarrow{\,g\cdot\,}A
 \longrightarrow B\longrightarrow0
\]
を考える。正規性により像は商の核 $(g)$ に等しい。
Hilbert 級数を取ると、仮定は左端の核の Hilbert 級数が零であることを意味する。
局所有限性により核は零である。右から掛ける写像についても、
$Ag=(g)$ を用いて同じ議論を適用できる。
\end{proof}

\begin{proposition}[被覆の超平面表示]
\label{prop:ambient-hyperplane}
固定された $B$ に対する、次数一の同型 $A_1\to B_1=V$ を持つ
型 $(1,4,6,4,1)$ の AS 正則二次被覆
$A\twoheadrightarrow B$ で、その核が次数二の正則正規元で生成されるものは、
\[
 R\in\operatorname{Gr}(6,K)\simeq\mathbb P(K^*)
\]
のうち、次の二条件を満たすものによって分類される。
\begin{enumerate}
 \item $A_R:=T(V)/(R)$ は型 $(1,4,6,4,1)$ の AS 正則代数である。
 \item $(K/R)\subset(A_R)_2$ を張る任意の非零元 $g_R$ は正規元である。
\end{enumerate}
ここで被覆の同型は $B$ への写像と可換な次数付き代数同型とし、
核の生成元自体のスカラー倍は区別しない。
第二条件の下では $g_R$ の両側正則性は自動的に従う。
第一条件の下で、第二条件は二つの線形写像
\[
 V\longrightarrow(A_R)_3,\quad v\longmapsto g_Rv,
 \qquad
 V\longrightarrow(A_R)_3,\quad v\longmapsto vg_R
\]
がともに階数 $4$ を持つことと同値である。
\end{proposition}

\begin{proof}
与えられた被覆から
$R=\ker(V\otimes V\to A_2)$ を取ると、$\dim A_2=10$ より
$\dim R=6$ であり、合成写像から $R\subset K$ である。
二次性により $A=T(V)/(R)$ である。
核の次数二成分は $K/R$ であり、仮定された正則正規元の張る直線に等しい。

逆に $R$ が二条件を満たすとする。
$B=T(V)/(K)$ なので、$A_R/(g_R)=B$ である。
Hilbert 級数の仮定と第一条件から
$H_B(t)=(1-t^2)H_{A_R}(t)$ であり、
補題~\ref{lem:normal-hilbert-regular} によって $g_R$ は両側正則である。
この二構成は互いに逆である。実際、$B$ への写像と可換な被覆の同型は
次数一の同一視 $A_1=V$ を保存しなければならず、$A$ は次数一で生成されるので
同じ $R$ を与える。

階数による判定について、$J=\ker(A_R\to B)$ と置くと
$\dim J_3=\dim(A_R)_3-\dim B_3=20-16=4$ である。
二つの写像の像はともに $J_3$ に含まれるので、ともに階数 $4$ ならば
$g_RV=J_3=Vg_R$ である。次数一生成性を用いて生成元を一つずつ動かすと
$g_RA_R=A_Rg_R$ が従う。
逆に正規性が成り立てば、既に示した両側正則性から両写像は単射であり、
その階数は $4$ である。
\end{proof}

\begin{remark}[有限次数の判定と未解決部分]
$A_R$ は次数一で生成されるため、
$g_RV=Vg_R$ が $(A_R)_3$ の部分空間として成立することと
$g_R$ の正規性は同値である。
一方の包含から生成元を一つずつ動かして $g_RA_R\subset A_Rg_R$ を得て、
逆の包含も同様に示せるからである。
これは正規性の有限次元線形代数による判定であるが、
有限次数での Hilbert 関数の一致だけで AS 正則性まで証明できるとは主張しない。
また、この許容集合が開集合である、滑らかである、あるいは一点であるという
主張はここには含まれない。
\end{remark}

\section{同じ可換 quadric を持つ非同一の ambient 被覆}
\label{sec:ore-ambient-family}

\begin{proposition}[正規被覆を忘れる写像の非自明なファイバー]
\label{prop:ore-ambient-family}
$a\in k^{\times}$ に対し
\[
 A_a=k\langle x,y,u,v\rangle/
 \bigl([u,v],[u,x],[u,y],[v,x],[v,y],
       yx-a\,xy-(1-a)uv\bigr),\qquad g=xy-uv
\]
と定める。このとき次が成り立つ。
\begin{enumerate}
 \item $A_a$ は noetherian な二次 AS 正則 domain であり、型は
       $(1,4,6,4,1)$ である。
 \item $g$ は次数二の正則正規元であり、
       $gx=a xg$、$gy=a^{-1}yg$、$gu=ug$、$gv=vg$ を満たす。
 \item 生成元を同一視する標準同型
       \[
        A_a/(g)\simeq k[x,y,u,v]/(xy-uv)=:B
       \]
       がある。
 \item $B_1$ の同一視を固定した被覆としては、$a\ne b$ ならば
       $A_a\twoheadrightarrow B$ と $A_b\twoheadrightarrow B$ は同型でない。
       $A_1$ は可換であり、$a\ne1$ ならば $A_a$ は非可換で $g$ は非中心である。
\end{enumerate}
\end{proposition}

\begin{proof}
$R_0=k[u,v,x]$ とし、
\[
 \sigma(u)=u,\quad\sigma(v)=v,\quad\sigma(x)=ax,
 \qquad
 \delta(u)=\delta(v)=0,\quad\delta(x)=(1-a)uv
\]
と置く。$a\ne0$ なので $\sigma$ は次数を保つ自己同型である。
$\delta$ は次数一の $\sigma$-derivation に延びる。
たとえば可換関係 $xu-ux$ に対して
\[
 \delta(xu-ux)=(1-a)uvu-u(1-a)uv=0
\]
であり、$xv-vx$、$uv-vu$ についても同様である。
従って表示から
$A_a=R_0[y;\sigma,\delta]$ である。

この Ore 拡大は、多項式環の自明な Zhang twist を基礎環とする
次数一の Ore 拡大であるため、
\cite[Proposition~2.2]{LV2026}（$n=3$、twist は恒等写像）の仮定を満たす。
同命題により $A_a$ は noetherian AS 正則 domain、global dimension $4$、
Koszul である。この適用に必要な条件は直前で確認した。
Ore 拡大の基底
\[
 \{u^i v^j x^r y^s\mid i,j,r,s\geq0\}
\]
から $H_{A_a}(t)=(1-t)^{-4}$ を得る。
Koszul 性により極小自由分解の第 $i$ 項は次数 $i$ で生成されるので、
Hilbert 級数の逆数
\[
 (1-t)^4=1-4t+6t^2-4t^3+t^4
\]
から分解の型が $(1,4,6,4,1)$ であると分かる。

基底より $g=xy-uv\ne0$ である。$A_a$ は domain なので $g$ は両側正則である。
$u,v$ の中心性と表示関係から
\[
 \begin{split}
 gx&=xyx-uvx
     =x\bigl(a\,xy+(1-a)uv\bigr)-xuv
     =a x(xy-uv)=a xg,\\
 yg&=yxy-yuv
     =\bigl(a\,xy+(1-a)uv\bigr)y-uvy
     =a(xy-uv)y=a gy
 \end{split}
\]
を得る。これは主張された正規性の関係を与える。

商で $g=0$ を課すと $xy=uv$ であり、表示関係は
$yx=a uv+(1-a)uv=uv=xy$ となる。
残る関係はすべて可換関係なので、商はちょうど上記の $B$ である。

最後に、$W\subset V\otimes V$ を表示中の最初の五つの可換関係の張る部分空間、
$r_a=y\otimes x-a x\otimes y-(1-a)u\otimes v$ と置く。
$R_a=W+kr_a$ が被覆の二次関係空間である。
もし $R_a=R_b$ なら、$y\otimes x$ の係数を比較して
$r_a-r_b=(b-a)(x\otimes y-u\otimes v)\in W$ となる。
右辺の括弧内の元は $W$ に含まれないので $a=b$ である。
命題~\ref{prop:ambient-hyperplane} の同型に関する議論により、
これは固定された $B$ 上の被覆の非同型性を意味する。
$a\ne1$ なら $[y,x]=(a-1)g\ne0$ であり、
$gx-xg=(a-1)xg\ne0$ なので、非可換性と非中心性も従う。
\end{proof}

\begin{remark}[この命題が分類について述べる範囲]
上の族は、固定された切断代数 $B$ から ambient 正規被覆が一意に復元される
という主張への反例である。同じ
$(\operatorname{coh}(\mathbb P^1\times\mathbb P^1),\mathcal O,s)$ を与えるので、
異なる $a$ を異なる $\mathcal O$-marked quadric と数えてはならない。
ただし、基底の代数 $B$ の自己同型まで許す場合の $a$ の同値関係を
この命題は分類していない。次数二の中心元だけを許す被覆という別の問題も、
上記の正規被覆の問題と区別する必要がある。
この具体例の文献上の新規性は未確認であり、新分類定理として扱わない。
\end{remark}


% !TEX root = ../nc-quadric-surface.tex
\subsection{二次被覆の有限判定}
\label{subsec:finite-ambient-test}

\begin{theorem}
\label{thm:finite-ambient-test}
$B=T(V)/(K)$ を左右 noetherian な連結二次 Koszul 環とし、
$H_B(z)=(1+z)/(1-z)^3$ とする。
超平面 $J\subset K$ に対して $A=T(V)/(J)$ と置き、
$K/J$ の非零元を $g$ とする。次は同値である。
\begin{enumerate}
\item $A$ は型 $(1,4,6,4,1)$ の環であり、$g$ は左右正則正規である。
\item 次の有限次元の条件が成り立つ。
\begin{enumerate}
\item $A_3$ で $gV=Vg$ である。
\item $A_i\to A_{i+2}$ の $g$ の左右乗法は $i=1,2$ で単射である。
\item 二次双対 $A^!=T(V^*)/(J^\perp)$ の次数零から五までの次元は
$1,4,6,4,1,0$ である。
\item $A^!_i\otimes A^!_{4-i}\to A^!_4$ の積の組合せが $i=1,2,3$ で非退化である。
\end{enumerate}
\end{enumerate}
従って $B$ の AS Gorenstein 性を先に仮定せず、適合被覆の存在は、
$\operatorname{Gr}(6,K)$ 上で有限個の階数条件を満たす点の有無で判定できる。
\end{theorem}

\begin{proof}
必要性は正規性と Koszul 双対の Hilbert 級数および Frobenius 性から従う。
逆向きに、一次生成性から $gA=Ag$ である。
$\operatorname{Tor}$ の環変換スペクトル系列を用いる。
右 $A$-加群 $B$ に対し $A(-2)\xrightarrow{g\cdot}A\to B\to0$ は極小表示である。
左乗法に核があると仮定し、その最小次数を $d$（$A$ の元の次数）とする。
$g\ne0$ と単射の条件から $d\ge3$ である。
極小分解の第二項から、$\operatorname{Tor}^A_2(B,k)$ の
最小非零内部次数は $d+2$ となる。
また $\operatorname{Tor}^A_1(B,k)\simeq k(-2)$ は
内部次数二だけに集中する左 $B$-加群である。

\[
 E^2_{p,q}=\operatorname{Tor}^B_p(k,\operatorname{Tor}^A_q(B,k))
 \ \Longrightarrow\ \operatorname{Tor}^A_{p+q}(k,k)
\]
において、$E^2_{0,2}$ の内部次数 $d+2$ は非零である。
実際、この最小次数には $B_{\ge1}$ による積が入らない。
この部分に届く可能性のある微分は $(2,1)$ および $(3,0)$ からであるが、
$B$ が Koszul なのでそれぞれ内部次数四および三に集中する。
$(0,2)$ から出る微分はなく、内部次数 $d+2\ge5$ の元は残る。
これは二次環 $A$ の $\operatorname{Tor}^A_2(k,k)$ が
内部次数二に集中することに反する。
反対環にも同じ議論を適用して、$g$ の左右正則性を得る。

従って $B$ の右 $A$-自由分解は長さ一であり、環変換スペクトル系列には
$q=0,1$ の二行のみが残る。$b_p=\dim(B^!)_p$ とすると、
$E^2_{p,0}$ は内部次数 $p$ に $b_p$ 次元、
$E^2_{p,1}$ は内部次数 $p+2$ に $b_p$ 次元ある。
全次数 $n$ の対角内部次数 $n$ に寄与するのは
$E^\infty_{n,0}=\ker(d_2:E^2_{n,0}\to E^2_{n-2,1})$ のみである。
二次環の対角 Tor の次元は二次双対の次元と一致する。
$A^!$ は一次生成なので、次数五の消滅からそれより上も消滅する。
よって仮定により
\[
 \dim E^\infty_{n,0}
 = [z^n](1+z)^4
 = b_n-b_{n-2},
\]
ただし負添字の $b_p$ は零とする。
最後の等号は $H_{B^!}(z)=(1+z)^3/(1-z)$ から従う。
従って全ての $d_2:E^2_{n,0}\to E^2_{n-2,1}$ は全射であり、
非対角の Tor は残らない。これで $A$ は Koszul となる。
同じ対角次元を持つ反対環にも適用できる。

正則正規元で割った環 $B$ が左右 noetherian なので、$A$ も左右 noetherian である。
実際、$(g)$ 進次数環は $B$ 上の skew polynomial 環であり、
斉次イデアルでは固定内部次数でフィルトレーションが有限なので
初期イデアルから有限個の生成元を持ち上げられる。
一般のイデアルには内部次数による最高次成分を用いる。
二次双対は最高次数四で積の組合せが完全なので graded Frobenius である。
Koszul 複体を双対化すると、この組合せによって左右の複体が
内部次数四と複体次数四の移動を除いて互いに同定される。
従って $\Ext_A^j(k,A)$ および反対環の Ext は、$j=4$ の $k(4)$ を除き零となる。
Koszul 性と二次双対の次元により自明加群の極小分解は
長さ四で所要の型となる。任意の有限生成次数付き加群の極小自由分解の長さも、
反対側の $k$ の分解を用いて Tor を計算すれば四以下である。
非次数付き有限生成加群には生成元からの次数フィルトレーションを入れ、
付随次数加群の自由分解を持ち上げると同じ上界を得る。
従って巡回加群に対する射影次元の上界から大域次元も四となる。
正則二次切断から $H_A(z)=(1-z)^{-4}$ が従う。
\end{proof}

% !TeX root = ../nc-quadric-classification.tex
% 2026-11-21 / v0.1
\section{可換二次曲面と平方根の非一意性}
\label{sec:commutative-roots}

この節では $k$ を標数 $0$ の代数閉体，
$Q=\mathbb P^1_k\times\mathbb P^1_k$ とする．
$\mathcal O_Q(a,b)$ を通常の線束とし，
$S_Q=(-)\otimes\mathcal O_Q(-2,-2)[2]$ と書く．
ここで分類する平方根は，$k$ 線型完全自己同値 $s$ であって
$s^2\cong S_Q^{-1}[2]$ を満たすものの同型類である．
この同型そのものは指定せず，Serre 双対性の trace と両立する
追加の巡回的 coherence 条件も課さない．
根の同値は，$F(\mathcal O_Q)\cong\mathcal O_Q$ を満たす
完全自己同値 $F$ と $Fs\cong s'F$ の存在を意味する．
この節は族や慣性の分類を主張しない．

\begin{lemma}[可換標識対象]
\label{lem:commutative-marking-line}
$Q$ 上の rank-one torsion-free exceptional coherent sheaf $E$ は線束である．
したがって $(\operatorname{coh}Q,E)$ は
$(\operatorname{coh}Q,\mathcal O_Q)$ と pointed exact 同値である．
\end{lemma}
\begin{proof}
$Q$ は滑らかな曲面なので $L=E^{\vee\vee}$ は線束であり，
$T=L/E$ は有限長である．$n=\operatorname{length}T$ と置く．
$E$ は完全複体として扱え，Riemann--Roch により
\[
 \chi(E,E)
 =\int_Q\operatorname{ch}(E^{\vee}_{\mathrm{der}})
             \operatorname{ch}(E)\operatorname{td}(Q)
 =\chi(\mathcal O_Q)-2n=1-2n.
\]
実際，$D=c_1(L)$ とすれば
$\operatorname{ch}(E)=e^D(1-n[\mathrm{pt}])$ であり，
導来双対の Chern character は $e^{-D}(1-n[\mathrm{pt}])$ である．
exceptionality から $\chi(E,E)=1$ なので $n=0$，従って $E=L$ である．
$(-)\otimes L^{-1}$ が求める pointed exact 同値を与える．
\end{proof}

\begin{theorem}[四つの平方根]
\label{thm:commutative-four-roots}
$\operatorname{coh}Q$ の $k$ 線型完全自己同値 $s$ が
\[
 s^2\cong(-)\otimes\mathcal O_Q(2,2),
 \qquad \dim_k\operatorname{Hom}(\mathcal O_Q,s\mathcal O_Q)=4
 \tag{$*$}\label{eq:commutative-root-conditions}
\]
を満たすとする．このような $s$ は pointed exact 同値を除いてちょうど四種類であり，
\[
 s_\varphi(F)=\varphi^*F\otimes\mathcal O_Q(1,1)
\]
において $\varphi$ を次のいずれかに取ったものである．
\begin{enumerate}
 \item 恒等写像．
 \item 一方の因子の非自明な対合 $\iota\times\mathrm{id}$．
 \item 両因子の非自明な対合 $\iota\times\iota$．
 \item 二因子の交換 $\epsilon$．
\end{enumerate}
ここで $\iota([u_0:u_1])=[u_0:-u_1]$ と取れる．
四者はいずれも $(\mathcal O_Q,s_\varphi)$ が Artin--Zhang ample であり，
対応する右規約の orbit algebra は $(14641)$ 型 quadratic AS-regular domain の
正則正規な二次斉次元による商である．
\end{theorem}

\begin{proof}
まず完全自己同値の形を決める．
Bondal--Orlov の定理 \cite[Theorem~3.1]{BondalOrlov2001} は，
反標準束が ample な滑らかな射影多様体の導来自己同値を
自己同型，線束による tensor，shift で記述する．
Abel 圏の完全自己同値から得られる導来自己同値は標準 heart を保つため shift は零である．
従って $s=\varphi^*(-)\otimes L$ と書ける．
平方根の条件から
\[
 \varphi^2=\mathrm{id},\qquad
 L\otimes\varphi^*L\cong\mathcal O_Q(2,2).
 \label{eq:commutative-root-norm}
\]
$\varphi$ が二つの ruling をそれぞれ保つ場合，Picard 群への作用は恒等なので
$L=\mathcal O_Q(1,1)$ である．
$\varphi$ が ruling を交換する場合は
$L=\mathcal O_Q(a,2-a)$ である．
このとき $h^0(L)=4$ はまず $0\le a\le2$ を強制し，さらに
\[
 (a+1)(3-a)=4
\]
から $a=1$ を強制する．従って，いずれの場合も $L=\mathcal O_Q(1,1)$ である．

$F(\mathcal O_Q)\cong\mathcal O_Q$ を保つ完全自己同値 $F$ は，
同じ定理により $\psi^*$ の形に限られる．
すべての $\psi\in\operatorname{Aut}(Q)$ は
$\mathcal O_Q(1,1)$ の同型類を保つので，求める同値類は
$\operatorname{Aut}(Q)$ における対合（恒等元を含む）の共役類である．
\[
 \operatorname{Aut}(Q)
 =\bigl(\operatorname{PGL}_2(k)\times\operatorname{PGL}_2(k)\bigr)
       \rtimes\langle\epsilon\rangle.
\]
$\operatorname{PGL}_2(k)$ の非自明な対合は，標数が $2$ でなく $k$ が代数閉体なので
すべて $\iota$ に共役である．
恒等成分に属する対合は二因子の対合の組であり，
因子交換も許すと，非自明な因子の数 $0,1,2$ の三種類になる．
他方，恒等成分の外にある対合は $(g,g^{-1})\epsilon$ と書け，
$(g,\mathrm{id})\epsilon(g,\mathrm{id})^{-1}=(g,g^{-1})\epsilon$
よりすべて $\epsilon$ に共役である．
前の三種類は ruling の作用と各因子の作用で区別でき，
最後の種類は ruling を交換するのでこれらと共役ではない．

残りの条件を確認する．各 $\varphi$ は $\mathcal O_Q(1,1)$ を保ち，
\[
 s_\varphi^{-n}\mathcal O_Q\cong\mathcal O_Q(-n,-n)
 \qquad(n\in\mathbb Z)
\]
である．$\mathcal O_Q(1,1)$ の ampleness による大域生成と Serre 消滅から，
この列は Artin--Zhang の生成条件と全射に対する Hom の最終的全射性を満たす．
以下の命題が orbit algebra の AS-regular cover を与える．
\end{proof}

\begin{proposition}[四平方根の明示的な cover]
\label{prop:commutative-twist-covers}
$V=H^0(Q,\mathcal O_Q(1,1))$，$P=\operatorname{Sym}(V)$ とし，
Segre 方程式を
$q=x_{00}x_{11}-x_{01}x_{10}\in P_2$ と書く．
定理\ref{thm:commutative-four-roots}の各 $\varphi$ について，
その $V$ 上の持上げ $\sigma$ を $\sigma^2=1$ となるように選ぶ．
$P$ と同じ次数付きベクトル空間に積
\[
 f\star g=\sigma^{\deg g}(f)g
 \qquad(f,g\text{ homogeneous})
\]
を入れた代数を $A_\sigma$ とする．
すると $A_\sigma$ は左右 noetherian な quadratic Koszul AS-regular domain であり，
解消型は $(14641)$ である．
$P_2=(A_\sigma)_2$ というベクトル空間の同一視で得られる元 $q$ は正則正規で，
右規約
$a*b=s_\varphi^{\deg b}(a)\circ b$
による orbit algebra は $A_\sigma/(q)$ に同型である．
\end{proposition}

\begin{proof}
上記の積は，$\sigma$ が次数付き自己同型であることから結合的である．
$\sigma$ の固有基底 $z_1,\ldots,z_4$ を取り，固有値を $e_i\in\{1,-1\}$ とすると，
\[
 z_i\star z_j=(e_i/e_j)z_j\star z_i.
\]
従って $A_\sigma$ は四変数の skew polynomial ring である．
順序付き単項式は PBW 基底をなし，自己同型による反復 Ore 拡大として
左右 noetherian かつ domain である．
通常の Koszul 複体の交換係数を上記の符号へ変更すると，
正規正則列 $z_1,\ldots,z_4$ に沿う標準的な skew Koszul 複体が得られ，
\[
 0\longrightarrow A_\sigma(-4)
 \longrightarrow A_\sigma(-3)^4
 \longrightarrow A_\sigma(-2)^6
 \longrightarrow A_\sigma(-1)^4
 \longrightarrow A_\sigma\longrightarrow k\longrightarrow0
\]
は完全である．双対複体のコホモロジーは最高次数の $k(4)$ だけであり，
反復 Ore 拡大の大域次元は $4$，Hilbert 級数は $(1-t)^{-4}$ である．
これが主張した AS-regularity と解消型を与える．

$\varphi$ は $Q$ を保つので $\sigma(q)=c q$ であり，$c\in\{1,-1\}$ である．
$x\in V$ に対して
\[
 q\star x=\sigma(q)x=c qx,
 \qquad x\star q=\sigma^2(x)q=xq
\]
だから $q\star x=c(x\star q)$ となる．
従って $q$ は正規であり，domain 性から正則でもある．
二側イデアル $(q)$ は通常の $qP$ と同じ次数付きベクトル空間なので，
$A_\sigma/(q)$ は Segre 切断環 $P/(q)$ の同じ積による twist である．
$\varphi^*\mathcal O_Q(1,1)\cong\mathcal O_Q(1,1)$ の同型を
$\sigma$ と整合的に選ぶと，orbit algebra の積はまさに
$a*b=\sigma^{\deg b}(a)b$ になる．
\end{proof}

\begin{remark}
上の四代表で，$\sigma$ はそれぞれ恒等作用，
$x_{10},x_{11}$ の符号変更，$x_{01},x_{10}$ の符号変更，
$x_{01}$ と $x_{10}$ の交換として取れる．
$q$ の固有値 $c$ は順に $1,-1,1,1$ である．
従って一因子対合のこの cover では $q$ は各一次元と反可換になり，中心元ではない．
これは正則正規な二次斉次元を許す定義を点検する具体例になる．
ここで $q$ は $P_2=(A_\sigma)_2$ の同一視で指定した元であり，
$\star$ 積で表した式の係数は通常の可換積の表示と異なり得る．
\end{remark}

\begin{example}[平方根と ampleness だけでは一次元数は決まらない]
\label{ex:commutative-unbalanced-root}
二因子の交換 $\epsilon$ を用いて
\[
 t(F)=\epsilon^*F\otimes\mathcal O_Q(2,0)
\]
と置くと $t^2\cong(-)\otimes\mathcal O_Q(2,2)$ であり，
$(\mathcal O_Q,t)$ は Artin--Zhang ample である．
実際，$t^{-2m}\mathcal O_Q\cong\mathcal O_Q(-2m,-2m)$，
$t^{-(2m+1)}\mathcal O_Q\cong\mathcal O_Q(-2m,-2m-2)$ である．
偶数部分列による生成と，両部分列に対する Serre 消滅から ampleness が従う．
しかし
\[
 \dim_k\operatorname{Hom}(\mathcal O_Q,t\mathcal O_Q)
 =h^0(\mathcal O_Q(2,0))=3.
\]
従って，その orbit algebra は四つの一次生成元を持つ AS-regular 環の
二次斉次元による商ではない．
平方根の存在と ampleness だけから $(14641)$ 型 cover の数値条件を結論することはできない．
\end{example}

\begin{remark}[対象と根の区別]
補題\ref{lem:commutative-marking-line}により，heart を
$\operatorname{coh}(\mathbb P^1\times\mathbb P^1)$ に固定した可換対象には，
pointed exact 同値を除いて一つの rank-one torsion-free exceptional 標識がある．
一方，定理\ref{thm:commutative-four-roots}はその同じ対象上で許容平方根が四つの同値類を持つことを示す．
この四種類を四つの $\mathcal O$-marked 曲面として数えてはならない．
この節の結果は既知の自己同値記述と明示的計算による基礎例であり，
独立した新規性は主張しない．
\end{remark}

% !TeX root = ../nc-quadric-classification.tex
% 2026-11-21 / v0.2。引用番号の版は本文で指定する。
\section{楕円型 Sklyanin 二次曲面の例}
\label{sec:sklyanin-examples}

この節では $k=\mathbb C$ とする。
$T$ を \cite[\S 7, (7.1), (7.2)]{BSW2010} の通常の四次元 Sklyanin 環とし、
\cite[\S 10.1]{SVdB2013} の楕円型データ $(E,L,\tau)$ に対応すると仮定する。
ここで $E$ は滑らかな楕円曲線、$\deg L=4$、$\tau$ は無限位数の平行移動である。
原点は同文献の規約に合わせ、$L$ による埋込みで四点が同一平面上にあることと
その和が零であることが同値になるように選ぶ。
したがって平行移動 $\tau$ は $E$ の点ともみなす。
この仮定は modified Sklyanin 環や退化した点スキームを含めない。

次数二の中央元からなる二次元空間 $\mathbb C\Omega_1+\mathbb C\Omega_2\subset T_2$ の
非零元による商は、$z\in E$ を用いて
\[
 R_z=T/(g_z),\qquad Q_z=\operatorname{Proj}R_z,
 \qquad Q_z=Q_{-z-2\tau}
\]
と表示される \cite[\S\S 10.2--10.4]{SVdB2013}。
同文献の \S 2 は右加群で定義される一方、\S 10 の幾何的記述には
$T/TW$ という左加群の表示が使われる。
この幾何的表示を本稿の右加群規約に合わせるため、
\[
 A=T^{\mathrm{op}},\qquad B_z=R_z^{\mathrm{op}}
       =A/(g_z),\qquad
 \cC_z=\qgr B_z,\qquad \OO_z=\pi B_z
\]
と置く。以下の反自己同型により、左右いずれの規約による結果も
構造対象と次数シフトを保って移送できる。
反対環への移行は以下で用いる左右 AS 正則性、Koszul 性、domain 性、
中央正則元、Calabi--Yau 性を保つ。

\begin{lemma}[左右規約の比較]
\label{lem:sklyanin-opposite}
通常の Sklyanin 座標において、
$\theta(x_0)=-x_0$、$\theta(x_i)=x_i$（$1\le i\le3$）は
$T$ の次数付き反自己同型を定め、中央二次 pencil を各点で固定する。
従って $R_z\simeq R_z^{\mathrm{op}}$ であり、左右の coherent module 圏の間には
構造対象と次数シフトを保つ $\mathbb C$-線形 exact 同値がある。
\end{lemma}
\begin{proof}
\cite[\S 7]{BSW2010} の関係式を $r_i,s_i$ と書くと、
積の順序を逆転して上記の符号を代入する写像は $r_i$ を $r_i$、$s_i$ を $-s_i$ に送る。
その二乗は恒等なので反自己同型である。
\cite[\S 1.1, pp.~38--39]{LS1993} の中央二次元空間は
$x_0^2,x_1^2,x_2^2,x_3^2$ の線型結合で張られ、各元は $\theta$ で固定される。
同文献の仮定 $\tau\notin E[4]$ は無限位数の仮定から従う。
よって $\theta$ は $g_z$ を固定して商へ降りる。
左加群 $M$ に右作用 $m\cdot r=\theta(r)m$ を入れることが所要の同値を与える。
\end{proof}

\begin{proposition}[滑らかな楕円型の対象]
\label{prop:sklyanin-h2}
上の仮定の下で、$z+\tau\notin E[2]$ ならば
$(\cC_z,\OO_z)\in\HH$ である。
存在条件の証人として $s=(1)$ と次数二の中央正則元による表示 $B_z=A/(g_z)$ を取れる。
特に切断環は高次数で一致するだけでなく、全非負次数で
\[
 B(\cC_z,\OO_z,(1))\cong B_z
\]
である。ここで $E[2]$ は $E$ の二捩れ部分群を表す。
\end{proposition}

\begin{proof}
以下 $B=B_z$、$\cC=\cC_z$、$\OO=\OO_z$ と略記する。
\cite[\S\S 10.1--10.2]{SVdB2013} により、$A$ は左右 noetherian な
四次元二次 AS 正則 Koszul domain で、$H_A(t)=(1-t)^{-4}$ である。
従って極小分解は型 $(1,4,6,4,1)$ となる。
$g_z$ は中央正則であり、$B$ も noetherian domain である。
商は AS Gorenstein 次元 $3$、Gorenstein parameter $2$ を持ち、
\[
 H_B(t)=\frac{1-t^2}{(1-t)^4}=\frac{1+t}{(1-t)^3}.
\]
この商の AS Gorenstein 性と次数については
\cite[\S 5, Definition 5.1 直後]{MU2021} も参照できる。
\cite[Theorem 10.2]{SVdB2013} により、仮定 $z+\tau\notin E[2]$ の下で
$\cC$ は滑らかであり、$\gldim\cC=2$ である。
また $B$ は深さ $3$ の Cohen--Macaulay 環であり、両側の条件 $\chi$ を満たす
\cite[\S\S 2.5--2.8]{SVdB2013}。
$\chi$ による coherent cohomology の有限性、構造対象の shift による生成、
および $\gldim\cC=2$ から $D^b(\cC)$ は Hom 有限である。

まず Serre 関手の自己同型による捩れを確認する。
\cite[arXiv v3, Proposition 7.1]{BSW2010} は $T$ が非捩れの
Calabi--Yau 環であることを示す。従って $A$ の Nakayama 自己同型は恒等である。
\cite[arXiv v2, Lemma 1.5]{RRZ2014} を中央正則元 $g_z$ に適用できる。
同補題の $\mu_A$-normal 条件は $\mu_A(g_z)=g_z$ により満たされ、
正規化自己同型も恒等なので、$\mu_B=\mathrm{id}$ となる。
従って $B$ の canonical bimodule は両側加群として $B(-2)$ である。
\cite[Example 3.6(2) および \S 5 の Proposition 5.2 直後]{MU2021} により
\[
 \mathbb S_{\cC}\simeq(-2)[2],\qquad
 s^2=(2)\simeq\mathbb S_{\cC}^{-1}[2].
\]
ここでは hypersurface 環 $B$ 自体が有限 global dimension を持つとは主張していない。

$\mathfrak m=B_{\ge1}$ とする。
\cite[(2.2)]{SVdB2013} と $H^0_{\mathfrak m}(B)=H^1_{\mathfrak m}(B)=0$ により
\[
 B\xrightarrow{\sim}
 \bigoplus_{n\in\mathbb Z}\Hom_{\cC}(\OO,\OO(n))
\]
である。この写像は本稿の切断環の積と整合する。
従って全次数での飽和性が成り立ち、$\End(\OO)=\mathbb C$ を得る。
同文献 (2.3) と深さの条件から
\[
 \Ext^1_{\cC}(\OO,\OO)=H^2_{\mathfrak m}(B)_0=0.
\]
さらに Serre 双対性と飽和性により
\[
 \Ext^2_{\cC}(\OO,\OO)
 \cong\Hom_{\cC}(\OO,\OO(-2))^*
 \cong B_{-2}^*=0.
\]
次数 $i>2$ の Ext も消えるので $\OO$ は exceptional である。
有限生成加群の十分高い切断に生成元を取ると、任意の対象は有限個の
$\OO(-n)$、$n\ge0$、の直和の商になる。
また $\chi$ と (2.3) は任意の $K\in\cC$ に対し
$\Ext^1(\OO(-n),K)=0$ を $n\gg0$ で与える。
全射の核にこれを適用すれば Artin--Zhang ample の第二条件も従う。

最後に定義~\ref{def:gk-rank} の Gabriel--Krull rank 条件を確認する。
これは $B$ の Gelfand--Kirillov dimension や商体上の rank の計算だけでは済まない。
有限生成次数付き $B$-加群 $M$ は有限生成 $A$-加群でもあり、
$A$ 上の有限次数付き自由分解を持つので、Hilbert 関数は最終的に多項式 $P_M(n)$ となる。
$M$ は有限個の $B$ の shift の商なので $\deg P_M\le2$ である。
$P_M$ は $\pi M$ のみで決まり、短完全列に対して加法的である。
零多項式となることと $\pi M=0$ は同値である。

$\mathcal D_0\subset\cC$ を Hilbert 多項式が定数である対象、
$\mathcal D_1\subset\cC$ を次数が高々 $1$ である対象の部分圏とする
（いずれも零対象を含める）。両者は Serre 部分圏である。
非零な定数 Hilbert 多項式の値は正整数なので、
$\mathcal D_0$ の対象の任意の真の部分商列の長さはその定数で抑えられる。
従って $\mathcal D_0\subset\cC_0$ である。
同様に、次数 $1$ の Hilbert 多項式の一次係数は正整数であり、
$\cC/\mathcal D_0$ において $\mathcal D_1$ の対象の長さを抑える。
exact な商関手 $\cC/\mathcal D_0\to\cC/\cC_0$ は有限長性を保つから
\[
 \mathcal D_1\subset\cC_1.
 \tag{\dagger}\label{eq:sklyanin-polynomial-filtration}
\]
ここではこれらの包含を等号とする必要はない。

$g_z$ と線形独立な次数二の中央元 $h$ を同じ pencil から取る。
$B$ におけるその像も $h$ と書く。
$h$ は非零で正則であり、$B/(h)$ は楕円曲線 $E$ の twisted homogeneous
coordinate ring の反対環である \cite[\S 10.2]{SVdB2013}。
従って $\qgr B/(h)\simeq\coh E$ は $\cC$ に exact かつ fully faithful に埋め込まれ、
$\pi(B/(h))$ とその次数 shift は $E$ 上の可逆層に対応する。
可逆層は、点 $p\in E$ における零点の位数を増やす部分可逆層の下降列を持つので、
この埋込み後も有限長ではない。
そこで $\OO$ の下降列
\[
 \OO\supset\pi(hB)\supset\pi(h^2B)\supset\cdots
\]
を考える。隣接する商は $\pi(B/(h))(-2j)$ なので $\cC_0$ に属さない。
この列は $\cC/\cC_0$ でも真の下降列であり、その中の $\OO$ の像は有限長でない。
従って $\OO\notin\cC_1$ である。

$0\ne N\subset\OO$ を任意の部分対象とする。
Serre 商における部分対象の持ち上げにより $N=\pi I$ と書ける。
ここで $I\subset B$ は非零な次数付き右イデアルである。
次数 $d$ の非零元 $f\in I$ を取ると、domain 性により
$B(-d)\xrightarrow{f\cdot}I$ は単射なので、$\OO(-d)\subset N$ である。
shift は $\cC_1$ を保つから $N\notin\cC_1$ となる。
一方、$B/I$ は $B/fB$ の商であり、
\[
 H_{B/fB}(t)=(1-t^d)H_B(t)
\]
の Hilbert 多項式は次数高々 $1$ である。
従って \eqref{eq:sklyanin-polynomial-filtration} により
$\OO/N\in\cC_1$ となる。
$q:\cC\to\cC/\cC_1$ とすると、$q\OO$ は非零であり、
その部分対象は $\OO$ の部分対象から持ち上がるので、以上から $q\OO$ は単純である。
また $\cC_1$ に属する非零部分対象を $\OO$ は持たないので
$\Hom(T,\OO)=0$（$T\in\cC_1$）である。
これで定義~\ref{def:gk-rank} と $\HH$ の全条件が従う。
\end{proof}

\begin{remark}[standard 性と既存分類への入口]
\label{rem:sklyanin-standard}
上の滑らかな Sklyanin 商は \cite[Theorem 5.5]{MUW2026} により standard である。
反対環の規約は補題~\ref{lem:sklyanin-opposite} で処理できる。
従って命題~\ref{prop:standard-pointed} により、
三次元 AS 正則 $\mathbb Z$-algebra への比較で標識と平方根を追跡できる。
\end{remark}

\begin{proposition}[可換二次曲面と同値でない対象]
\label{prop:sklyanin-genuinely-noncommutative}
上と同じ $T$ に対して $z=\tau$ と取ると、
$(\cC_\tau,\OO_\tau)\in\HH$ であり、
$\cC_\tau$ は $\coh(\PP^1\times\PP^1)$ と $\mathbb C$-線形 exact 同値でない。
\end{proposition}

\begin{proof}
$\tau$ が無限位数なので $2\tau\notin E[2]$ であり、
命題~\ref{prop:sklyanin-h2} が適用できる。
注意~\ref{rem:sklyanin-standard} の standard 性により、以下では
一般の中心二次商についての誤った standard 性の主張を用いない。
\cite[\S 10.5, p.~854]{SVdB2013} では、
$\omega\in E[2]$、$i\in\mathbb N$ に対して、$Q_{\omega+i\tau}$ 上の
単純対象 $F_{\omega+i\tau}$ を構成し、
同文献 Proposition 9.3 の Euler 形式を用いて
$\chi(F_{\omega+i\tau},F_{\omega+i\tau})=2$ を計算している。
同文献の (10.1) でこの単純対象は coherent な line object の商として与えられるので、
$\coh Q_{\omega+i\tau}$ に属する。
ここで $\omega=0$、$i=1$ とすれば $\cC_\tau$ に単純対象 $F$ が存在して
\[
 \chi(F,F):=\sum_{j=0}^2(-1)^j\dim_{\mathbb C}\Ext^j(F,F)=2
\]
となる。
一方、$\coh(\PP^1\times\PP^1)$ の単純対象は閉点の skyscraper 層 $k(p)$ であり、
滑らかな二次元局所環の Koszul 分解からその自己 Ext の次元は $(1,2,1)$ である。
従って $\chi(k(p),k(p))=0$ となる。
exact 同値は単純性と全 Ext を保つので、両圏は同値でない。
\end{proof}

この例は可換二次曲面の同じ Abel 圏に異なる平方根や twisted coordinate ring を
選ぶことでは説明されない。
一方、この命題は全ての滑らかな Sklyanin pencil の成員について非同値性を主張するものでも、
楕円型の全対象を分類するものでもない。
引用した Sklyanin 環・pencil・単純対象の構成は既知の結果であり、
本節はそれらが本稿の標識付き定義、とくに全次数の切断環と
Gabriel--Krull rank 条件を満たすことを確認したものである。

% !TEX root = ../nc-quadric-surface.tex
\section{反 Serre 環による同値判定}
\label{sec:antiserre-ring}

\begin{theorem}[反 Serre 環と pointed exact 同値]
\label{thm:antiserre-invariant}
$(\cC,\OO)\in\HH$ に対し、$t=S^{-1}[2]$ の heart 上の代表を用い、
\[
 R(\cC,\OO)=\bigoplus_{n\ge0}\Hom(\OO,t^n\OO)
\]
に定義~\ref{def:ample} の右規約で積を入れる。
その次数付き同型類は代表、Serre 関手および平方根の選択に依存しない。
任意の適合根の切断環 $B$ に対し $R\simeq B^{(2)}$ であり、
$(\cC,\OO,t)\simeq(\qgr R,\pi R,(1))$ となる。
二つの $\HH$ 対象が pointed exact 同値であるための必要十分条件は
反 Serre 環が次数付き $k$-代数として同型となることである。
\end{theorem}

\begin{proof}
$t$ は導来圏上で $s^2$ と自然同型なので heart を保ち、制限は exact 自己同値である。
heart から導来圏への埋込みの充満忠実性により自然同型を制限できる。
一般に $F:\cC\to\cC'$ が $k$-線形同値で、
$u:F\OO\xrightarrow{\sim}\OO'$ および $\eta:Ft\xrightarrow{\sim}t'F$
があるとする。反復同型 $\eta_n:Ft^n\to t'^nF$ を
$\eta_0=1$、$\eta_{n+1,X}=t'(\eta_{n,X})\circ\eta_{t^nX}$ で定める。
この定義から
$\eta_{m+n,X}=t'^n(\eta_{m,X})\circ\eta_{n,t^mX}$ となる。
次数 $m$ で
\[
 a\longmapsto t'^m(u)\circ\eta_{m,\OO}\circ F(a)\circ u^{-1}
\]
と置くと、$\eta_n$ の自然性を $a:\OO\to t^m\OO$ に適用して
積 $t^n(a)\circ b$ と整合することが分かる。
$F$ の充満忠実性により各次数で線形同型である。

この構成を恒等関手と自然同型 $s^2\simeq t$ に適用すると
$R\simeq B^{(2)}$ を得る。Serre 関手の一意性から代表の変更にも同じ構成を使える。
pointed exact 同値 $F$ の導来同値は Serre 関手を自然同型まで移すので、
heart 上で $Ft\simeq t'F$ となり、必要性が従う。

$B^{(2)}$ は左右 noetherian である。実際、右イデアル $I$ を $B$ へ拡張すると
$IB\cap B^{(2)}=I$ である。積の偶数成分への射影は $I$ の元との左乗法と可換であり、
$IB$ の任意の元の偶数成分は $I$ に入る。左イデアルにも同様に適用できる。
補題~\ref{lem:ample-square-from-generation} により $(\OO,t)$ は ample なので、
\cite[Theorem 2.6(2)]{MU2021} が主張した triple の再構成を与える。
逆に次数付き環同型は右加群の移送を通じて $\pi R$ を保つ $\qgr$ の exact 同値を与える。
\end{proof}

\begin{proposition}[反 Serre 環の二次関係]
\label{prop:antiserre-quadratic}
$B=T(V)/(K)$ を適合根に対応する切断環とし、
$U=B_2=(V\otimes V)/K$ と書く。このとき
\[
 R\simeq T(U)/(J),\qquad
 J=\ker(U\otimes U\to B_4)
   =\operatorname{im}(V\otimes K\otimes V\to U\otimes U).
\]
$\dim U=9$、$\dim J=56$ であり、
$\dim R_n=(2n+1)^2$、$H_R(z)=(1+6z+z^2)/(1-z)^3$ となる。
\end{proposition}
\begin{proof}
$T(V)^{(2)}=T(V\otimes V)$ と同一視する。
次数 $2n$ での関係イデアルは $2n$ 文字の中に $K$ を一箇所挿入した部分空間の和である。
$K$ の前の文字数が偶数なら $K$ は二文字一組の内部にあり、一次関係 $K$ が生成する。
奇数なら挿入箇所の前後に各一文字を加えて二組を取り出せるため、
$V\otimes K\otimes V$ が生成する。最初の関係で割った後は $T(U)$ となり、
残る生成関係は次数二の上記の像である。
命題~\ref{prop:reconstruction} より $\dim B_2=9$、$\dim B_4=25$ だから
$\dim J=81-25=56$。同じ Hilbert 関数の偶数項を足せば級数を得る。
\end{proof}

従って二つの適合表示から計算した関係空間 $J,J'$ について、
$(h\otimes h)J=J'$ を満たす線形同型 $h:U\to U'$ の存在が
pointed exact 同値の必要十分条件である。
この判定では $B$ の基底や被覆を保持する必要はない。
体上で実現する軌道は定理~\ref{thm:complete-h2-classification} で全て記述する。
これを族のモジュライへ拡張するには、関係空間の実現部分を底変更と整合する
代数的 locus として構成し、相対再構成を別途証明する必要がある。

\begin{corollary}[導来圏と構造対象からの heart の認識]
\label{cor:pointed-derived-exact}
二つの $\HH$ 対象の有界導来圏の $k$-線形三角同値が
$\OO$ を $\OO'$ に送るなら、その同値は標準の heart に関して $t$-exact である。
\end{corollary}
\begin{proof}
$t=S^{-1}[2]$ とする。任意の $M\in\cC$ に対し、
$t^{-n}\OO$ から $M$ への正次数 Ext は $n\gg0$ で消える。
実際、豊富性により $M$ への全射の始域を有限個の $t^j\OO$ の直和に取れる。
従って $\Hom(M,t^{-n-1}\OO)=0$ は負切断の消滅から従い、
Serre 双対性で $\Ext^2(t^{-n}\OO,M)=0$ を得る。
$\Ext^1$ の消滅は二次の正則正規商 $B$ の AS Gorenstein 性を用いる
\cite[Lemma 2.4]{MU2022}。
深さ三により $\Ext^1_{\qgr B}(\pi B(i),\pi B(j))=0$ である
（\cite[Lemma 5.9(2)]{MU2021} と同じ打切りの議論）。
上の全射の核を $K$ とすると、$\Ext^1(t^{-n}\OO,M)$ は
$\Ext^2(t^{-n}\OO,K)$ に埋まり、後者は既に十分大きい $n$ で消える。

従って $L\in D^b(\cC)$ に対し、十分大きい $n$ では
$\Hom(t^{-n}\OO,L[i])\simeq\Hom(t^{-n}\OO,H^i(L))$ である。
これは標準打切りに伴う Ext のスペクトル系列から従う。
豊富性の生成条件により、非零な $H^i(L)$ への非零射は
任意に大きい $n$ のある $t^{-n}\OO$ から存在する。
従って $L$ が heart に属する条件は、全ての $i\ne0$ について
$\Hom(t^{-n}\OO,L[i])$ が $n\gg0$ で消えることである。
三角同値は Serre 関手とシフトを保つので、$\OO$ を保てば
この条件を双方向に保つ。
\end{proof}

\subsection{可換 quadric と楕円型 Sklyanin pencil の判別}

可換 quadric の四つの根が与える反 Serre 環はいずれも
$\bigoplus_{n\ge0} H^0(Q,\OO_Q(2n,2n))$ である。
これは $t=-\otimes\OO_Q(2,2)$ から直接従い、
命題~\ref{prop:ore-ambient-family} の被覆の選択にも依存しない。

\begin{proposition}[滑らかな Sklyanin 例の内在的な楕円商]
\label{prop:sklyanin-antiserre-center}
命題~\ref{prop:sklyanin-h2} の範囲で $R=B_z^{(2)}$ とする。
pencil の残りの中心二次元 $h$ の像を $R_1$ で考えると
$Z(R)=k[h]$ であり、$R/(h)=(B_z/(h))^{(2)}$ となる。
特にこの範囲の全対象は可換 quadric と pointed exact 同値でなく、
商に現れる楕円曲線は pointed exact 同値類から回復できる。
\end{proposition}
\begin{proof}
$h$ は中心正則である。偶数成分を取ると主張した商が得られる。
既知の楕円商表示（命題~\ref{prop:sklyanin-h2} の証明）により
この商は楕円曲線の twisted homogeneous coordinate ring の反対環である。
偶数次だけを取ると、平行移動 $\sigma$ は $\sigma^2$ に、
次数一の線束 $L$ は $L\otimes\sigma^*L$ に替わる。
$\sigma^2$ は無限位数で、この線束の次数は八である。

このような楕円環の中心は $k$ である。実際、線束を関数体 $k(E)$ 上で
自明化すると、環は skew polynomial ring $k(E)[u;\sigma^2]$ に入る
（反対環を取っても中心は同じである）。
中心に非零斉次元 $f u^m$（$m>0$）があれば、
次数一の任意の二つの非零切断 $a u,b u$ との可換性を割り算して
$(\sigma^2)^m(a/b)=a/b$ を得る。
次数八の線束は very ample なので、その比で関数体が生成される。
これは $\sigma^2$ の無限位数性に反する。
従って $R/(h)$ の中心は $k$ であり、$R$ の中心にある任意の正次数斉次元は
$h$ の倍数となる。$h$ が中心正則なので割った元も中心にあり、次数による帰納法で
$Z(R)=k[h]$ が従う。

従って $R$ は非可換であり、定理~\ref{thm:antiserre-invariant} と
可換 quadric の $R$ の表示から、pointed exact 非同値である。
さらに $Z(R)_1$ の生成するイデアルで割る操作は内在的であり、
商の非零斉次元を逆にした次数付き商環の零次体は $k(E)$ となる。
実際、商は skew Laurent 環（またはその反対環）に埋め込まれ、
次数一切断の比が $k(E)$ を生成する。
斉次 Ore 局所化は noetherian domain 性から存在する。
斉次元同士の共通倍数は非零な斉次成分を取って選べる。
従って $k(E)$ は同型不変であり、その滑らかな射影模型 $E$ も同型まで定まる。
さらに商環の次数一の可逆元による零次体上の共役作用は、
零次体が可換なので元の選択に依存しない。この作用は上の表示で
$\sigma^2$（反対環の規約ではその逆）に対応する。
\end{proof}

% !TEX root = ../nc-quadric-surface.tex
\section{中心 standard 表示との比較}
\label{sec:standard-comparison}

\begin{proposition}[標識とシフトの移送]
\label{prop:standard-pointed}
$B=A/(g)$ を $\HH$ 対象の適合表示とし、$g$ は中心元、
$B$ は \cite{MN2025} の意味で standard とする。
$X$ を二つの ruling に対応する極大 Cohen--Macaulay 加群の一つとする。
すると \cite[Theorem 6.13]{MN2025} の三次元 AS 正則
$\mathbb Z$-algebra $D$（Gorenstein parameter 四）への同値 $F$ は
\[
 F(\pi B(i))\simeq\pi e_{-2i}D,\qquad
 F(\pi X(i))\simeq\pi e_{-2i-1}D
\]
を満たす。特に構造対象は $\pi e_0D$ に移る。
$D$ の二添字の周期移動から得る自己同値 $T$ に対して
$F(1)\simeq TF$ であり、$T^2\simeq S_D^{-1}[2]$ となる。
\end{proposition}

\begin{proof}
命題~\ref{prop:reconstruction} により $B$ は domain であり、
命題~\ref{prop:gldim-two} により $\qgr B$ の大域次元は有限である。
よって中心性と standard 性を加えれば
\cite[arXiv v1, Theorem 6.13]{MN2025} の仮定を満たす。
その構成は \cite[Theorem 5.16]{MU2021} の列
$E_{2i}=\pi B(i)$、$E_{2i+1}=\pi X(i)$ を用いる。
$D_{ij}=\Hom(E_{-j},E_{-i})$ の積は合成であり、
$j$ 成分が $\Hom(E_{-j},M)$ の右 $D$-加群を tails に送ると $F$ になる
\cite[arXiv v1, Definition 5.1, Theorems 5.8 and 6.4]{MN2025}。
従って $F(E_{-i})\simeq\pi e_iD$ である。

シフト $s=(1)$ は $s E_i\simeq E_{i+2}$ を与える。
これらの同型を使い、射に $s$ を作用させることで
$\alpha:D_{ij}\xrightarrow{\sim}D_{i-2,j-2}$ を得る。
右加群上で $(TN)_j=N_{j+2}$ とし、作用を $\alpha^{-1}$ で引き戻す。
写像
\[
 \Hom(E_{-j-2},M)\longrightarrow\Hom(E_{-j},sM),
 \qquad f\longmapsto s(f)
\]
（始域には上記の同一視を用いる）は合成と可換であり、
$TF\simeq Fs$ を与える。導来同値による Serre 関手の移送から
$T^2\simeq S_D^{-1}[2]$ が従う。

具体的に $p\le q$ に対して $B_{q-p}\to D_{2p,2q}$ は
$a\mapsto s^{-q}(a)$ であり、$p\le q\le r$ での合成は
$s^{-q}(a)\circ s^{-r}(b)=s^{-r}(s^{r-q}(a)\circ b)$ となる。
従って偶数添字部分は右規約の $\check B$ であり、
$R_n$ は $D_{0,4n}$ に対応する。積は周期移動と合成で与えられる。
\end{proof}

命題~\ref{prop:sklyanin-h2} の滑らかな Sklyanin 例では
\cite[Theorem 5.5]{MUW2026} が standard 性を保証するので、
命題~\ref{prop:standard-pointed} を適用できる。
両方の ruling を別々に用いて得た比較は、どちらも同じ $(\qgr B,\pi B)$ を回復する。
逆向きに任意の cubic $\mathbb Z$-algebra から出発した場合、
二添字の周期移動だけから $T^2\simeq S_D^{-1}[2]$ を結論することはできない。

% !TeX root = ../nc-quadric-surface-classification.tex

\subsection{三次正則表示を与える ruling}
\label{subsec:cubic-ruling}

\begin{proposition}[三次型の判定]
\label{prop:cubic-ruling}
命題~\ref{prop:standard-pointed} の中央標準表示 $B$ に対し、
二つの次数零生成 MCM 加群の一方 $X$ を選び、他方を $Y$ とする。
$E_{2n}=\pi B(n)$、$E_{2n+1}=\pi X(n)$ から作る
$D_{ij}=\Hom(E_{-j},E_{-i})$ は、次の積が同型である場合にちょうど
Van den Bergh の三次三次元 AS 正則 $\mathbb Z$-algebra となる。
\[
 \mu_X:\Hom(\pi X,\pi B(1))\otimes_k\Hom(\pi B,\pi X)
       \longrightarrow \Hom(\pi B,\pi B(1)),\qquad h\otimes x\longmapsto h\circ x .
\]
\end{proposition}

\begin{proof}
必要性は \cite[arXiv v3, \S 5.1, (Q2)--(Q3)]{VdB2011} の
距離二の積に関する記述から従う。三つの Hom はそれぞれ次元 $2,2,4$ である
\cite[Lemmas 5.6, 5.9]{MU2021}。

十分性を示す。$X(n)$ を標的とする射について、
\[
 0\longrightarrow Y(n-1)\longrightarrow B(n)^2\longrightarrow X(n)
   \longrightarrow0
\]
を使う \cite[Proposition 5.5(1)]{MU2021}。
tails で $E_j$（$j<2n+1$）から $\Hom$ を取ると、
$\Ext^1(E_j,\pi Y(n-1))=0$ なので、各射は
直前の $\pi B(n)$ を経由する射の和となる。消滅は
\cite[Lemmas 5.6(8), 5.9(2), 5.11]{MU2021} による。
$B$ の一次生成性と $\mu_X$ の全射性から、$\pi B(a)$ から $\pi B(b)$
（$b>a$）への射は、隣接する helix 対象を介した合成で張られる。
残る $\pi X(a)$ から $\pi B(b)$（$b>a$）への射に対しては、もう一方の列
$0\to X(-1)\to B^2\to Y\to0$ を $\Hom_B(-,B)$ で双対化する。
$Y$ は MCM なので $B^2\to\Hom_B(X,B)(1)$ は全射となる。
ここで次数付き内部 Hom を用いる。従って $X$ から $B$ への内部 Hom は、
次数一の二つの射で左 $B$-加群として生成され、その左作用は
$B$ のシフト間の射との合成である。tails との Hom の比較
\cite[Lemma 5.9(1)]{MU2021} を使えば、全ての射が隣接する helix 対象を
介した合成で張られる。

よって $D$ は隣接成分で生成され、各隣接成分は二次元である。
右単純加群の極小分解の初めは $P_{i+1}^2\to P_i$ であり、
左単純加群でも同じ隣接生成性を使える。
Mori--Nyman の AS 条件で大域次元 $3$、Gorenstein parameter $4$ なので、
$P_i=e_iD$ とすると右単純加群の最後の項は $P_{i+4}$ である。
双対複体が添字 $i+4$ の左単純加群の極小分解になることから、
残る項は $P_{i+3}^2$ である
\cite[arXiv v1, Theorem 4.5 and Remark 4.6]{MN2025}。
$D$ の成分次元は \cite[Lemmas 5.6, 5.9]{MU2021} の
Hilbert 級数から添字の差に関して二次で抑えられる。
これは \cite[arXiv v3, Definitions 4.1.1--4.1.2]{VdB2011} の三次正則型を与える。
\end{proof}

\begin{proposition}[Sklyanin pencil での判定]
\label{prop:sklyanin-cubic-ruling}
命題~\ref{prop:sklyanin-h2} の滑らかな対象に対し、
ruling のラベルを $w\in\{z,-z-2\tau\}$ と書く。
この ruling に属する線加群を $\operatorname{coker}(X(-1)\to B_z)$
と表示して得られる $X$ に対し、
\[
 \rank\mu_X=
 \begin{cases}
  4 & w\notin E[2],\\
  3 & w\in E[2].
 \end{cases}
\]
従って少なくとも一方の選択で、標識と平方根を保持した三次正則表示が得られる。
\end{proposition}

\begin{proof}
命題~\ref{prop:standard-pointed} の仮定は
\cite[Lemma 5.1 and Theorem 5.5]{MUW2026} と命題~\ref{prop:sklyanin-h2} により満たされる。
右加群での ruling と secant の対応には
\cite[Lemma 3.20, Proposition 3.21 and Lemmas 5.3--5.4]{MUW2026} を用いる。
ここでは補題~\ref{lem:sklyanin-opposite} の次数付き同型で
$B_z$ を通常の右加群表示 $T/(g_z)$ に移してから $w$ を読む。
固定した ruling は、$p+q=w$ を満たす secant $pq$ の線加群の族である。
実際、同じラベルの二つの線加群は Lemma 5.4 の列で共通の補線を取り、
異なるラベルの二つには同じ補題の列を直接適用できる。
標準性と Proposition 3.21 により、同じラベルと同じ ruling が対応する。

異なる一般の二つの secant $pq,p'q'$ を取り、四点は相異なるとする。
対応する線加群を定める $X(-1)\to B_z$ の二射は
二次元の Hom 空間の異なる一次元部分空間を与えるので、基底をなす。
それぞれの次数一像 $W,W'\subset (B_z)_1$ は secant 上で消える
線型形式の二次元空間であり、$\operatorname{im}\mu_X=W+W'$ である。
$W\cap W'\ne0$ は四点が同一平面上にあることと同値である。
楕円曲線の次数四の埋込みでは、四点が同一平面上にあることと
$p+q+p'+q'=0$ が同値である \cite[\S 10.1]{SVdB2013}。
従って共通部分が非零となるのは $2w=0$ の場合であり、
二つの secant は相異なるため、その場合の共通部分は一次元となる。
両ラベルが $E[2]$ に属すると $4\tau=0$ となり、
$\tau$ の無限位数に反する。
命題~\ref{prop:cubic-ruling} と命題~\ref{prop:standard-pointed} から結論を得る。
\end{proof}

$z=0$ は $\tau$ が無限位数のとき滑らかな例であり、
$w=0$ の選択では距離二の積が階数三、$w=-2\tau$ の選択では階数四である。
同じ構造対象と平方根から出発しても、どちらの ruling を取るかで
三次正則型の判定結果が変わる。


\subsection{制約を保つ Zhang twist}

\begin{lemma}[次数付き自己同型が tails 上で自明となる条件]
\label{lem:scalar-tails}
適合切断環 $B$ の次数付き自己同型 $\sigma$ に対し、
右作用の引き戻し $(-)_\sigma$ が $\qgr B$ 上で恒等と自然同型となるための
必要十分条件は、ある $\lambda\in k^\times$ に対して
$\sigma(b)=\lambda^d b$ ($b\in B_d$) となることである。
この次数スカラーを $\xi_\lambda$ と書く。
\end{lemma}
\begin{proof}
$\pi B(i)_\sigma\to\pi B(i)$ の基準同型は $b\mapsto\sigma^{-1}(b)$ である。
$\End(\pi B(i))=k$ より、恒等への任意の自然同型の $\pi B(i)$ 成分は
この基準同型の非零定数 $c_i$ 倍である。
$a\in B_1$ の左乗法 $\pi B(i)\to\pi B(i+1)$ に対する自然性から
$c_{i+1}\sigma^{-1}(a)=c_i a$ を得る。ここで左乗法は
$B_1\to\Hom(\pi B(i),\pi B(i+1))$ で単射である。
実際、非零 $a$ の左乗法は domain 性により非零の十分高い次数においても単射であり、
tails で零射にならない。一次生成により $\sigma$ は次数スカラーとなる。
逆に $\sigma=\xi_\lambda$ なら次数 $j$ の成分に $\lambda^{-j}$ を掛けることで
$M_\sigma\to M$ の自然同型が得られる。
\end{proof}

\begin{proposition}[固定表示での中心化と平方根性]
\label{prop:central-twist-obstruction}
適合表示 $B=A/(g)$ に対し、$ag=g\nu(a)$ で正規化自己同型 $\nu$ を定める。
$\sigma\in\Aut_{\mathrm{gr}}(A)$ が $\sigma(g)=c g$ を満たすとし、
積 $a\star b=a\sigma^{\deg a}(b)$ の Zhang twist を取る。
twist 後に $g$ が中心元となる条件は $\sigma^2=\xi_c\nu$ である。
一方 twist 後の商の次数シフトが反 Serre の平方根となる条件は
$\bar\sigma^2$ が $B$ の次数スカラーとなることである。
従って両条件が成り立つなら $\nu=\xi_1$ または $\nu=\xi_{-1}$ である。
\end{proposition}
\begin{proof}
$g$ の正則正規性が $\nu$ の存在と一意性を与える。
$a\in A_d$ に対し $a\star g=c^d g\nu(a)$、$g\star a=g\sigma^2(a)$
なので、中心性の判定は左から $g$ を消去すればよい。
これは $g$ の直線を保つ場合の
\cite[Remark 4.4]{MUW2026} の正規化公式である。

右加群の twist 同値 $F$ は作用を、次数 $j$ の $m$ に対して
$m\star b=m\bar\sigma^j(b)$ と変える。これは regular module を保つ。
$F(M)(1)$ の作用では $M(1)$ における次数 $j$ の元に元の次数 $j+1$ を使うので、
$F(M)(1)\simeq F((M(1))_{\bar\sigma})$ となる。
従って新しいシフトを $\qgr B$ に戻すと $s(-)_{\bar\sigma}$ であり、
その二乗は $s^2(-)_{\bar\sigma^2}$ となる。
$s^2\simeq S^{-1}[2]$ と Serre 関手の移送、および
補題~\ref{lem:scalar-tails} により平方根性の条件を得る。
両条件下で $\nu$ は $B$ 上で次数スカラー、従って
$A_1=B_1$ と一次生成性から $A$ 上でも次数スカラーである。
正規化の定義から $\nu(g)=g$ なので、その一次のスカラーは $\pm1$ である。
\end{proof}

\begin{example}
\label{ex:ore-central-twist}
命題~\ref{prop:ore-ambient-family} の $A_a$ では
$\nu(x)=a^{-1}x$、$\nu(y)=ay$、$\nu(u)=u$、$\nu(v)=v$ である。
$b^2=a$ を満たす $b$ を選べば
$(x,y,u,v)\mapsto (b^{-1}x,by,u,v)$ は関係と $g$ を保ち、
二乗が $\nu$ なので中心化できる。しかし $a\ne1$ なら $\nu$ は次数スカラーでなく、
この表示からはどの次数付き自己同型 twist でも
中心性とシフトの平方根性を同時には得られない。
同じ標識対象には $a=1$ の中心表示が存在するため、
中心表示の存在一般に対する反例は与えない。
\end{example}

\subsection{二添字周期が平方根となる条件}

\begin{proposition}
\label{prop:periodicity-gauge}
$D$ を命題~\ref{prop:standard-pointed} で構成した $\mathbb Z$-algebra とする。
より一般に \cite[Theorem 6.8]{MN2025} の coherent $\mathrm{ASF}^{++}$ 正則性の仮定を満たす
次元三・Gorenstein parameter 四の連結な $D$ で、
$D_{ij}\to\Hom(\pi e_jD,\pi e_iD)$ が同型である場合にも以下は成立する。
Nakayama 同型を $\nu:D_{ij}\to D_{i-4,j-4}$ とする。
二添字移動 $\alpha:D_{ij}\to D_{i-2,j-2}$ から得る自己同値 $T_\alpha$ が
$T_\alpha^2\simeq S_D^{-1}[2]$ を満たす必要十分条件は、
ある $c_i\in k^\times$ に対し
\[
 (\nu^{-1}\alpha^2)(a)=c_i c_j^{-1}a\qquad (a\in D_{ij})
\]
となることである。
また $\pi e_0D$ の反 Serre 環は次数 $n$ で $D_{-4n,0}$、
積は $a*b=\nu^n(a)b$ ($a\in D_{-4m,0}$、$b\in D_{-4n,0}$) となる。
\end{proposition}
\begin{proof}
命題~\ref{prop:standard-pointed} の $D$ に対する仮定は
\cite[Theorems 6.4, 6.13 and Proposition 4.11]{MN2025} による。
\cite[Theorem 6.8]{MN2025} の右規約では
$S_D^{-1}[2](N)=N(4)_{\nu^{-1}}$ であり、
$\pi e_iD$ を $\pi e_{i-4}D$ に移し、左乗法の射には $\nu$ を作用させる。
$T_\alpha^2$ は同じ添字移動を行い、射には $\alpha^2$ を作用させる。
二つの関手の比は添字を固定する自己同型 $\delta=\nu^{-1}\alpha^2$ による
作用の引き戻しである。
$\End(\pi e_iD)=k$ と標準 Hom 同型の下で
各左乗法に自然性を適用すると、$\delta$ が成分ごとに $c_i c_j^{-1}$ 倍となる。
逆にこの形なら、作用を $\delta^{-1}$ で引き戻した右加群から元の加群への同型は
次数 $j$ の成分で $c_j^{-1}$ を掛けて与えられる。従って tails 上でも自然同型を得る。

積は $S_D^{-1}[2]$ を $n$ 回左乗法に作用させてから合成するので、
表示された公式となる。次数 $n$ ごとに $\nu^{-n}$ で移すと
$D_{0,4n}$ 上で積は $a*b=a\nu^{-m}(b)$ となる。
特に適合表示が存在する場合、二次関係空間は
$D_{0,4}\otimes D_{0,4}\to D_{0,8}$、$a\otimes b\mapsto a\nu^{-1}(b)$
の核である。
\end{proof}

% !TEX root = ../nc-quadric-surface.tex
\section{正規切断の標識付き比較}
\label{sec:normal-comparison}

\begin{theorem}[正規切断からの四周期 helix]
\label{thm:normal-helix}
$\HH$ の任意の適合表示を $B=A/(g)$ とする。
二つの次数零生成非自由極大 Cohen--Macaulay 型を $X,Y$ とし、
ラベル集合 $\{X,Y\}$ の交換を $c$、$Z\mapsto\Omega Z(1)$ の置換を $p$ と書く。
任意の $X_0\in\{X,Y\}$ から $X_{i-1}=pc(X_i)$ ($i\in\mathbb Z$) と定めると、
\begin{equation}
 E_{2i}=\pi B(i),\qquad E_{2i+1}=\pi X_i(i)
\end{equation}
は ample な四周期の full geometric helix をなす。
従って $D_{ij}=\Hom(E_{-j},E_{-i})$ ($i\le j$) から
右 noetherian な次元三・Gorenstein parameter 四の AS 正則
$\mathbb Z$-algebra を得る。標識は $\pi e_0D$ に移り、
反 Serre 関手は $s^2$ から誘導される四添字の移動に移る。
standard の場合は $X_i$ が一定であり、non-standard の場合は二型が交互に現れる。
\end{theorem}

\begin{proof}
正則正規切断に対する Rees の補題により $B$ は次元三・parameter 二の
AS Gorenstein 環である
\cite[Lemma 2.4]{MU2022}。parameter の変化は切断次数二による。
$B$ は domain なので $g$ は二つの次数一元の積に分解できず、
\cite[Definitions 4.10 and 5.9]{MU2022} の $n=4$ の high rank 条件を満たす。
$A$ は同論文の quantum polynomial algebra であり、$\qgr B$ の大域次元は二である。
従って \cite[Theorem 5.5, Lemma 5.13(4)]{MU2022} により
$C(B)\simeq M_2(k)\times M_2(k)$ となる。
ここで用いる定義は AS 正則性と Hilbert 級数を要求し、
Auslander 正則性と Cohen--Macaulay 性の追加条件を課さない。
同論文の Lemma 4.13(5) により、
安定圏の次数零生成部分の二つの非零単純型に対応する非自由加群を $X,Y$ と取れる。
Lemma 5.11 の証明で、対応する単純加群の次元二は
線形行列分解の階数二に等しい。この分解から極小自由分解の
各項が $B(-j)^2$ となり、$\Omega Z(1)$ も次数零生成の非自由 MCM 型である。
$\Omega(1)$ は安定圏の自己同値なので二型を置換する。
自由直和因子を除いた極小代表を取っているため、これは加群の同型型の置換 $p$ を与える。
$p^2=1$ から $X_{i+2}=X_i$ である。
$Z\in\{X,Y\}$ の極小自由被覆は
\begin{equation}
 0\longrightarrow p(Z)(-1)\longrightarrow B^2\longrightarrow Z
 \longrightarrow0
 \label{eq:normal-ruling-evaluation}
\end{equation}
を与える。
$Z$ と $p(Z)$ を交換した同様の列を双対化すると、
内部 $\Hom_B(Z,B)$ は次数一に二元で生成され、特に次数零以下は零である。
ここでは $\Ext_B^1(p(Z),B)=0$ を使った。

\cite[Lemma 4.13(3)--(5)]{MU2022} の安定圏の導来同値は、
二型を半単純環の二つの単純加群に送る。
安定圏で $\Omega$ は懸垂の逆であり、内部シフト $(1)$ は
$p$ と懸垂 $[1]$ の合成となる。従って二型の安定 Hom は
内部シフト零でのみ非零になり、同じ型の場合に $k$ である。
次数零以下での $Z$ から自由加群を経由する射は、
$\Hom_B(Z,B(j))=0$ ($j\le0$) と $Z'$ の負次数の消滅から零となる。
従って未シフトの二型の次数零 Hom も $\delta_{ZZ'}k$ である。
また $\Ext_B^1(Z,Z'(d))$ は安定 Hom
$\underline{\Hom}_B(p(Z),Z'(d+1))$ であり、
$d=-1$ かつ $Z'=p(Z)$ の場合にのみ非零となる。

自由加群と二型はいずれも深さ三である。実際、
\eqref{eq:normal-ruling-evaluation} を $p(Z)$ にも適用して
$Z$ を自由加群へ埋め込む列を得る。その余核も二型のシフトである。
この操作を繰り返すと、自由加群の深さ三と長完全列から
$\Ext^j_B(k,Z)=0$ ($0\le j<3$) が従う。
従って $\Hom$ と $\Ext^1$ は、\cite[Lemma 5.9(1)--(2)]{MU2021}
の打切りの議論により graded 加群から tails に移しても変わらない。
$\Ext^1$ はどちらかが自由加群なら零である。
奇数項同士では前向きの射の $\Ext^1$ は零であり、
四連続項の中の逆向きでは、$X_{i-1}=pc(X_i)\ne p(X_i)$
により零である。$\Hom$ の逆向き消滅は次数と
$\Hom_B(Z,B)$ の次数一以上への集中から従う。
$\Ext^2(E_i,E_j)$ は Serre 双対性により $\Hom(E_j,E_{i-4})$ の双対であり、
$i\le j$ または $0<i-j<4$ のとき逆向きの Hom なので零となる。
高次は大域次元から零である。
この結果と $E_{i+4}\simeq s^2E_i$ が geometric helix の条件を与える。
ここで $s^{-2}$ は導来圏上で $S[-2]$ となる exact 自己同値なので、
この圏の canonical bimodule を与える。

完全性を示す。Orlov の分解と MCM 安定圏の記述により、
$D^b(\qgr B)$ は $\pi B(-1),\pi B,\pi X,\pi Y$ で生成される
\cite[Theorem 2.5(i)]{Orlov2009}。
実際、$\Hom_B(Z,B(j))=0$ ($j\le0$) なので、
Orlov の埋込みは二型の特異圏での像を $\pi X,\pi Y$ に送る
\cite[Lemma 4.1]{MU2016Stable}。二型は半単純な安定圏の導来表示を生成する。
これを任意に次数シフトする。$Y_i=c(X_i)$ と書くと、
\eqref{eq:normal-ruling-evaluation} の $Y_i(i)$ の核が $X_{i-1}(i-1)$ なので、
生成列 $(\pi B(i-1),\pi B(i),\pi Y_i(i),\pi X_i(i))$
の中の隣接二項を対応する三角で入れ替えて、偶数添字で始まる四項の完全性を得る。
奇数添字で始まる場合は、補型 $c(X_i)$ を同じ方法で選んだ
偶数添字から始まる四項
$(\pi B(i+1),\pi Y_{i+1}(i+1),\pi B(i+2),\pi Y_{i+2}(i+2))$
で、二つの隣接対をそれぞれ同じ三角で置き換える。

豊富性の生成条件は偶数項から従う。
残る条件に必要な $\Ext^1(\pi Z(-n),K)=0$ ($n\gg0$) を示す。
任意の $K\in\qgr B$ について、
$0\to L\to P\to K\to0$ を短完全列とし $P$ を有限個の $\pi B$ のシフトの直和とする。
$\Ext^1(\pi Z(-n),P)=0$ なので $\Ext^1(\pi Z(-n),K)$ は
$\Ext^2(\pi Z(-n),L)$ に埋まる。
同様の自由被覆 $P'\twoheadrightarrow L$ と大域次元二を用いると、
$\Ext^2(\pi Z(-n),L)$ は $\Ext^2(\pi Z(-n),P')$ の商となる。
Serre 双対性と負次数の消滅から
$\Ext^2(\pi Z(-n),P')=\Hom(P',\pi Z(-n-2))^*=0$ ($n\gg0$) となる。
よって十分負の奇数項からの $\Hom$ は全射を保つ。
以上で \cite[Definition 5.6]{MN2025} の豊富性が従う。

\cite[Theorems 5.8, 6.4 and 6.5]{MN2025} を適用する。
構造対象と反 Serre 関手の移送は命題~\ref{prop:standard-pointed}
と同じ Hom の構成で得られる。non-standard の場合、
ここで使う表示に二添字の移動があることは要求していない。
\end{proof}

\begin{corollary}[三次正則型からのずれ]
\label{cor:normal-cubic-defect}
定理~\ref{thm:normal-helix} の $D$ は隣接成分と、距離二の偶数添字間の成分で生成される。
選んだ列について
\[
 r_i=\rank\bigl(
 \Hom(\pi X_i,\pi B(1))\otimes\Hom(\pi B,\pi X_i)
 \longrightarrow\Hom(\pi B,\pi B(1))
 \bigr)
\]
と置くと、$r_i\in\{3,4\}$ であり、距離二の偶数添字間で
追加の生成元を $4-r_i$ 個要する。$D$ が Van den Bergh の
三次正則型となる条件は全ての $r_i$ が四となることである。
また $r_i$ は $i$ に依存しない。
\end{corollary}
\begin{proof}
$\pi B(i)$ を始域かつ $\pi B(j)$ を終域とする射は、$B$ の一次生成性から
距離二の積で生成される。
奇数項を終域とする射で始域がそれより前にある場合は、
\eqref{eq:normal-ruling-evaluation} のシフトに $\Hom$ を適用して
直前の偶数項を経由する射で張られる。
実際、核のシフトは始域のシフトより小さくない
（始域が同じシフトなら自由加群である）ので、核への $\Ext^1$ は零である。
奇数項を始域とし偶数項を終域とする射については、内部 $\Hom_B(X_i,B)$ の
次数一生成性を使う。よって主張した距離で生成される。
各隣接空間は二次元である。

下限 $r_i\ge3$ を示す。$Z=X_i$ とする。$Z$ と $p(Z)$ は
\eqref{eq:normal-ruling-evaluation} からともに自由加群へ埋め込める。
noetherian domain $B$ の商斜体上で正の rank を持ち、
同じ完全列から両者の rank の和は二なので、各 rank は一である。
従って任意の非零斉次射 $Z(-1)\to B$ は単射で、
その像は次数一の二次元部分空間で生成される。
一次独立な二射の像の次数一成分が一致したとすると、
次数一生成性から像全体が一致する。このとき二射は
同じ像への同型であり、$\End_{\mathrm{gr}}(Z)=k$ に反する。
従って二つの二次元部分空間の和の次元は三以上となる。
この和が上の積の像である。
残る距離二の偶数添字間の生成元の個数は、その積の余次元 $4-r_i$ となる。

距離二の奇数項間の Hom は四次元である。
実際、$X_{i+1}(i+1)$ の被覆列の核 $p(X_{i+1})(i)$ に対し、
$\Hom(\pi X_i(i),-)$ を適用する。$X_i\ne p(X_{i+1})$ による
次数零 Hom の消滅と $\Ext^1$ の消滅から、
直前の偶数項を経由する積が同型であることが分かる。

$D$ の添字 $j$ から始まる距離二の生成元の個数を $d_j$ とする
（奇数 $j$ では零）。右単純加群の極小分解は
\[
 0\longrightarrow P_{j+4}\longrightarrow
 P_{j+3}^{2}\oplus P_{j+2}^{d_{j+2}}\longrightarrow
 P_{j+1}^{2}\oplus P_{j+2}^{d_j}\longrightarrow P_j\longrightarrow S_j
 \longrightarrow0
\]
となる。右側の第一項は生成元から分かり、双対極小分解では同じ議論が
左単純加群に使える。左側から見た距離二の生成元の個数も、
同じ隣接積の余次元なので $d_{j+2}$ である。
分解の次数 $j+2$ 成分の Euler 標数を取ると、
$4-2\cdot2-d_j+d_{j+2}=0$、すなわち $d_j=d_{j+2}$ を得る。
従って $r_i$ は $i$ に依存しない。特に non-standard では
二つの MCM 型の階数が等しい。

全ての $r_i$ が四なら分解は三次型の形になる。
成分次元の多項式上界は、
$B(i)^2\twoheadrightarrow X_i(i)$ と各 $X_j(j)$ の線形自由分解による
Hilbert 級数 $2(1-t)^{-3}$ から従う。
逆に三次正則型なら距離二の積は全射である。
\end{proof}

\begin{remark}[階数三の行列表示]
\label{rem:rank-three-matrix}
階数三の場合、$\Hom_B(X_i,B(1))$ の基底と $(X_i)_0$ の基底を選ぶと、
合成の四成分は $A_1=B_1$ の元を成分とする行列 $\Phi$ になる。
双対評価列と線形行列分解において、これは残りの行列 $\Psi$ と
$\Phi\Psi=gI$ を満たす側の行列として取れる。
四成分の間の唯一の線形関係を表す $2\times2$ 係数行列は階数二である。
実際、もし階数一なら基底変換で $\Phi$ の一成分を零にできる。
この場合、零成分を含む行と積 $\Phi\Psi$ の対角成分を見ると
$g$ が二つの次数一元の積となり、$B$ の domain 性に反する。
従って行と列の基底変換でこの関係を trace が零という形にできる。
そのように取ると
\[
 \Phi=\begin{pmatrix} a&b\\ c&-a\end{pmatrix},
 \qquad a,b,c\in A_1\ \text{は線形独立}
\]
となる。階数三の領域では、この形と対応する $\Psi$ の積の式が
残りの階数および別表示との比較の出発点となる。
\end{remark}

\begin{corollary}[non-standard 型の正規化条件]
\label{cor:nonstandard-square-matrix}
$p(Z)=Z$ のとき、対応する評価行列を $M\in M_2(A_1)$ として
\[
 M^2=gC,\qquad \nu(M)=CMC^{-1},\qquad C\in\GL_2(k)
\]
を満たす形に取れる。ここで $\nu$ は $ag=g\nu(a)$ の規約の正規化自己同型である。
特に $g$ が中心なら $C$ はスカラーであり、$M$ を定数倍して $M^2=gI$ とできる。
\end{corollary}
\begin{proof}
\eqref{eq:normal-ruling-evaluation} の射を
$\iota:Z(-1)\to B^2$ と $q:B^2\to Z$ とする。
積 $\iota q(-1):B(-1)^2\to B^2$ を $\bar M$ と書くと、
$q\iota=0$ により $\bar M^2=0$ となる。
これは $Z$ の極小線形表示であり、
それを与える行列分解の第一行列は、自由加群の基底変換の後、
$\bar M$ の $A_1$ への一意な持上げ $M$ と一致する。
従って $M$ は $A$ 上で単射である。
$\bar M^2=0$ と $\ker(A_2\to B_2)=kg$ から $M^2=gC$ となる。
$C$ が非可逆なら非零の定数列 $v$ で $Cv=0$ があり、
$M$ の単射性に反して $M^2v=0$ となる。
最後に $M(gC)=(gC)M$ に $Mg=g\nu(M)$ を代入し、
正則元 $g$ を左から消去して共役の式を得る。
中心切断では $C$ は $M$ の各係数行列と可換である。
系~\ref{cor:normal-cubic-defect} により係数行列の張る空間は三次元以上だが、
非スカラーな二次正方行列の中心化環は二次元なので、$C$ はスカラーである。
基礎体は代数閉なので必要な平方根を取り、$M$ を定数倍できる。
\end{proof}

% !TEX root = ../nc-quadric-surface.tex
\subsection{二型がともに階数三となる非標準表示}
\label{sec:rank-three-example}

\begin{proposition}
\label{prop:nonstandard-rank-three-example}
標数零の代数閉体上に、二つの MCM 型の行列成分がともに三次元を張る
non-standard な $\HH$ の適合表示が存在する。
この表示と定理~\ref{thm:normal-helix} のいずれの初期型を用いても、
得られる四周期表示は Van den Bergh の cubic 型にならない。
\end{proposition}

\begin{proof}
$k$ の元 $\mathrm i$ を $\mathrm i^2=-1$ となるように選び、
\[
 A=k[u,v]\otimes_k k\langle x,y\rangle/(yx+xy),\qquad
 g=\tfrac14(x^2+y^2)-uv,\qquad B=A/(g)
\]
とする。$A$ は四変数の skew polynomial ring であり、
命題~\ref{prop:commutative-twist-covers} で用いた skew Koszul 複体と同様に
型 $(1,4,6,4,1)$ である。$g$ は中央正則である。
$B$ は $k[u,v,x]$ 上で $y$ に関する Ore 拡大を取り、
モニックな関係 $y^2=4uv-x^2$ を課した環なので、
左 $k[u,v,x]$ 加群として基底 $1,y$ を持つ。
特に $u$ は正則である。$u$ を可逆にすると $v$ を消去でき、
\[
 B[u^{-1}]\simeq k[u,u^{-1}]\otimes_k k\langle x,y\rangle/(yx+xy)
\]
となるため、$B$ は domain である。

滑らかさは $B$ の Koszul 双対から判定できる。
$U,V,X,Y$ を $u,v,x,y$ の双対基底とすると、二次関係の直交補空間は
\[
\begin{gathered}
 U^2=V^2=0,\quad X^2=Y^2=w,\quad UV+VU=-4w,\quad XY=YX,\\
 UX=-XU,\quad UY=-YU,\quad VX=-XV,\quad VY=-YV
\end{gathered}
\]
を与える。ここで $w=X^2$ は中央であり、$B^!/(w)=A^!$ となる
正則元の代表である \cite[Lemma 4.11(2)]{MU2022}。$P=(U+V)/(2\mathrm i)$、$Q=(U-V)/2$ と置けば
$P^2=Q^2=w$、$PQ=-QP$ である。
$C(B)=B^![w^{-1}]_0$ は $P X^{-1},Q X^{-1},Y X^{-1}$ によって生成される。
最初の二つはそれぞれ平方が $-1$ で互いに反交換し、最後は中央で平方が $1$ である。
従って $M_2(k)\times M_2(k)$ から $C(B)$ への全射を得る。
\cite[Lemma 4.13(1)]{MU2022} により $C(B)$ は八次元なのでこの全射は同型であり、
同論文 Theorem 5.5 によって $\qgr B$ は滑らかである。

$B$ は正則二次切断なので次元三・parameter 二の AS Gorenstein 環である
\cite[Lemma 2.4]{MU2022}。canonical bimodule を与える右作用の twist を計算する。
\cite[arXiv v2, Example 5.5]{RRZ2014} により、$A$ の Nakayama 自己同型は
$u,v$ を固定し $x,y$ に符号 $-1$ を掛ける。
これは $g$ を固定するため同論文 Lemma 1.5 を適用でき、
$B$ の Nakayama 自己同型もその誘導 $\mu$ となる。
$\mu^{-1}=\mu$ なので、\cite[Example 3.6(2)]{MU2021} の
右 canonical bimodule の規約では
$S_{\qgr B}^{-1}[2]\simeq(-)_\mu(2)$ となる。

ここで graded 自己同型
\[
 \sigma(u)=-u,\qquad \sigma(v)=v,\qquad
 \sigma(x)=\mathrm i x,\qquad \sigma(y)=\mathrm i y
\]
を取る。$\sigma(g)=-g$ なので $B$ に誘導され、
その二乗は $\mu$ である。命題~\ref{prop:central-twist-obstruction} の証明中の
作用比較を用いると、Zhang 同値 $F:\qgr B\to\qgr B^\sigma$ は $\pi B$ を
$\pi B^\sigma$ に送り、新しい次数シフトを旧圏では $(-)_\sigma(1)$ として実現する。
その二乗は $(-)_\mu(2)$ なので、新しい次数シフトは反 Serre の平方根である。
$A^\sigma$ は型 $(1,4,6,4,1)$ を保ち \cite[Lemma 4.1]{MUW2026}、
$g$ はその中で左右正則正規であり、$B^\sigma=A^\sigma/(g)$ である。

$\pi B^\sigma$ に対する切断と豊富性を確認する。
正規正則切断の AS Gorenstein 性と深さ三から、
全次数での飽和性、$\operatorname{Ext}^1(\pi B^\sigma,\pi B^\sigma)=0$ が従う。
すなわち次数付き有限長加群 $L$ について $\operatorname{Ext}^j(L,B^\sigma)=0$
($j<3$) を用い、自由加群の打切りと Hom の極限を取る
\cite[Lemma 5.9(1)--(2)]{MU2021}。
$\operatorname{Ext}^2$ は Serre 双対性により負次数の切断の双対なので零である。
高次も大域次元二により零であり、$\End(\pi B^\sigma)=k$ と合わせて exceptionality を得る。
任意の対象は有限個の $\pi B^\sigma(-n)$ の直和の商であり、
その指数は打切りにより任意に大きく取れる。
全射の核 $K$ について $\operatorname{Ext}^1(\pi B^\sigma(-n),K)=0$
($n\gg0$) を得るには、$K$ の自由被覆を取って、
その核への $\operatorname{Ext}^2$ に埋め込む。
その核の自由被覆と大域次元二から、最後の群は
自由加群の像への $\operatorname{Ext}^2$ の商となり、Serre 双対性と
負次数の切断の消滅により零である。
従って $(\pi B^\sigma,(1))$ は ample であり、その切断環は $B^\sigma$ である。

Gabriel--Krull 条件は同値で不変なので $\qgr B$ で確認する。
命題~\ref{prop:sklyanin-h2} と同じく $A$ 上の有限自由分解により
Hilbert 関数は最終的に多項式となる。定数なら有限長、次数が高々一なら
$\cC_1$ に入ることは \eqref{eq:sklyanin-polynomial-filtration} と同じ議論で従う。
一方 $B/uB$ 上で $v$ は正則であり、$B/(u,v)$ は無限次元である。
従って $v$ の冪による列は $\pi(B/uB)$ の真の下降列となり、
この対象は有限長でない。さらに $u$ の冪による列の各相続商は
$\pi(B/uB)$ のシフトであるから、$\pi B\notin\cC_1$ である。
$\pi B$ の任意の非零部分対象は、非零斉次元 $f\in B$ に対する
$\pi(fB)\simeq\pi B(-\deg f)$ を含む。
$B/fB$ の Hilbert 多項式は次数が高々一なので、非零部分対象による商は
$\cC_1$ に入る。よって $\pi B$ は rank one torsion-free である。

最後に $c=(x+y)/2$、$d=(x-y)/2$ とし、$r=c,d$ に対し
\[
 M_r=\begin{pmatrix}r&u\\ -v&-r\end{pmatrix}
\]
と置く。$c^2=d^2=(x^2+y^2)/4$ により $M_r^2=gI$ であり、
これらは reduced な線形行列分解である。
成分の張る空間はそれぞれ $\langle u,v,c\rangle$ と
$\langle u,v,d\rangle$ である。これらは異なり、
行と列の定数基底変換で保たれるので、二つの余核は同型でない。
各分解では二つの連続行列が同じなので、$\Omega(1)$ は対応する余核を固定する。
この二つが全 MCM 型であることは $C(B)\simeq M_2(k)\times M_2(k)$ と
\cite[Lemma 4.13(4)--(5)]{MU2022} による。
実際、余核は二元生成で、表示行列は非零の最小関係を与えるため自由ではない。
非自由直和因子は少なくとも二元で生成されるので、各余核は直既約である。

$D=\operatorname{diag}(\mathrm i,1)$ とすると
$\sigma(M_r)=\mathrm i D M_r D^{-1}$ なので、右作用の $\sigma$ による引戻しは
両方の型を固定する。
Zhang 同値は極小自由被覆とその核を移送し、
$F(M)(1)\simeq F(M(1)_\sigma)$ であるから、
$B^\sigma$ でも $\Omega(1)$ は両方の型を固定する。
最小表示行列の成分は twist 後も同じ次元の空間を張る。
実際、$F(B(-1))\simeq B^\sigma(-1)$ を $m\mapsto\sigma^{-1}(m)$ で同定すると、
次数一成分が $a$ の左乗法は $a\star(-)$ となる。
従って $B^\sigma$ は non-standard であり、双方の成分空間が三次元である。
系~\ref{cor:normal-cubic-defect} により、どちらの初期型を選んでも
cubic 型は得られない。
\end{proof}

任意の別表示を含めた cubic 型の排除は、
\ref{thm:cubic-coverage} で証明する。

\begin{corollary}[この skew polynomial 環の中央切断の許容域]
\label{cor:skew-rank-three-central-locus}
上の $A$ の非零中央二次元 $g$ を
\[
 g=a x^2+b y^2+q(u,v)
\]
と書く。$A/(g)$ が domain かつ滑らかとなるための必要十分条件は
$ab\ne0$ かつ二次形式 $q$ が非退化となることである。
この領域では全ての中央切断およびその Zhang twist は
同じ pointed exact 同値類を与える。この類は $\HH$ に属し、
定理~\ref{thm:cubic-coverage} の cubic 表示を持たない類である。
\end{corollary}

\begin{proof}
次数二の中心は表示された五つの二次元で張られる。
これは skew polynomial の標準単項式を用い、
$x,y$ と交換する元の $x,y$ の次数がともに偶数となることから従う。
$ab\ne0$ かつ $q$ 非退化なら、$x,y$ のスカラー変更と
$u,v$ の線形基底変更で命題~\ref{prop:nonstandard-rank-three-example} の
$g$ に移る。この変更は $A$ の自己同型である。
Zhang 同値は構造対象を保存するので、この条件が十分である。

必要性を示す。$a=b=0$ なら非零の $q$ は一次元の積に分解し、
二因子は商で非零なので商は domain でない。
残りでは $x,y$ を交換して $a\ne0$ としてよい。
$B=A/(g)$ とし、$X,Y$ を $x,y$ の双対基底とする。
$B^!$ において $w=X^2$ は中央正則元の代表であり、
$XY=YX$、$Y^2=(b/a)w$ となる。
$u,v$ の双対に由来する二つの生成元は $X,Y$ と反交換し、
その平方と相互の反交換子は $w$ に $q$ の係数を掛けたものである。
これらはいずれも $B$ の二次関係空間との直交性を評価して得られる。
従って $C(B)=B^![w^{-1}]_0$ は
$Y X^{-1}$ と、$u,v$ の双対元に $X^{-1}$ を掛けた二元で生成される。

$b=0$ なら $Y X^{-1}$ は非零の中央平方零元なので $C(B)$ は
半単純でない。非零性は $w$ の正則性による。
また $q$ が退化しているなら、$u,v$ の双対側でその二次形式の根基から
非零元 $L$ を取れる。$L$ は平方が零であり、他の双対二次形式の
生成元との反交換子も零である。
よって $L X^{-1}$ は非零の平方零元で、
他の生成元と交換または反交換する。従ってそれが生成する
両側イデアルは非零で平方零となる。
いずれの場合も \cite[Theorem 5.5]{MU2022} により $\qgr B$ は滑らかでない。
\end{proof}

\begin{theorem}[四変数符号型の対角切断]
\label{thm:sign-quadric-diagonal}
$\epsilon_{ij}\in\{1,-1\}$ とし、
$A_\epsilon=k\langle x_1,\dots,x_4\rangle/
(x_jx_i-\epsilon_{ij}x_ix_j)_{i<j}$ とする。
非零中央元 $g=\sum_i a_i x_i^2$ による切断を考える。
$A_\epsilon/(g)$ が滑らかな domain となるのは、全ての $a_i$ が非零であり、
\[
 \epsilon_{ij}\epsilon_{ik}\epsilon_{jk}=-1
 \qquad(i<j<k)
\]
となる三つ組の個数が零または二の場合にちょうど限る。
前者は可換二次曲面の標識類、後者は
命題~\ref{prop:nonstandard-rank-three-example} の標識類を与える。
この二類はいずれも $\HH$ に属し、pointed exact 同値でない。
\end{theorem}

\begin{proof}
$a_j\ne0$ を一つ選ぶ。
二次双対において $w=X_j^2$ は中央正則で、
$X_i^2=(a_i/a_j)w$、$X_iX_l=-\epsilon_{il}X_lX_i$ ($i<l$) である。
$C(B)$ は $X_iX_j^{-1}$ ($i\ne j$) で生成され、
これらは相互に符号付きで交換する。
$a_i=0$ なら $X_iX_j^{-1}$ は零でない平方零元であり、
これが生成する両側イデアルも平方零である。
従って滑らかさに必要な半単純性が失われる
\cite[Theorem 5.5]{MU2022}。

以後 $a_i\ne0$ とする。対角符号自己同型 $x_i\mapsto s_i x_i$ による
Zhang twist は $\epsilon_{ij}$ を $s_i s_j\epsilon_{ij}$ に替える。
$g$ の新しい各平方に対する係数は非零のままなので、
変数のスカラー変更で正規化できる。
$g$ は自己同型で固定され、商の twist は構造対象を保つ。
また domain 性は斉次元の積に関してこの twist で不変である。

この符号変更で特定の頂点と残りの全ての間の符号を正にできる。
その後に残る三辺の符号は、その頂点を含む三つ組の積で決まる。
従って三つ組の四つの符号が各ラベルで一致するなら
二つの符号行列は twist で移り合う。
四符号の積は一であり、頂点の置換は三つ組の置換を全て実現する。
従って負の三つ組の個数は零、二、四のいずれかであり、
それぞれ全辺正、一辺だけ負、全辺負を代表として取れる。

零の場合の中央切断は滑らかな可換二次曲面である。
二の場合は可換変数二つと反交換する変数二つを持つ
前の系の非退化切断である。
四の場合には $\sum_i a_i x_i^2=(\sum_i \sqrt{a_i}x_i)^2$ となる。
一次の元は商で非零であり、平方が零なので domain ではない。
二つの残る標識類の非同値は、
可換類が cubic 表示を持ち、後者が持たないこと
（定理~\ref{thm:cubic-coverage}）から従う。
\end{proof}

% !TEX root = ../nc-quadric-surface-classification.tex
\subsection{符号型 skew polynomial 環の全正規切断}
\label{subsec:sign-normal-quadrics}

\begin{theorem}
\label{thm:sign-quadric-normal}
四変数符号型 skew polynomial 環 $A_\epsilon$ の非零正規二次元 $g$ を考える。
$B=A_\epsilon/(g)$ と置く。$(\qgr B,\pi B)$ が $\HH$ に属する場合の
pointed exact 同値類は、次の三つである。
\begin{enumerate}
\item 滑らかな可換二次曲面。
\item 命題~\ref{prop:nonstandard-rank-three-example} の非 cubic 型。
\item $g=x_1x_2+x_3x_4$、
$(\epsilon_{12},\epsilon_{34})=(-1,1)$、
$(\epsilon_{13},\epsilon_{14},\epsilon_{23},\epsilon_{24})=(1,-1,1,-1)$ の切断。
この類は cubic 型であり、反 Serre 環は非可換である。
\end{enumerate}
\end{theorem}

\begin{proof}
まず中央切断は定理~\ref{thm:sign-quadric-diagonal} に帰着する。
実際、$\epsilon_{ii}=1$ として符号行列を対称に書くと、
二次単項式 $x_i x_j$ ($i\ne j$) は第 $i,j$ 行が一致する場合に限り中央である。
各変数との交換は標準単項式ごとに検査できる。
同じ行を持つ変数は相互に可換であり、他の変数との交換符号も一様である。
従って同じ行の変数内での任意の線形基底変更は自己同型を与える。
中央二次形式はこれらの可換ブロック上の二次形式の和なので、
ブロックごとに対角化すればよい。

次に非中央の正規 $g$ を扱う。$x_i g=g\nu(x_i)$ と書く。
符号型の標準単項式の任意の辞書式順序で最高次項を比較すると、
積の最高指数は二因子の最高指数の和である。
従って $\nu(x_i)$ の最高単項式はどの順序でも $x_i$ であり、
$\nu(x_i)=\chi_i x_i$ となる。ここで $\chi_i=\pm1$ である。
$g$ の各単項式は同じ交換指標 $\chi=(\chi_i)_i$ を持つ。
$\chi\ne1$ なので平方は現れず、$x_i x_j$ が現れる条件は
第 $i,j$ 行の成分ごとの積が $\chi$ となることである。
同じ行のブロックとそれに $\chi$ を掛けたブロックは相異なる。
従って $g$ は互いに素なブロック対に属する双線形形式の和である。
それぞれを左右の基底変更で階数標準形にすると、四変数では
一つの積または二つの互いに素な積の和となる。
一つの積の商は domain でないので、
$g=x_1x_2+x_3x_4$ と取る。

$p=\epsilon_{12}$、$q=\epsilon_{34}$、$r=\epsilon_{13}$ とすると
両単項式の指標が一致するため
\[
 (\epsilon_{13},\epsilon_{14},\epsilon_{23},\epsilon_{24})
       =(r,pr,qr,pqr),\qquad
 \chi=(p,p,q,q).
\]
$p=q$ なら全ての三つ組の符号積は正である。
全ての交換符号を正にする対角符号 twist は、
両方の積を同じ符号 $p$ 倍するので $g$ の直線を保つ。
その可換商は滑らかな Segre 二次曲面を与える。
$p\ne q$ なら二組の交換と、必要なら第二組内の交換で
$(p,q,r)=(-1,1,1)$ を選べる。項の係数は変数のスカラー変更で戻せる。
これが (3) に掲げた切断である。

以下 (3) の $B=A_\epsilon/(g)$ を確認する。
$x_1$ は正規であり、$A_\epsilon/(x_1)$ は domain で $g$ の像も非零である。
$x_1 b=ga$ とすれば $a\in (x_1)$ となり、
$A_\epsilon$ 上で $x_1$ を消去して $b\in(g)$ を得る。
従って $x_1$ は $B$ 上で正則である。
$x_1$ を可逆にすると $x_2=-x_1^{-1}x_3x_4$ を代入でき、
残る三変数の skew Laurent 拡大に埋まるので $B$ は domain である。

二つの線形行列
\[
 \Phi=\begin{pmatrix}x_1&x_3\\x_4&x_2\end{pmatrix},\qquad
 \Psi=\begin{pmatrix}x_2&x_3\\x_4&-x_1\end{pmatrix}
\]
は $\Phi\Psi=gI$、$\Psi\nu(\Phi)=gI$ を満たす。
従ってどちらから始めても reduced な正規行列分解を得る。
可換多項式と見なした二つの行列式は
$x_1x_2-x_3x_4$ と $-x_1x_2-x_3x_4$ であり、比例しない。
従って定数の行・列変換で互いに移らず、
極小表示が異なる二つの二元生成 MCM 型を得る。
$g$ は一次元の積に分解しないので
\cite[Lemma 5.11]{MU2022} により、
対応する二次元の $C(B)$-加群は単純である。
ここで使う階数と次元の対応は同補題の証明および
同論文 Lemma 4.13 による。
$\dim C(B)=8$ なので半単純商の二つの四次元成分で全次元を使い、
$C(B)=M_2(k)\times M_2(k)$ となる。
従って $\qgr B$ は滑らかである。
両行列の成分空間は四次元である。

この符号行列では各行の符号積が正なので
$A_\epsilon$ の Nakayama 自己同型は恒等である
\cite[Example 5.5]{RRZ2014}。
$B$ の canonical twist は $\nu=\nu^{-1}$ である
\cite[Lemma 1.5]{RRZ2014}。
canonical bimodule と Serre 関手との対応には \cite[Example 3.6(2)]{MU2021} を用いる。
$\sigma$ を各変数に順に $\mathrm i,-\mathrm i,1,1$ を掛ける対角自己同型とすると、
$\sigma(g)=g$、$\sigma^2=\nu$ である。
従って $B^\sigma$ のシフトは旧圏では $(1)(-)_{\sigma}$ であり、
その平方が反 Serre となる。$A_\epsilon^\sigma$ は型 $(1,4,6,4,1)$ を保ち、
$B^\sigma$ はその正則正規二次商である。
飽和性、例外性、豊富性は命題~\ref{prop:nonstandard-rank-three-example} と同じ
正則切断および Serre 双対性の議論で従う。

Gabriel--Krull 条件については $B/(x_1)$ が
$B/(x_1,x_3)\simeq k_{-1}[x_2,x_4]$ に全射を持つ。
後者の正則元 $x_2$ の冪による下降列は、
隣接商が $k[x_4]$ のシフトなので tails で真に下降する。
従って $\pi(B/x_1B)$ は有限長でなく、
$x_1$ の冪による列により $\pi B\notin\cC_1$ となる。
domain 性と Hilbert 多項式から、命題~\ref{prop:nonstandard-rank-three-example}
と同様に rank one torsion-free 性を得る。
従って (3) は $\HH$ に属する。
twist 後も極小表示の成分空間の次元は保たれるので、
定理~\ref{thm:cubic-coverage} により (3) は cubic 型である。

最後に $B^\sigma$ の偶数 Veronese において、
$x_1x_3$ と $x_2x_3$ は反交換する。
旧積で既に反交換し、偶数の新しい積では
両方に同じ $\sigma^2$ の符号 $-1$ が働くからである。
domain 性から積は非零なので反 Serre 環は非可換である。
従って (3) は可換類とも非 cubic 類とも異なる。

逆に、どの切断の標識対も、$\HH$ に入るなら滑らかである。
また正則正規二次切断の AS Gorenstein 性により
元の切断環は飽和している。
元の次数シフトの $\pi B$ の軌道も rank one torsion-free なので、
切断の非零積は単射の合成となり、元の $B$ は domain である。
このため上で排除した非 domain または非滑らかな切断は、
別の平方根を選んでも新たな $\HH$ 類を与えない。
\end{proof}

% !TEX root = ../nc-quadric-classification.tex
\section{多パラメータ skew polynomial 環上の正規二次切断}
\label{sec:multiparameter-quadrics}

ここでは $q_{ij}\in k^\times$、$q_{ji}=q_{ij}^{-1}$、$q_{ii}=1$ とし、
$A=k\langle x_1,x_2,x_3,x_4\rangle/(x_jx_i-q_{ij}x_ix_j)_{i<j}$ を考える。
$Q_r$ ($r\in k^\times$) を、次のパラメータの skew polynomial 環の
$g=x_1x_2+x_3x_4$ による商とする。
\[
 q_{12}=q_{34}=1,\qquad (q_{13},q_{14},q_{23},q_{24})=(r,r^{-1},r^{-1},r).
\]

\begin{theorem}[多パラメータ切断の分類]
\label{thm:multiparameter-quadric-classification}
上の $A$ の正規二次切断から得られる $\HH$ の pointed exact 同値類は、
命題~\ref{prop:nonstandard-rank-three-example} の非 cubic 型一つと、
$Q_r$ の族である。$Q_r$ はすべて cubic 型であり、
\[
 (\qgr Q_r,\pi Q_r)\simeq(\qgr Q_{r'},\pi Q_{r'})
 \quad\Longleftrightarrow\quad
 r'^2=r^2\ \text{または}\ r'^2=r^{-2}
\]
が成り立つ。可換類は $r^2=1$ に対応する。
\end{theorem}

\begin{proof}
正規切断に対してその元の切断環が domain かつ滑らかであることは、
定理~\ref{thm:sign-quadric-normal} の末尾と同じ理由で必要である。
正規化自己同型が対角となる最高単項式の議論も、交換係数が非零ならそのまま使える。
各変数の交換指標を $\rho_i=(q_{ij})_j$ と書く。
$g$ の各単項式は同じ指標 $\chi$ を持ち、
$x_i x_j$ の指標は $\rho_i\rho_j$ である。
等しい指標を持つ変数のブロック内では変数は可換であり、
ブロックごとの基底変更は $A$ の自己同型である。
$\rho\mapsto\chi/\rho$ が固定するブロックには二次形式が現れ、
別のブロックへ移す場合にはブロック対に双線形形式が現れる。
それぞれ対角化または階数標準形にすることで、$g$ は
重複のない変数の平方と、互いに素な変数対の積の和になる。

この形で $g$ に現れない変数があれば商は滑らかでない。
実際、その変数の双対 $X$ は $(A/(g))^!$ で平方零であり、
他の双対生成元と非零スカラーを介して交換する。
このため $(X)$ は非零平方零イデアルである。
$w$ の代表は、$g$ に現れる一つの平方の双対平方、
または一つの積に対応する双対の二項和として取れる。
いずれも同じ変数の積からなるため、$Xw$ と $wX$ はスカラー倍で一致する。
\cite[Lemma 4.11]{MU2022} の正則正規二次元 $w$ を逆にしても消えず、
その局所化の次数零部分にも非零の平方零イデアルを与える。
実際 $Xw\ne0$ なので、$w$ を二文字の和に展開すると
ある双対生成元 $Y$ に対して $XY\ne0$ となり、$XYw^{-1}$ がその部分に入る。
従って $C(B)$ は半単純でなく、\cite[Theorem 5.5]{MU2022} により
商は滑らかでない。

全変数が現れる場合、可能な形は四平方、二平方と一対の積、二対の積である。
四平方なら各成分で $\chi_j=q_{jj}^2=1$ なので全交換係数が符号となり、
定理~\ref{thm:sign-quadric-normal} が適用できる。
二平方の場合に $g=x_1^2+x_2^2+x_3x_4$ とする。
指標の等式から $a=q_{13}$ とすれば
$q_{14}=a^{-1}$、$q_{23}=\pm a$、$q_{24}=q_{23}^{-1}$、
$q_{34}=a^{-2}$、$q_{12}=\pm1$ である。
対角自己同型 $(1,1,a,a^{-1})$ による Zhang twist は
$g$ の直線を保存し、交換係数をすべて符号にする。
従ってこの場合も符号型の分類に帰着する。

二対の積の場合は $g=x_1x_2+x_3x_4$ と取る。
$p=q_{12}$、$q=q_{34}$、$a=q_{13}$ とすると
\[
 q_{14}=p/a,\qquad q_{23}=1/(qa),\qquad q_{24}=qa/p
\]
である。$b^2p=c^2q=1$ を満たす $b,c$ を選び、
対角自己同型 $(b,b^{-1},c,c^{-1})$ で twist すれば、
新しい $p,q$ はともに一になり、$a$ は $r=ba/c$ となる。
二つの積の係数は非零なので変数のスカラー変更で $g$ を正規化できる。
これにより $Q_r$ が得られる。
符号型の混合符号の第三類もこの操作により $r^2=-1$ の $Q_r$ へ移る。

次に全ての $Q_r$ が $\HH$ に属することを示す。
$x_1$ は正則正規であり、それを逆にすると
$x_2=-x_1^{-1}x_3x_4$ を消去できるので、商は domain である。
正則性と消去は定理~\ref{thm:sign-quadric-normal} の議論に従い、
交換係数に上の $r$ を入れて確かめられる。
$r^2=1$ の場合は符号 twist で可換例になる。
$r^2\ne1$ の場合、
\[
 \Phi=\begin{pmatrix}x_1&x_3\\x_4&-r x_2\end{pmatrix},\qquad
 \Psi=\begin{pmatrix}x_2&x_3\\x_4&-r^{-1}x_1\end{pmatrix}
\]
は積が両方向とも $gI$ となる線形行列分解を与える。
可換多項式の行列式は比例しないため、二つの MCM 型は相異なる。
$C(Q_r)$ が $M_2(k)\times M_2(k)$ となることと滑らかさは、
定理~\ref{thm:sign-quadric-normal} と同じ次元の議論で従う。
両方の行列の成分は四次元を張る。
各変数の交換係数の積が一で $g$ は中央なので、
Nakayama 自己同型は $Q_r$ でも恒等である
\cite[Example 5.5 and Lemma 1.5]{RRZ2014}。
従って通常シフトが反 Serre の平方根となる。
飽和性、例外性および豊富性は正則切断の議論で確認できる。
$x_1$ で割った後さらに $x_3$ で割ると二変数の skew polynomial 環が得られ、
その tails の構造対象に無限下降列がある。
従って同じ $x_1$ の冪の議論で Gabriel--Krull 条件も従う。
以上で $\HH$ に入り、定理~\ref{thm:cubic-coverage} から cubic 型となる。

最後に同値を判定する。$r\mapsto-r$ は $g$ の直線を保つ
対角符号 twist で実現でき、$r\mapsto r^{-1}$ は $x_1,x_2$ の交換で実現できる。
必要性には命題~\ref{prop:first-orbit-finite-equivalence} の
添字付き Hom 環の同型を使う。
関係の双線形化が定める点対 $(p,p')\in\PP^3\times\PP^3$ では、
交換関係から $p,p'$ の非零座標の集合が一致する。
さらに三座標が非零なら対応する三角形の交換比の積が一でなければならない。
この積は $r^2$ または $r^{-2}$ なので、$r^2\ne1$ なら非零座標は高々二つである。
$g$ の双線形化は座標対 $12,34$ を排除するため、
各射影の像は座標対 $13,14,23,24$ の四直線からなる正方形である。
$r^2=1$ の可換類では像は滑らかな二次曲面なので、
二つの領域は添字付き線形同型で交わらない。

両方が $r^2\ne1$ の領域にあれば、その同型の各添字の射影線形変換は
四直線の交点、すなわち四座標点を置換する。
点対が各座標点を同じ座標点へ結ぶので、この置換 $\pi$ は添字に依存しない。
従って第 $i$ 添字で $x_a$ を $c_{i,a} x_{\pi(a)}$ へ送ると書ける。
隣接二添字で交換関係を送ったとき、任意の異なる $a,b$ について
$q'_{\pi(a)\pi(b)}/q_{ab}$ は
$c_{i+1,b}/c_{i,b}$ を $c_{i+1,a}/c_{i,a}$ で割った値である。
これは反対辺の対 $12,34$ についても成り立つ。
その二つの単項式は相互に比例するが零ではなく、
各交換関係の像は一方の対の単項式だけを含むからである。
よって三角形を回ると定数倍の比は相殺し、
交換比の三角積が保存される。三角積は $r^2$ と $r^{-2}$ なので
$r'^2=r^2$ または $r'^2=r^{-2}$ が必要となる。
\end{proof}

% !TEX root = ../nc-quadric-surface-classification.tex
\section{交差項を持つ Ore 環の正規二次切断}
\label{sec:ore-normal-quadrics}

$u,v$ を中央の次数一元、$q\in k^\times$、$0\ne h\in k[u,v]_2$ とし、
\[
 A=k[u,v]\langle x,y\rangle/(yx-qxy-h)
\]
を考える。これは $k[u,v,x]$ 上の Ore 拡大であり、型 $(1,4,6,4,1)$ を持つ
（命題~\ref{prop:ore-ambient-family} と同じ Ore 拡大定理を用いる）。

\begin{proposition}
\label{prop:ore-normal-quadrics}
$A$ の非零正規二次元は、スカラー倍を除いて次の形に限る。
\begin{enumerate}
\item $q=1$ の場合は $k[u,v]_2$ の元。
\item $q\ne 1,-1$ の場合は $k[u,v]_2$ の元または $xy+h/(q-1)$。
\item $q=-1$ の場合は $a x^2+b y^2+H(u,v)$ または $xy-h/2$。
\end{enumerate}
従って $q\ne-1$ のとき、これらの切断が与える $\HH$ の標識類は、
$q\ne1$ かつ $h$ が非退化である場合の可換二次曲面一つに限る。
$q=-1$ の非中央切断も同じ条件で同じ標識類を与える。
\end{proposition}

\begin{proof}
$g$ を正規二次元とし $ag=g\nu(a)$ と書く。中央元 $u,v$ は $\nu$ で固定される。
$x,y$ の全次数を先に比較する次数フィルトレーションを取ると、
その付随次数環は $yx=qxy$ の skew polynomial 環である。
その中で $x,y$ の辞書順序をどちらから始めてもよく、
積の最高単項式では指数が加算される。
$xg=g\nu(x)$ と $yg=g\nu(y)$ の最高指数を比較すると、
$\nu(x)=\lambda x+L(u,v)$、$\nu(y)=\mu y+M(u,v)$
となり、$\lambda,\mu\ne0$ である。

$q\ne1$ なら $\nu$ を Ore 関係に適用する。
$x,y$ の一次の係数から $(1-q)\lambda M=(1-q)\mu L=0$ となり、
$L=M=0$ である。また $h\ne0$ なので $\lambda\mu=1$ である。
$g$ の $x,y$ 次数が最大の層にある単項式 $x^m y^n$ を取ると、
$x,y$ との交換から $\lambda=q^{-n}$、$\mu=q^m$ となる。
従って $q^{m-n}=1$ である。
$q\ne1,-1$ なら $m+n\le2$ と合わせて $m=n$ となる。
最大層が $xy$ の場合、$x,y$ 次数二の比較で
$u,v$ を係数とする一次の項が消える。
実際、$x$ との交換で $x$ 項は消え、$y$ との交換で $y$ 項は消える。
$g=a xy+H(u,v)$ と書くと
$xg=gq^{-1}x$ は $(q-1)H=ah$ と同値である。
最大層が次数零なら $g\in k[u,v]_2$ である。

$q=-1$ の場合、最大層が $xy$ なら同じ計算を適用する。
他の最大層では $m-n$ が偶数で、平方の層または次数零しかなく、
$\lambda=\mu=1$ となる。$x^2,y^2$ は中央元であり、
残る $x,y$ の一次の項は中心に入らないので、
中央二次式は表示した形になる。

$q=1$ なら交換子は $x,y$ 次数を下げる。
最高層の比較で $\lambda=\mu=1$ となり、
$[x,g]=gL$、$[y,g]=gM$ の次数を比較すれば $L=M=0$ である。
従って $g$ は中央である。
PBW 表示で $[y,-]$ は $h\,\partial/\partial x$、
$[x,-]$ は $-h\,\partial/\partial y$ として作用する。
標数零かつ $h\ne0$ なので中心は $k[u,v]$ である。

$k[u,v]_2$ の非零元は一次式の積に分かれ、商は domain でない。
一方 $g=xy+h/(q-1)$ で割ると $xy=yx=h/(1-q)$ となり、
商は四変数可換多項式環の二次切断である。
これは $h$ が非退化である場合にちょうど滑らかで、
その場合は可換二次曲面を与える。
元の切断環の domain 性と tails の滑らかさは、
定理~\ref{thm:sign-quadric-normal} の末尾と同じく $\HH$ への所属に必要である。
従って上の標識類以外は、この範囲では現れない。
\end{proof}

\subsection{反交換 Ore 環の中央切断}
\label{subsec:ore-central-anticommuting}

\begin{theorem}
\label{thm:ore-anticommuting-central}
$q=-1$ の中央切断 $g=a x^2+b y^2+H(u,v)$ が
$\HH$ の標識類を与えるための必要十分条件は、
$a,b$ がともに非零であり、以下の二つの二次形式 $f_1,f_2$ が非退化となることである。
$x,y$ をスカラー変更して $a=b=1$ とし、その規約で $f_1=-H+h$、$f_2=-H-h$ とする。
許容される標識類は、$\PP^1$ 上の相異なる二点の非順序対を二つ取り、
対の交換と $\PGL_2$ で割ったデータで判定される
（二つの対は互いに一致してもよい）。どの類も標識付き cubic 表示を持たない。
\end{theorem}

\begin{proof}
$a=b=0$ なら商は domain でない。
一方だけ零、例えば $b=0$ なら $B^!$ で $w=X^2$ は中央正則で、
$Y^2=0$、$XY=YX$ である。$X,Y$ は $u,v$ の双対生成元と反交換する。
従って $C(B)$ 中で $YX^{-1}$ は非零中央平方零元であり、
この場合は滑らかでない。

次に $a=b=1$ とする。$c=x+y,d=x-y$ と置けば
\[
 B=k[u,v]\langle c,d\rangle/(c^2-f_1,d^2-f_2).
\]
交互に並ぶ $c,d$ の語が $k[u,v]$ 上の基底である。
書換え $c^2\mapsto f_1$、$d^2\mapsto f_2$ は長さを下げ、
三文字の同じ文字の重なりは係数の中心性で解消するからである。
元の Ore 環の関係は $c^2-d^2=f_1-f_2$ である。
双対で $CD=DC=0$、$w=C^2+D^2$ は正則正規二次元の代表である。
$C^2,D^2$ は中央なので、$C(B)$ は直交冪等元 $C^2/w,D^2/w$ により分解する。
第一成分では $C$ が対応する局所化の角で可逆となり、
$UC^{-1},VC^{-1}$ が二次形式が $f_1$ に比例する二元 Clifford 環を生成する。
第二成分では $D$ と $f_2$ を用いた同じ計算になる。
実際、$C,D$ は $U,V$ と反交換し、
$U,V$ の平方と反交換子は $f_1,f_2$ の係数を掛けた $C^2,D^2$ の和である。
従って各成分は高々四次元で、$\dim C(B)=8$ により両方が四次元である。
標数零なので二元 Clifford 環は非退化の場合にちょうど半単純であり、
退化の場合には根基のベクトルが非零平方零イデアルを生成する。
従って滑らかさは $f_1,f_2$ がともに非退化であることと同値である
\cite[Lemma 4.13 and Theorem 5.5]{MU2022}。

以後 $f_1,f_2$ を非退化とする。$B$ は domain である。
$F=k(u,v)$ に移し、中央元 $z=cd+dc$ を取ると、
交互語の基底から $F[z]$ 上で $1,c,d,cd$ が自由基底になる
（$z$ の冪の最高長さの交互語を比較する）。
さらに $F(z)$ に移し $e=d-zc/(2f_1)$ とすれば、
$c^2=f_1$、$ec=-ce$、$e^2=-(z^2-4f_1 f_2)/(4f_1)$ の四元数環となる。
$f_1,f_2$ はともに $F$ の平方でない。
$z^2-4f_1 f_2$ の任意の既約因子の素点では、$f_1$ は剰余体でも平方ではない。
$f_1 f_2$ が平方なら剰余体は $F$ であり、そうでなければ $F(\sqrt{f_1 f_2})$ である。
後者で $f_1$ が平方なら $F$ で $f_1$ または $f_2$ が平方になってしまう。
従ってこの素点における $\sqrt{f_1}$ の拡大は非分岐で剰余次数二となる。
ノルムの付値は偶数だが $e^2$ の付値は一である。
よって $e^2$ はノルムでなく、この四元数環は斜体である。
自由性により元の $B$ も domain となる。

適合根と被覆を構成する。
$f_1,f_2$ の二つの根の対を互いに交換する $\PP^1$ の対合を取れる。
根の対が一致するときは恒等を使う。
一根を共有するとき、座標で二対を $\{0,\infty\},\{0,1\}$ とすれば
$t\mapsto t/(t-1)$ が働く。四根が相異なるときは
$\{0,\infty\},\{1,\lambda\}$ とし、$t\mapsto\lambda(t-1)/(t-\lambda)$ を使う。
二乗が恒等の線形持上げ $L$ を選び、
$Lf_1=\eta f_2,Lf_2=\eta^{-1}f_1$ とする。
$\alpha^2=\eta$ を選び、$\sigma$ を $u,v$ 上で $L$、
$c\mapsto\alpha d$、$d\mapsto-\alpha^{-1}c$ とする。
$r_c=c^2-f_1,r_d=d^2-f_2$ と書けば $\sigma$ は
$r_c-\eta r_d$ の直線と $r_c+\eta r_d$ の直線を保つ。
$A'$ を中央 $u,v$ と第一の関係 $r_c-\eta r_d$ だけで定め、
第二の元 $g'=r_c+\eta r_d$ で切断すれば $B$ が得られる。
$c\pm\sqrt\eta\,d$ を新しい変数に取ると、
$A'$ は反交換 Ore 環である。型 $(1,4,6,4,1)$ を持ち、
$c^2,d^2$ が中央であるため $g'$ は中央正則である。

$A'$ の Nakayama 自己同型は $u,v$ を固定し $c,d$ に負号を掛ける。
Ore 表示の基礎環での作用は \cite[Theorem 3.3]{LWW2014} により、
残る変数の作用は Ore 関係との整合から従う。
Ore 関係の $k[u,v]$ 項が非零なら、基礎環の非中央変数に関する次数を比較して
Ore 関係を保つために平行移動項は零となり、係数は $-1$ である。
Ore 関係の $k[u,v]$ 項が零なら符号型の Nakayama 公式を使える。
従って中央切断の公式 \cite[Lemma 1.5]{RRZ2014} により、$B$ の Nakayama 自己同型も
$u,v$ を固定し $c,d$ を反転する。これは $\sigma^2$ と一致する。
よって $B^\sigma=A'^\sigma/(g')$ の通常シフトが適合根となる。
飽和性、例外性、豊富性は前の符号型の場合と同じ正則切断の議論による。

GK 条件には、$B$ が $k[u,v]$ 上自由で、
$B/(u,v)=k\langle c,d\rangle/(c^2,d^2)$ が無限次元であることを使う。
$B/uB$ 上の $v$ の冪の列は tails で真に下降するので、
$B/uB$ の像は有限長でない。
$u$ の冪の列と domain 性を使うと、前と同様に rank one torsion-free 性が従う。

$f_1=-\ell_1\ell_2$ と一次式の積に書けば、
行列 $\left(\begin{smallmatrix}c&\ell_1\\ \ell_2&-c\end{smallmatrix}\right)$ の二乗は
$r_c I$ であり、$A'$ 上で $g'$ の非零スカラー倍になる。
$f_2$ と $d$ に対しても同様である。
二型の成分空間はそれぞれ $\langle u,v,c\rangle$ と
$\langle u,v,d\rangle$ で、両方が階数三である。
twist は階数を保つので標識付き cubic 表示は存在しない。

最後に二つの許容切断を比較する。
適合環の添字付き代数は、それぞれ twist 前の $B$ の添字付き代数と
同型である。従って命題~\ref{prop:first-orbit-finite-equivalence} を使える。
上の $B$ の双線形化関係は、$(u,v)\ne(0,0)$ では
\[
 E=\{c^2=f_1(u,v),\ d^2=f_2(u,v)\}\subset\PP^3
\]
の上の対角対応を与える。$u,v$ のいずれかを使う交換関係から
次の点が比例するためである。$(u,v)=(0,0)$ では、
$c,d$ の二つの座標点が互いに移る。これらは $E$ には含まれない。
二つの多項式の根の対が異なれば $E$ は整で $\PP^3$ を張る。
実際、$\PP^1$ の関数体上で独立な二平方根を付加する被覆であり、
式は基底 $1,c,d,cd$ の有限自由被覆を与える。
根の対が一致すれば $E$ は異なる二平面にある二つの滑らかな円錐曲線であり、
二点で交わる。いずれの場合も $E$ を点ごとに固定する射影線形変換は恒等である。
整の場合は固有空間のいずれか一つに全体が入ることから、
二成分の場合は二平面上のスカラーが交点で一致することから従う。

従って添字付き同型の連続する射影線形変換は $E$ 上で一致するので、
スカラー倍まで同じである。よって同値は $B$ の次数付き同型に帰着する。
次数一の中心は $\langle u,v\rangle$ である。
$c$ の像を $a c'+b d'+\ell(u,v)$ と書くと、平方が中央変数の多項式であるため
交互語の基底から $ab=0,\ell=0$ を得る。$d$ に対しても同様である。
従って次数付き同型は $u,v$ の基底変更と、$c,d$ の交換およびスカラー変更
に尽きる。これはちょうど二つの根の対に対する主張した同値関係である。
逆に任意の二対を定める非退化な二次形式を取り、
必要なら一方をスカラー倍して $f_1-f_2\ne0$ とすれば、この Ore 族から実現できる。
\end{proof}

\begin{corollary}
多パラメータ skew polynomial の切断と、この節の Ore 環の切断を合わせると、
$\HH$ に属する標識類は、$r^2$ を逆元で割った cubic 族と、
定理~\ref{thm:ore-anticommuting-central} の非 cubic 族である。
後者は根の対が一致する一点、根を一つ共有する一点、
および四根が異なる場合の $\lambda\in k\setminus\{0,1\}$ を
$\lambda\leftrightarrow\lambda^{-1}$ で割った族からなる。
\end{corollary}
\begin{proof}
前者は定理~\ref{thm:multiparameter-quadric-classification} による。
同定の際、多パラメータ切断の非 cubic の一点は
命題~\ref{prop:nonstandard-rank-three-example} を $c,d$ の表示で書けば
両平方が同じ二次形式となる場合である。
後者の四根が異なる場合は二対を $\{0,\infty\},\{1,\lambda\}$ へ移せる。
各対内の交換は交比を逆にし、二対そのものの交換は交比を保存する。
一致と一根共有の場合は、それぞれ二点または三点に対する
$\PGL_2$ の推移性から一点となる。
\end{proof}

% !TEX root = ../nc-quadric-surface.tex
\subsection{反 Serre Hilbert 多項式による最初の軌道対象の認識}
\label{sec:first-orbit-object}

\begin{proposition}
\label{prop:first-orbit-object}
$(\cC,\OO)\in\HH$ とし、$t=S^{-1}[2]$、適合根を $s$ とする。
Gabriel--Krull rank one torsion-free な exceptional 対象 $M$ が
\[
 \chi(\OO,t^n M)=(2n+2)^2\qquad(n\gg0)
\]
を満たせば、$M\simeq s\OO$ である。
従って、どの適合根 $s'$ を用いても $s'\OO\simeq s\OO$ である。
\end{proposition}

\begin{proof}
定理~\ref{thm:normal-helix} の四連続項
$(E_0,E_1,E_2,E_3)=(\OO,\pi X_0,s\OO,\pi X_1(1))$
の類は $K_0(\cC)$ の基底である。Euler 形式の行列は
\[
 G=\begin{pmatrix}
 1&2&4&6\\
 0&1&2&4\\
 0&0&1&2\\
 0&0&0&1
 \end{pmatrix}.
\]
実際、$\Hom(\OO,\pi X_i)$ の次元は二元生成から $2$ であり、
$\Hom(\pi X_0,s\OO)$ の次元も双対評価列から $2$ である。
同じ評価列を次数一で取ると $\dim (X_1)_1=2\dim B_1-\dim(p(X_1))_0=6$ であり、
$\dim B_1=4$ である。
残る $\Hom(\pi X_0,\pi X_1(1))$ は
系~\ref{cor:normal-cubic-defect} の奇数項間の積の同型から四次元である。
他の成分は full geometric helix の Ext 消滅と例外性から従う。

$G$ は $Q=\PP^1\times\PP^1$ 上の
$\OO_Q,\OO_Q(1,0),\OO_Q(1,1),\OO_Q(2,1)$ の Euler 行列と一致する。
これらは $K_0(Q)$ の整数基底である。例えば積の標準基底
$\OO_Q,\OO_Q(1,0),\OO_Q(0,1),\OO_Q(1,1)$ から、
$[\OO_Q(2,1)]=2[\OO_Q(1,1)]-[\OO_Q(0,1)]$ によって移る。
従って上記の順序で $K_0(\cC)$ を $K_0(Q)$ と Euler 形式を保って同定できる。
Serre 関手の $K_0$ 上の作用は
$\chi(F,G)=\chi(G,SF)$ と Euler 形式の非退化性で一意に定まる。
従ってこの同定で $t$ の作用は $\OO_Q(2,2)$ をテンソルする作用である。

$H_1,H_2$ を $Q$ の二つの因子から来る超平面類とし、
$[M]$ に対応する類の Chern 指標を
$r+aH_1+bH_2+cH_1H_2$ と書く。
$Q$ 上の Riemann--Roch により
\[
 \chi(\OO,t^n M)=r(2n+1)^2+(a+b)(2n+1)+c,\qquad
 \chi(M,M)=r^2+2rc-2ab.
\]
第一式の係数比較から $r=1$、$a+b=2$、$c=1$ を得る。
例外性から第二式は $1$ なので $ab=1$、従って $a=b=1$ である。
Chern 指標の有理同型と $K_0(Q)$ の無捩れ性から $[M]=[s\OO]$ が従う。

$L=s\OO$ と置く。$\chi(L,M)=1$ より $\Hom(L,M)$ と
$\Ext^2(L,M)$ の少なくとも一方は非零である。
後者が非零なら Serre 双対性により $M\to t^{-1}L$ に非零射がある。
両対象は rank one torsion-free なので、この射は単射である。
実際、\ref{lem:sections-monic} の証明と同じく、商圏で非零となるため
商圏で同型となり、核は torsion-free 性によって零である。
従って $\Hom(\OO,t^n M)\hookrightarrow\Hom(\OO,t^{n-1}L)$ を得る。
ここで $\dim\Hom(\OO,t^n M)=\chi(\OO,t^n M)$ は $n\gg0$ で成り立つ。
定理~\ref{thm:normal-helix} の証明末尾の自由被覆の議論を偶数項に適用すれば
$\Ext^1(t^{-n}\OO,M)=0$ を得る。
また $\Ext^2(t^{-n}\OO,M)=\Hom(M,t^{-n-1}\OO)^*=0$ は、
$M$ の自由被覆と負次数切断の消滅から従う。
よって上の単射は $(2n+2)^2\le (2n)^2$ を与え、矛盾する。
従って $\Hom(L,M)$ に非零射があり、同じ理由でこれは単射である。
余核の $K_0$ 類は零である。
$B=A/(g)$ 上の有限生成加群の Hilbert 多項式は、
$A$ 上の有限自由分解により存在し、tails 上で加法的である。
非零の tails の対象は零 Hilbert 多項式を持たないので、余核は零となる。

適合根 $s'$ に対しては
$\dim\Hom(\OO,t^n s'\OO)=(2n+2)^2$ が切断環の Hilbert 関数から従い、
前段の Ext 消滅により漸近的に Euler 特性と一致する。
$s'\OO$ は $\OO$ と同じ rank one torsion-free 性と例外性を持つため、
第一の主張を適用できる。
\end{proof}

\begin{corollary}[異なる適合根における階数の保存]
\label{cor:first-orbit-rank-invariance}
同じ標識対象の二つの適合切断環について、二つの次数零生成 MCM 型の
最小線形表示行列の成分空間の次元は、二つの値の多重集合として一致する。
特に命題~\ref{prop:nonstandard-rank-three-example} の例では、
どの適合根から定理~\ref{thm:normal-helix} の方法を用いても
二つの階数はともに三である。
\end{corollary}
\begin{proof}
$s^2\simeq (s')^2\simeq t$ と命題~\ref{prop:first-orbit-object} から
$s^j\OO\simeq(s')^j\OO$ が全整数 $j$ について従う。
これらの対象の添字を保つ同型を選ぶと、合成を積とする
$\Hom(s^{-j}\OO,s^{-i}\OO)$ の $\mathbb Z$-algebra が同型になる。
各切断環 $B$ とその $\mathbb Z$-algebra の間では、
添字 $(i,j)$ 成分を $B_{j-i}$ と同定する
$a\mapsto s^{-j}(a)$ が合成を保つ。
従って得られる次数付き右加群圏の同値は
全添字の自由加群 $B(i)$ を対応する自由加群に移し、
最小自由被覆、Ext、次数零生成性を保存する。
特に自由加群の全シフトを係数に取った Ext 消滅で定まる MCM 型を移送する。
最小線形表示に現れる $B(-1)\to B$ の Hom 空間が線形同型で移るので、
その中で行列成分が張る空間の次元は保たれる。
\end{proof}

\begin{proposition}[四生成元による有限同値判定]
\label{prop:first-orbit-finite-equivalence}
二つの $\HH$ 対象の適合切断環を、
次数一の基底を選んでそれぞれ $B=T(k^4)/(K)$、$B'=T(k^4)/(K')$
とする。$G=\GL_4(k)$ とし、
\[
 E=\{(g,h)\in G\times G\mid (g\otimes h)K=K'\}
\]
を作る。$C_0=G$ から始め、射影 $\operatorname{pr}_i$ を用いて
\[
 C_{n+1}=\operatorname{pr}_1\bigl(E\cap(G\times C_n)\bigr)
         \cap\operatorname{pr}_2\bigl(E\cap(C_n\times G)\bigr)
\]
と定める。この下降列は有限段で安定し、安定集合が非空であることと
二つの対象の pointed exact 同値性は同値である。
各段は $\GL_4$ の座標環での有限の代数的消去により定まる。
機械的な計算可能性を述べる場合には、係数体が有効に与えられていることを前提とする。
\end{proposition}

\begin{proof}
命題~\ref{prop:first-orbit-object} から、pointed exact 同値は
各添字の軌道対象を保ち、系~\ref{cor:first-orbit-rank-invariance} の証明と同じ
indexed Hom の構成により $\check B\simeq\check B'$ を与える。
逆にこのような添字を保つ $\mathbb Z$-algebra 同型は、
各次数成分を各添字の加群成分と同定することにより、
次数付き右加群の圏の間の同値を誘導する。全ての自由シフトと
次数ごとのベクトル空間を保つため、有限生成性と有限次元性を保つ。
$B,B'$ は右 noetherian なので有限生成次数付き加群の圏は Abel 圏であり、
その torsion 対象は有限次元である。この部分圏を割ると exact 同値を得て、
自由加群の添字零は $\OO$ に対応する。

$B,B'$ は二次環なので、$\check B$ と $\check B'$ はそれぞれ
各隣接成分を $k^4$ とするテンソル経路環から、各二段の関係 $K$、$K'$ を課して得られる。
従って添字を保つ同型は、各隣接成分上の $g_i\in G$ により、
$(g_i,g_{i+1})\in E$ ($i\in\mathbb Z$) として必要十分に記述される。

$G\times G$ は $V^{\otimes2}$ の部分空間の Grassmannian に作用する。
$E$ は $K$ の $K'$ への transporter なので、空集合か、
閉安定化部分群 $H_E\subset G\times G$ の剰余類である。
ここで有限停止には、一般の構成可能集合の下降列ではなく
\emph{剰余類の構造が再帰で保たれる}ことが重要である。
実際、$C_n$ が非空な閉部分群 $H_n\subset G$ の剰余類なら、
$E\cap(G\times C_n)$ と $E\cap(C_n\times G)$ は
それぞれ非空な場合には
$H_E\cap(G\times H_n)$ と $H_E\cap(H_n\times G)$ の剰余類である。
それらの $\operatorname{pr}_1,\operatorname{pr}_2$ による像は、
代数群準同型による閉部分群の像の剰余類だから閉集合である。
二つの剰余類の共通部分も、非空なら閉部分群の剰余類である。
$C_0=G$ からの帰納法により、全ての $C_n$ は空集合または
閉代数部分群の剰余類となる。
従って $C_n$ は閉集合の下降列であり、$G$ の座標環の noetherian 性から
有限段で安定する。ここでは任意の閉集合の射影が閉とは主張しない。

また $g\in C_n$ となることは、$g$ を中央とする長さ $2n$ の有限列を
作れることと同値である。両側への延長可能性が再帰式を与える。
一度 $C_{n+1}=C_n$ となれば再帰式により以後は同じである。
安定集合が非空なら、その中の任意の元から左右に安定集合内の隣を選び続けられるので、
両側無限列を構成できる。逆向きには両側無限列の中央項がすべての $C_n$ に入る。

上記の剰余類に限れば射影像は閉である。
従って行列式の逆元も変数として含めた座標環で消去を行い、
各像と共通部分の根基イデアルを求めればよい。
根基の一致および全環になるかどうかを比較すると停止と非空性を判定できる。
\end{proof}

\begin{corollary}[二つの行列による pointed 同値判定]
\label{cor:two-matrix-pointed-equivalence}
前命題の仮定の下で、pointed exact 同値が存在することは
\[
 \exists\,g,h\in\GL_4(k),\qquad
 (g\otimes h)K=K',\qquad(h\otimes g)K=K'
\]
と同値である。
\end{corollary}
\begin{proof}
pointed exact 同値 $F$ を取り、両側の適合根を $s,s'$、Serre 関手を $S,S'$ とする。
$D^b(F)$ は Serre 双対を移すので、表現対象の一意性から自然同型
$D^b(F)S\simeq S'D^b(F)$ を得る。$s^{-2}\simeq S[-2]$ と
$(s')^{-2}\simeq S'[-2]$ を使い、heart に制限して
$\alpha:F s^{-2}\xrightarrow{\sim}(s')^{-2}F$ を一つ固定する。
$P_i=s^{-i}\OO,\ P'_i=(s')^{-i}\OO'$ と置く。pointing と
命題~\ref{prop:first-orbit-object} から $u_0:F P_0\simeq P'_0$ と
$u_{-1}:F P_{-1}\simeq P'_{-1}$ を選べる。
各偶奇の列において
\[
 u_{i+2}=(s')^{-2}(u_i)\circ\alpha_{P_i}
\]
で正負両方向へ延長する。逆向きにも $(s')^{-2}$ は全忠実なので一意に延長できる。

$a\in B_1$ に対し $a_i=s^{-(i+1)}(a):P_{i+1}\to P_i$ と書き、
$a'_i$ も $B'_1$ に対して同様に定める。一次成分の同型 $g_i$ は
$(g_i a)'_i=u_iF(a_i)u_{i+1}^{-1}$ で与えられる。$\alpha$ の自然性から
\[
 u_{i+2}F(a_{i+2})u_{i+3}^{-1}
   =(s')^{-2}\bigl(u_iF(a_i)u_{i+1}^{-1}\bigr)
   =(g_i a)'_{i+2}.
\]
従って $g_{i+2}=g_i$ である。右加群規約では
$a_i b_{i+1}=s^{-(i+2)}(s(a)\circ b)$ なので、
二段の関係を保つ条件は順に $(g_0\otimes g_1)K=K'$ と
$(g_1\otimes g_0)K=K'$ となる。
逆に解 $(g,h)$ から偶数辺に $g$、奇数辺に $h$ を置けば全ての二次関係を保つ。
逆行列も逆向きの関係を保つので、添字付き環の同型となり、
前命題の tails 構成で pointed exact 同値を得る。
$\alpha$ は存在証明のための選択であり、根を保つ同値を仮定する必要はない。
\end{proof}

\begin{example}[二行列を一行列に制限できない例]\label{ex:two-matrix-segre}
Segre 環 $B=k[x_{00},x_{01},x_{10},x_{11}]/(\det x)$ で、第一行の符号を変える
次数付き対合 $\gamma$ を取る。右 twist $B'=B^\gamma$ の積を
$a\star b=a\gamma^{\deg a}(b)$ とすると $K'=(1\otimes\gamma)K$ である。
$\gamma\otimes\gamma$ は $K$ を保つので $(1,\gamma)$ は上の二条件を満たす。
両辺とも適合表示である。実際、$\gamma$ を四変数多項式環へ同じ対合として
持ち上げると、twist 後も型 $(1,4,6,4,1)$ の AS 正則被覆を持ち、
$\det x$ の直線は安定なので正則正規二次切断を与える。
右加群の twist 同値は自由加群を自由加群に送り、$\gamma^2=1$ より
二重シフトと自然に可換する（右作用 $m\star b=m\gamma^{\deg m}(b)$ を比較すればよい）。
ゆえに通常シフトの平方は両側で反 Serre 移動となる。
一方、$x_{00}\star x_{10}=x_{00}x_{10}=-x_{10}\star x_{00}\ne0$ である。
従って可換な $B$ から $B'$ への次数付き環同型はなく、$g=h$ の解は存在しない。
\end{example}

% !TEX root = ../nc-quadric-classification.tex
\section{cubic 型による被覆の限界}
\label{sec:cubic-coverage}

\begin{theorem}[cubic 表示の判定]
\label{thm:cubic-coverage}
$\HH$ 対象 $(\cC,\OO)$ の任意の適合表示から
定理~\ref{thm:normal-helix} の二つの初期 MCM 型を取る。
次の条件は同値である。
\begin{enumerate}
\item $(\cC,\OO)$ はある Van den Bergh の三次三次元 AS 正則
$\mathbb Z$-algebra $D$ に対する $(\qgr D,\pi e_0D)$ と pointed exact 同値である。
\item 二つの初期型の少なくとも一方で、偶数添字間の隣接合成が階数四になる。
\end{enumerate}
特に命題~\ref{prop:nonstandard-rank-three-example} の対象は
いかなる cubic 表示でも添字零の構造対象を保って実現できない。
\end{theorem}

\begin{proof}
十分性は定理~\ref{thm:normal-helix} と系~\ref{cor:normal-cubic-defect}
から従う。必要性のため、cubic 表示があると仮定し、
$\pi e_iD$ の像を $Q_i\in\cC$ と書く。
cubic 正則性と \cite[Proposition 4.11, Theorem 6.8]{MN2025} により、
$tQ_i\simeq Q_{i-4}$、および
$\Hom(Q_j,Q_i)=D_{ij}$ である。
この適用に必要な正則性と周期性は \cite[arXiv v3, Definition 4.1.1,
Proposition 5.6.1]{VdB2011} と
\cite[Theorem 4.19]{MN2025} から従う。

$M=Q_{-2}$ と置く。二添字移動は $D$ の自己同型として実現されるので
$Q_{-2}$ は $\OO=Q_0$ の exact 自己同値による像である。
従って $M$ は rank one torsion-free exceptional である。
さらに cubic 型の成分次元から
$\Hom(\OO,t^nM)=D_{-4n-2,0}$ は $(2n+2)^2$ 次元である。命題~\ref{prop:first-orbit-object} の証明の Ext 消滅により、
$n\gg0$ では Euler 特性も同じ値を取る。
命題~\ref{prop:first-orbit-object} により $M\simeq s\OO$ となる。
従ってすべての整数 $j$ で $Q_{-2j}\simeq s^j\OO$ であり、
これらの同型は Hom と合成で定めた二つの添字付き環の同型
$D_{\mathrm{even}}\simeq\check B$ を与える（偶数添字は二で割る）。
ここでは cubic 側の二添字移動そのものが反 Serre の平方根であることは使わない。

$P_i=e_iD$ の偶数成分への制限を $P_i^{\mathrm{ev}}$ と書く。
上の同型により次数付き右 $B$-加群とみなせる。
奇数 $i$ に対する右単純加群の極小分解を偶数成分に制限すると
\[
 0\longrightarrow P_{i+4}^{\mathrm{ev}}
 \longrightarrow (P_{i+3}^{\mathrm{ev}})^2
 \longrightarrow (P_{i+1}^{\mathrm{ev}})^2
 \longrightarrow P_i^{\mathrm{ev}}\longrightarrow0
\]
となる。偶数の添字の制限は $B$ の自由シフトである。
$i$ を前後に四ずつ変えると、$P_i^{\mathrm{ev}}$ は自由加群への埋込みを
繰り返し持ち、その余核は次の自由加群へ埋め込める。
従ってこの列の二回の次元移動を反復し、
$B$ の有限単射次元を使えば $\Ext_B^j(P_i^{\mathrm{ev}},B)=0$ ($j>0$) となる。

$X=P_{-1}^{\mathrm{ev}}$ は次数零に二元生成され、
$H_X(z)=2/(1-z)^3$ であり、自由加群ではない。
定理~\ref{thm:normal-helix} の二型の記述により、非自由 MCM の
直和因子は少なくとも二元を必要とする。従って $X$ は
定理~\ref{thm:normal-helix} の二型の一方である。
$D_{-2,-1}$ は左乗法と制限により $X\to B(1)$ の射を与え、
$D_{-1,0}$ は $B\to X$ の射を与える。
それらの積は $D_{-2,0}\simeq B_1$ を張る。
実際、cubic 型では距離二に関係がなく、この隣接積は同型である。
これらは tails への移行後も同じ Hom と合成になるので、
$X$ の偶数項間の合成は階数四である。
\end{proof}

% !TEX root = ../nc-quadric-surface-classification.tex
\subsection{階数三の四周期表示の生成元と関係}
\label{subsec:rank-three-relations}

\begin{proposition}
\label{prop:rank-three-relations}
定理~\ref{thm:normal-helix} の階数三の場合に、
矢を $E_i\to E_j$ の向きで書き、道の積を合成で定める。
隣接する二元の矢に加えて、$E_{2i}\to E_{2i+2}$ ごとに一つの追加矢を取る。
最小関係は以下の三種類である。
\begin{enumerate}
\item $E_{2i}\to E_{2i+2}$ に一つ。二つの隣接矢の空間を通る合成の核で、
その二因子テンソルの階数は二である。
\item $E_{2i}\to E_{2i+3}$ に二つ。
$X_{i+1}$ の極小評価列の核の次数一部分を、
$\Hom(\pi B(i),\pi B(i+1))$ と $\Hom(\pi B(i+1),\pi X_{i+1}(i+1))$
を経由する道として表す。
\item $E_{2i+1}\to E_{2i+4}$ に二つ。
$X_i$ の双対評価列の核の最初の非零次数部分を、
$\Hom(\pi X_i(i),\pi B(i+1))$ と $\Hom(\pi B(i+1),\pi B(i+2))$
を経由する道として表す。
\end{enumerate}
各距離三の関係は、距離二の関係を先に課した後の道の空間から
一意な二次元部分空間として得られる。
\end{proposition}

\begin{proof}
四周期表示の単純加群の極小分解
（系~\ref{cor:normal-cubic-defect} の証明）により、最小関係の個数は
偶数項間の距離二で一、距離三で各二であり、他にはない。
距離二の偶数項間では階数三の合成に第四の基底を加えてあるので、
唯一の関係の追加矢における係数は零である。
この関係が階数二になることは
注意~\ref{rem:rank-three-matrix} の行列分解の議論による。

次に $E_{2i}\to E_{2i+3}$ の道を考える。
距離二の関係を先に課すと、道の空間は
$\Hom(\pi B(i),\pi B(i+1))\otimes\Hom(\pi B(i+1),\pi X_{i+1}(i+1))$
と同定できる。実際、距離三には三つの隣接矢からなる道と
追加矢を一つ含む道があり、前者に距離二の関係を適用すれば
最初の二矢の積と追加矢を合わせて $B_1$ の基底が得られる。
従って合成の核を求めるには、$X_{i+1}$ の評価列
\eqref{eq:normal-ruling-evaluation} を次数一で取ればよい。
その核は $(p(X_{i+1}))_0$ であり、二次元である。

$E_{2i+1}\to E_{2i+4}$ についても、
後二つの隣接矢の積と追加矢を合わせると $B_1$ の基底が得られる。
ここで $B$ への内部 Hom を取った評価列
\[
 0\to\Hom_B(p(X_i),B)\to B^2
       \to \Hom_B(X_i,B)(1)\to0
\]
を使うと、合成の核は左端の次数一成分である。
これも二次元であり、主張した全ての距離三の関係を与える。
\end{proof}

比較先の四周期環を逆に構成するとき、
この表示で距離二の関係を階数二に正規化してから
距離三の各二次元部分空間を変化させられる。
ここから得る候補の全次数の許容条件は、
極小分解が実際に完全であることと、周期偶数側が
命題~\ref{prop:first-orbit-object} の軌道を与えることを含む。
これらは有限次元の関係の個数からは従わない。

% !TeX root = ../nc-quadric-classification.tex
\section{階数三支持超平面の共通二変数環}
\label{sec:common-plane}

基礎体 $k$ は標数零の代数閉体とする。
次数一で生成された連結二次次数付き domain $B=T(V)/(K)$ について
\[
 \dim V=4,\qquad K\subset V\otimes V,\qquad \dim K=7
\]
を仮定する。従って $\dim B_2=9$ である。
線形行列 $M\in\operatorname{Mat}_2(V)$ の成分が張る部分空間を $U(M)$ と書く。

\begin{lemma}[行列積の四成分の独立性]
\label{lem:w-matrix-products}
$L_1,L_2\subset\operatorname{Mat}_2(k)$ を三次元部分空間とする。
trace pairing に関する $L_2^\perp$ が可逆行列で生成されるなら、
行列の乗法
\[
 L_1\otimes L_2\longrightarrow\operatorname{Mat}_2(k)
\]
は全射である。
従って $P\in\operatorname{Mat}_2(U_1)$ と $Q\in\operatorname{Mat}_2(U_2)$ の係数行列が
それぞれ $L_1,L_2$ の基底をなすとき、$PQ$ の四成分は
形式的テンソル空間 $U_1\otimes U_2$ で一次独立である。
\end{lemma}
\begin{proof}
全射でないなら、非零行列 $D$ があって
$\operatorname{tr}(DXY)=0$ が全ての $X\in L_1,Y\in L_2$ に対して成立する。
従って $DL_1\subset L_2^\perp=kC$ であり、$C$ は可逆である。
$D$ が可逆なら $\dim DL_1=3$ なので矛盾する。
$\operatorname{rank} D=1$ なら $DX$ は全て特異行列であるため $DL_1=0$ となるが、
左乗法 $X\mapsto DX$ の核は二次元であり、三次元の $L_1$ を含めない。
後半はこの全射の双対を取ったものである。
\end{proof}

\begin{lemma}[平方零行列の成分の関係]
\label{lem:w-invertible-annihilator}
$M\in\operatorname{Mat}_2(V)$ が $B$ 上で $M^2=0$ を満たし、$\dim U(M)=3$ とする。
四成分の唯一の一次関係を表す二次正方係数行列は可逆である。
従って $M^2$ の四成分を $V\otimes V$ の元として解釈すると、
それらは四次元部分空間 $R(M)\subset K\cap U(M)^{\otimes2}$ を張る。
\end{lemma}
\begin{proof}
関係の係数行列の階数が一なら、非零の定数行ベクトル $a$ と
列ベクトル $b$ によって $aMb=0$ と書ける。
$ab\ne0$ なら同時基底変換、すなわち $M$ の共役変換によって
この式を $M_{11}=0$ にできる。
このとき $M=\left(\begin{smallmatrix}0&x\\y&z\end{smallmatrix}\right)$
と書け、$M^2=0$ の第一対角成分から $xy=0$ となる。
domain 性により $x=0$ または $y=0$ で、成分空間の次元三に反する。
$ab=0$ なら同時基底変換で $M_{12}=0$ にできる。
このとき $M=\left(\begin{smallmatrix}x&0\\y&z\end{smallmatrix}\right)$
であり、二つの対角成分から $x^2=z^2=0$ となって同様に矛盾する。
関係の係数行列の転置は係数行列空間の trace-annihilator を生成するので、
補題~\ref{lem:w-matrix-products} を適用できる。
\end{proof}

\begin{lemma}[二次元部分空間にある関係]
\label{lem:w-binary-relations}
二次元部分空間 $W\subset V$ に対し、$\dim(K\cap W^{\otimes2})\le1$ である。
等号が成り立つ場合、非零の関係は二因子テンソルとして階数二である。
\end{lemma}
\begin{proof}
階数一の非零テンソル $w\otimes w'$ が $K$ にあれば、
$B$ で $ww'=0$ となって domain 性に反する。
もし $K\cap W^{\otimes2}$ に独立な二つの元 $r_0,r_1$ があれば、
$\det(sr_0+tr_1)$ は $(s,t)$ の斉次二次式である。
代数閉性により、この式が恒等的に零の場合も含めて、
非零の $(s,t)$ で行列式が零になる。
$sr_0+tr_1$ は非零なので階数一となり、矛盾する。
\end{proof}

\begin{theorem}[non-standard 二型の共通二変数環]
\label{thm:w-common-plane}
$M_1,M_2\in\operatorname{Mat}_2(V)$ が $B$ 上で $M_i^2=0$ を満たすとする。
$U_i=U(M_i)$ は三次元で、$U_1\ne U_2$ と仮定する。
$W=U_1\cap U_2$ および $R_i=R(M_i)$ と置けば、
\[
 \dim W=2,\qquad
 R_1\cap R_2=K\cap W^{\otimes2}=kr,\qquad
 \operatorname{rank} r=2,\qquad K=R_1+R_2.
\]
特に二つの平方零行列は $B$ の全ての二次関係を与える。
\end{theorem}
\begin{proof}
補題~\ref{lem:w-invertible-annihilator} により $\dim R_i=4$ である。
$R_i\subset K$ と $\dim K=7$ から $\dim(R_1\cap R_2)\ge1$ となる。
$U_1,U_2$ に適合する $V$ の基底を取れば
\[
 (U_1\otimes U_1)\cap(U_2\otimes U_2)
   =(U_1\cap U_2)\otimes(U_1\cap U_2)=W^{\otimes2}
\]
である。従って $R_1\cap R_2\subset K\cap W^{\otimes2}$ であり、
補題~\ref{lem:w-binary-relations} によって両辺は一次元に等しい。
さらに $\dim(R_1+R_2)=4+4-1=7=\dim K$ となる。
\end{proof}

\paragraph{正規切断への適用.}
系~\ref{cor:nonstandard-square-matrix} により、
正規切断 $B=A/(g)$ で $p(X)=X,p(Y)=Y$ のとき、評価行列は
$M_X^2=gC_X$, $M_Y^2=gC_Y$ ($C_X,C_Y\in\operatorname{GL}_2(k)$) と取れる。
$B$ に移せば平方零なので、双方が階数三かつ支持が異なる場合には
定理~\ref{thm:w-common-plane} が適用される。
$g$ が中央であることも $C_X,C_Y$ がスカラーであることも使用していない。

\begin{remark}[交差積の場合にも使える部分]
\label{rem:w-crossed-products}
二つの行列 $P,Q$ の支持が異なる三次元空間 $U_1,U_2$ で、
それぞれの係数行列空間の trace-annihilator が可逆行列で生成され、
$PQ=QP=0$ が $B$ で成立する場合も同じ結論が成り立つ。
この場合、四次元の関係空間は
$R_{12}\subset U_1\otimes U_2$ と $R_{21}\subset U_2\otimes U_1$ にあり、
それらの交わりが $W^{\otimes2}$ に含まれることを使う。
中央切断の通常の行列分解 $PQ=QP=gI$ はこの状況の候補となるが、
一般の正規切断では一方の積に正規化自己同型が入る。
後者への適用には、twist された支持の一致を別に確定する必要がある。
\end{remark}

\subsection{全次数での量子平面・Jordan 平面の埋込み}

\begin{proposition}
\label{prop:w-embedded-plane}
$W\subset B_1$ が二次元で $K\cap W^{\otimes2}=kr\ne0$ とする。
このとき $W$ の基底 $u,v$ を適切に選ぶと、$W$ が生成する $B$ の部分環は
次のいずれかと同型である。
\[
 Q_q=k\langle u,v\rangle/(vu-quv),\quad q\in k^\times,
 \qquad
 Q_J=k\langle u,v\rangle/(vu-uv-u^2).
\]
これは二次関係だけの同定でなく、自然な写像 $Q_q\to B$ または $Q_J\to B$
の全次数での単射を主張する。
\end{proposition}
\begin{proof}
補題~\ref{lem:w-binary-relations} により $r$ の係数行列は可逆である。
$r$ を対称部分と反対称部分に分ける。
対称部分の階数が二なら、代数閉体上の二変数非退化対称形式を
$uv+vu$ に正規化できる。従って $r=\alpha uv+\beta vu$ と書け、
係数行列の可逆性から $\alpha\beta\ne0$ である。これは $Q_q$ の関係である。
対称部分が零なら $r$ は $vu-uv$ の定数倍である。
対称部分の階数が一なら、$r=\alpha u^2+\beta(vu-uv)$ と書ける。
可逆性から $\beta\ne0$ なので、基底を定数倍して Jordan 平面の関係となる。

次に単射性を証明する。いずれの環 $Q$ も $k[u]$ の Ore 拡大で、
単項式 $u^iv^j$ ($i,j\ge0$) が基底をなし、domain である。
$u$ は正規なので可逆化できる。
$z=u^{-1}v\in Q[u^{-1}]$ と置けば、
\[
 zu=u\tau(z),\qquad
 \tau(z)=qz\ \ (Q=Q_q),\qquad
 \tau(z)=z+1\ \ (Q=Q_J)
\]
となる。従って次数 $n$ の任意の非零斉次元 $f$ は一意に
\[
 f=u^n p(z),\qquad 0\ne p\in k[z],\qquad \deg p\le n
\]
と書ける。実際、$u^{n-j}v^j$ をこの形に直すと、
$p$ は次数 $j$ の非零先頭係数を持つ多項式になり、三角的に基底を与える。
代数閉性から $p(z)=a\prod_{i=1}^d(z-\lambda_i)$ と書く。
恒等式
\[
 \prod_{i=1}^{d}\bigl(u f_i(z)\bigr)
 =u^d\prod_{i=1}^{d}\tau^{d-i}\bigl(f_i(z)\bigr)
\]
に $f_i(z)=\tau^{-(d-i)}(z-\lambda_i)$ を入れれば、
$u^dp(z)$ は $Q_1$ の非零元 $d$ 個の積の定数倍となる。
各 $f_i$ が一次式なので $u f_i(z)\in ku+kv$ である。
残る $u^{n-d}$ も合わせると、$f$ は非零一次元 $n$ 個の積である。

自然な全射 $Q\to k\langle W\rangle\subset B$ は次数一で単射である。
その核が非零なら非零斉次元 $f$ を取り、上の一次因子分解を写すと、
$B$ で非零一次元の積が零になる。これは domain 性に反する。
\end{proof}

\begin{remark}[ここからはまだ正規化できない]
$Q_q$ の $u,v$ は $Q_q$ の中で正規であり、$Q_J$ の $u$ は $Q_J$ の中で正規である。
しかし $B$ 内での正規性は、この埋込みの命題からは従わない。
また $q\ne1$ の量子平面や Jordan 平面の場合、$W$ の線形基底変更だけで
二つの生成元を $B$ の中心に入れることはできない。
そのような基底があれば $W$ の全ての元が互いに交換するが、
命題の全次数単射に反するからである。
Zhang twist や pointed 同値によって別の表示へ移る可能性は残る。
\end{remark}

\subsection{同一支持と中央化についての障害}

\begin{lemma}[同一支持の場合に確定すること]
\label{lem:w-equal-support}
三次元部分空間 $U\subset B_1$ について
$\dim(K\cap U^{\otimes2})\le4$ である。
従って定理~\ref{thm:w-common-plane} の平方零行列で $U_1=U_2=U$ なら、
\[
 R_1=R_2=K\cap U^{\otimes2}.
\]
\end{lemma}
\begin{proof}
$\mathbb P(U)\times\mathbb P(U)$ の Segre 像は
$\mathbb P(U^{\otimes2})=\mathbb P^8$ 内の四次元射影多様体である。
$K\cap U^{\otimes2}$ が五次元以上なら、その射影化の次元は四以上であり、
射影次元定理によりこの Segre 像と交わる。
この交点は非零の階数一テンソルを与え、domain 性に反する。
後半は各 $R_i$ が四次元であることから従う。
\end{proof}

\paragraph{残る点.}
$R_1=R_2$ から二つの余核が同型になるかは、この次元論だけでは分からない。
従って二つの異なる MCM 型で支持が異なるという仮定を外してはいけない。
次の具体的課題は、三次元の同じ成分空間と同じ四次元の平方関係空間を持つ
二つの行列について、定数行列による左右同値を証明するか、反例を得ることである。

\paragraph{中央化は適合表示の中ですら自動でない.}
命題~\ref{prop:nonstandard-rank-three-example} で使う
\[
 B_0=k[u,v]\langle c,d\rangle/(c^2-uv,d^2-uv)
\]
を考える。この表示では $u,v$ は中央である。
$t\in k^\times$ に対して
\[
 \sigma_t(u)=t^2u,\quad \sigma_t(v)=v,\quad
 \sigma_t(c)=tc,\quad\sigma_t(d)=td
\]
は次数付き自己同型である。
$a\star b=a\sigma_t^{\deg a}(b)$ の規約の Zhang twist を $B_t$ と書くと、
関係は
\[
\begin{gathered}
 vu=t^2uv,\qquad cu=tuc,\qquad cv=t^{-1}vc,\\
 du=tud,\qquad dv=t^{-1}vd,\qquad
 c^2=tuv,\qquad d^2=tuv.
\end{gathered}
\]
ここでは積記号を再び通常の並置で書いた。
行列
\[
 M_r=\begin{pmatrix}r&u\\-tv&-tr\end{pmatrix},\qquad r=c,d,
\]
はそれぞれ $M_r^2=0$ を満たす。
この二つの支持は $\langle u,v,c\rangle$, $\langle u,v,d\rangle$ であり、
交わりはパラメータ $q=t^2$ の量子平面を生成する。
$t^2\ne1$ なら、同じ $B_t$ の中で $W$ の基底を変えて中央化することは不可能である。

特に $t=\mathrm i$ は命題~\ref{prop:nonstandard-rank-three-example} の反 Serre 平方根に適合した表示そのものであり、
このとき $vu=-uv$ である。
従って「$H_2$ の適合表示なら $W$ を線形基底変更だけで中央化できる」
という主張は、既存例によって既に否定される。
一方 $u,v$ はこの例では $B_t$ 全体でも正規なので、
一般の $W$ の正規化という課題の反例にはなっていない。
任意の $t$ について次数シフト自体が反 Serre の平方根になるとは主張していない。

% !TeX root = ../nc-quadric-surface-classification.tex
\section{standard 型の帰着と非 cubic 部分の完全分類}
\label{sec:standard-rank-three}

\begin{proposition}[standard 二型の両側零積]
\label{prop:standard-mutual-zero-matrices}
$B=A/(g)$ を $\HH$ の適合表示とし、二つの非自由 MCM 型を $X,Y$ とする。
$\Omega X(1)\simeq Y$, $\Omega Y(1)\simeq X$ の場合、
両者の極小線形表示行列 $M,N\in M_2(B_1)$ は
\[
 MN=NM=0
\]
を同時に満たすように取れる。
両合成階数が三なら、両行列の成分間の唯一の一次関係の係数行列は可逆である。
\end{proposition}
\begin{proof}
評価列を
\[
 0\longrightarrow Y(-1)\xrightarrow{i_Y}F_X
 \xrightarrow{q_X}X\longrightarrow0,
 \qquad
 0\longrightarrow X(-1)\xrightarrow{i_X}F_Y
 \xrightarrow{q_Y}Y\longrightarrow0
\]
と取る。ここで $F_X,F_Y$ はそれぞれ二元の次数零自由加群である。
$M=i_Yq_Y(-1):F_Y(-1)\to F_X$、
$N=i_Xq_X(-1):F_X(-1)\to F_Y$ と定めれば、
$q_Xi_Y=q_Yi_X=0$ より、適切なシフトを付した二つの積は零である。
また $q_X,q_Y$ は全射なので、これらの余核はそれぞれ $X,Y$ である。
従ってこれは極小線形表示を与える。

仮に $M$ の唯一の一次関係の係数行列の階数が一なら、
$M'=PMQ$、$N'=Q^{-1}NP^{-1}$ と同時に基底変更し、
$M'_{11}=0$ とできる。この変更後も $M'N'=N'M'=0$ である。
$M'=\left(\begin{smallmatrix}0&a\\b&c\end{smallmatrix}\right)$
と書くと、三次元の成分空間より $a\ne0$ である。
$M'N'=0$ の第一行と domain 性から $N'$ の第二行が零となり、
$N'$ の成分空間が三次元であることに反する。
$N$ の関係も同様である。
\end{proof}

\begin{remark}[正規切断での持上げの範囲]
$A_1=B_1$ による一意な持上げを同じ $M,N$ で表す。
各行列は線形行列分解の最初の行列であるから $A$ 上で単射であり、
\[
 MN=gC_X,\qquad NM=gC_Y,\qquad C_X,C_Y\in\mathrm{GL}_2(k)
\]
となる。定数行列が特異なら積が非零の定数列を殺すので単射性に反する。
$ag=g\nu(a)$ の規約で結合則を使うと
\[
 \nu(M)=C_XMC_Y^{-1},\qquad
 \nu(N)=C_YNC_X^{-1}
\]
を得る。一般の正規切断で二つの定数行列を同時に恒等へ直したとは主張しない。
以下の二次関係の議論は $B$ 上の両側零積のみを使うため、この障害を受けない。
\end{remark}

\begin{proposition}[standard 型の支持が異なる場合]
\label{prop:standard-crossed-common-plane}
命題~\ref{prop:standard-mutual-zero-matrices} の両階数三の場合で、
$U=U(M)$、$U'=U(N)$ が異なるとする。
$R_{12}$ を形式的な $MN$ の四成分の張る二次関係空間、
$R_{21}$ を $NM$ の同様の空間と置く。
このとき $W=U\cap U'$ は二次元であり、
\[
 \dim R_{12}=\dim R_{21}=4,\qquad
 R_{12}\cap R_{21}=K\cap W^{\otimes2}=kr,\qquad
 \operatorname{rank}r=2,\qquad K=R_{12}+R_{21}
\]
が成り立つ。$W$ が生成する部分環は全次数で量子平面または Jordan 平面である。
\end{proposition}
\begin{proof}
前命題の非退化な一次関係と行列乗法の全射性
（補題~\ref{lem:w-matrix-products}）により、各 $R$ は四次元である。
また $R_{12}\subset U\otimes U'$、$R_{21}\subset U'\otimes U$ であり、
両テンソル空間の交わりは $W^{\otimes2}$ である。
$\dim K=7$ と、domain の二次元部分空間には高々一つの非退化二次関係しか
存在しないことから、主張する次元と等式が従う。
全次数の単射は命題~\ref{prop:w-embedded-plane} による。
\end{proof}

\begin{lemma}[一つの支持を持つ零積行列は conic の twist を与える]
\label{lem:standard-one-support-conic}
$B$ を連結次数付き domain とし、$M,N\in M_2(B_1)$ が $MN=0$ を満たすとする。
$U(M)=U(N)=U$ が三次元であると仮定する。
このとき $U$ が生成する部分環は可換滑らか conic の
次数付き自己同型による Zhang twist である。
また $M,N$ は定数行列による左右変換で互いに移り、
$B$ 上の次数付き余核は同型である。
この命題には $NM=0$ も $\dim B_1=4$ も必要ない。
\end{lemma}
\begin{proof}
まず、各行列の四成分の唯一の一次関係は可逆な係数行列を持つ。
実際、$M$ の関係の係数行列の階数が一なら、
$M\mapsto PMQ$, $N\mapsto Q^{-1}NP^{-1}$ の同時基底変更で
$M_{11}=0$ にできる。
$M=\left(\begin{smallmatrix}0&a\\b&c\end{smallmatrix}\right)$ と書くと
$a\ne0$ であり、$MN=0$ の第一行と domain 性から $N$ の第二行が零となる。
これは $N$ の成分空間が三次元であることに反する。
$N$ の関係の場合は $N_{11}=0$ とし、積の第一列を使えば同様である。

それぞれの係数行列空間の trace-annihilator を用いて
\[
 M=M_0G,\qquad N=N_0H,\qquad G,H\in\mathrm{GL}_2(k),
 \qquad \operatorname{tr}M_0=\operatorname{tr}N_0=0
\]
と書く。
$M_0=\left(\begin{smallmatrix}a&b\\c&-a\end{smallmatrix}\right)$ とすると
$a,b,c$ は $U$ の基底である。
$\psi:U\to U$ を
\[
 \psi(M_0)=G N_0G^{-1}
\]
で定める。右辺も三次元の traceless 成分空間を持つので、$\psi$ は可逆である。
形式的テンソル積として
\[
 MN=M_0\psi(M_0)GH
\]
であるから、$MN$ の四成分の張る関係空間は
\[
 R=(1\otimes\psi)K_0,\qquad
 K_0=\bigwedge^2U+kq,\qquad q=a^2+bc
\]
である。ここでは $q$ を対称テンソルとして代表させている。
$K_0$ は可換 conic $C_0=k[a,b,c]/(a^2+bc)$ の四次元二次関係空間である。

$\sigma=\psi^{-1}\in\mathrm{GL}(U)$ を三変数可換多項式環
$P=\operatorname{Sym}(U)$ の次数付き自己同型へ延長する。
この段階では $\sigma$ が $q$ を保つとは仮定しない。
$A=P^\sigma$ と置けば、$A$ の三つの二次交換関係は
$(1\otimes\psi)\bigwedge^2U$ であるから、$A\to B$ の準同型がある。
underlying vector space $A_2=P_2$ の中で、残る関係はちょうど元 $q$ である。
domain の三次元部分空間の二次関係は高々四次元である
（補題~\ref{lem:w-equal-support} の Segre 次元論）。
従ってこの準同型の次数二の核はちょうど $kq$ である。

$\sigma$ が conic を保つことを、domain 性から導く。
任意の非零 $u\in U$ に対し、$A=P^\sigma$ の積の規約から
\[
 u\star q=u\,\sigma(q)
 =\sigma(q)\,u
 =\sigma(q)\star\sigma^{-2}(u).
\]
中央の積は underlying 可換多項式環 $P$ の積である。
$B$ では $q=0$ であり、$\sigma^{-2}(u)$ の像は非零の次数一元である。
従って domain 性により $\sigma(q)$ も零に写る。
次数二の核が $kq$ なので
\[
 \sigma(q)\in k^\times q
\]
となる。よって $\sigma$ は $C_0$ の自己同型に降り、
$A/(q)=C_0^\sigma$ から $B$ への写像を得る。

この写像は次数一で単射であり、全次数でも単射である。
実際、$C_0=k[s,t]^{(2)}$ の全斉次元は代数閉性により次数一元の積に分解でき、
$C_0^\sigma$ でも因子に $\sigma$ の逆冪を掛ければ同様である。
核の非零斉次元は $B$ における非零一次元の零積を与えてしまう。
従って $k\langle U\rangle=C_0^\sigma$ である。

最後に余核を比較する。
$MN=0$ を可換環 $C_0$ の積へ戻すと $M\sigma(N)=0$ である。
$M$ も $\sigma(N)$ も非零なので、商体上では両方の階数が一であり、
従って $\det M=\det N=0$ となる。
$a=st,b=s^2,c=-t^2$ と書くと、両行列は二変数二次式からなる階数一行列であり、
成分は全ての二次式の三次元空間を張る。
一意分解性により各行列は一次式の二成分列と一次式の二成分行の積に分解する。
それぞれの成分は独立である。
具体的には、成分の共通因子は $s^2,t^2$ を張ることから定数であり、
列または行の次数が零なら成分空間が高々二次元になるので、両方の次数は一である。
従って各行列は
\[
 \begin{pmatrix}st&s^2\\-t^2&-st\end{pmatrix}
 =\begin{pmatrix}s\\-t\end{pmatrix}\begin{pmatrix}t&s\end{pmatrix}
\]
と定数行列による左右変換で同値である。
この一次元の等式は $B$ 上でも成立し、余核の同型を与える。
\end{proof}

\begin{corollary}[standard 二型の支持も異なる]
\label{cor:standard-supports-distinct}
$\HH$ の standard 適合表示で、二つの MCM 型の合成階数がともに三なら、
それらの支持超平面は異なる。
\end{corollary}
\begin{proof}
評価列から得た行列 $M,N$ は $MN=NM=0$ を満たし、
余核は異なる二つの MCM 型である。
同じ支持を持てば補題~\ref{lem:standard-one-support-conic} が余核の同型を与えて矛盾する。
\end{proof}

\begin{theorem}[standard 型の異なる支持の同時 untwist]
\label{thm:standard-distinct-untwist}
$B=T(V)/(K)$ を noetherian 連結二次次数付き domain とし、
$\dim_kB_n=(n+1)^2$ とする。
三次元の異なる成分空間 $U,U'\subset V$ を持つ行列
$M,N\in M_2(B_1)$ が $MN=NM=0$ を満たすと仮定する。
このとき $B$ の次数付き自己同型 $\tau$ が存在して、
$B^\tau$ は中央二重平方環
\[
 k[u,v]\langle z_1,z_2\rangle/
 (z_1^2-f_1(u,v),\ z_2^2-f_2(u,v))
\]
と同型になる。ここで $u,v$ は中央、$f_1,f_2$ は非退化二次形式である。
従って standard な両階数三表示で支持が異なる場合も、
pointed exact 同値の意味で定理~\ref{thm:ore-anticommuting-central} の族に帰着する。
\end{theorem}
\begin{proof}
命題~\ref{prop:standard-mutual-zero-matrices} と同じ議論で、
両係数行列空間の trace-annihilator は可逆である。
それらを $kD_M,kD_N$ と書く。
$H^2=D_MD_N$ となる可逆行列 $H$ を取り、$Q=HD_N^{-1}$ と置けば
$QD_NQ=D_M$ である。
$M\mapsto MQ$, $N\mapsto Q^{-1}N$ により、二つの trace-annihilator を一致させる。
この同時変更は両側零積も各支持も保つ。

共通の一次関係から、成分の対応が線形同型
$\varphi:U\to U'$、$\varphi(M)=N$ を定める。
共通の係数 hyperplane の trace-annihilator を $kD$ とし、
$G=D^{-1}$ と置くと
\[
 M=M_0G,\qquad
 M_0=\begin{pmatrix}a&b\\c&-a\end{pmatrix},\qquad
 C_0=k[a,b,c]/(a^2+bc)
\]
と書ける。
$\sigma_0(M_0)=G^{-1}M_0G$ は $C_0$ の次数付き自己同型を定める。
$K_0\subset U^{\otimes2}$ を $C_0$ の二次関係空間とする。
形式的な $M^2$ の四成分は
$R_0=(1\otimes\sigma_0^{-1})K_0$ を張る。
これは $C_0^{\sigma_0}$ における
$M\star M=M\sigma_0(M)=M_0^2G^2=0$ と、
四成分の独立性から従う。
ここで $M^2=0$ を $B$ で仮定しているわけではない。

従って両側零積の二次関係は
\begin{equation}
 R_{12}=(1\otimes\varphi\sigma_0^{-1})K_0,
 \qquad
 R_{21}=(\varphi\otimes\sigma_0^{-1})K_0.
 \label{eq:standard-cross-conic-kernels}
\end{equation}
$W=U\cap U'$ と置くと、命題~\ref{prop:standard-crossed-common-plane} の次元論より
$K=R_{12}+R_{21}$ および
$R_{12}\cap R_{21}=K\cap W^{\otimes2}=kr$ であり、$r$ は階数二である。
$Q_W=k\langle W\rangle$ は全次数で量子平面または Jordan 平面である。

可換 conic を $C_0=k[s,t]^{(2)}$ と書く。
まず $W\subset U=k[s,t]_2$ が基点を持たないことを示す。
仮に $W=h\,k[s,t]_1$ とする。
$r$ を \eqref{eq:standard-cross-conic-kernels} の第一の kernel で解釈すると、
\[
 \ker\left(W\otimes\sigma_0\varphi^{-1}(W)
       \longrightarrow k[s,t]_4\right)\ne0
\]
である。$r$ の階数二より第二因子の空間も
$\sigma_0\varphi^{-1}(W)=h'k[s,t]_1$ と書ける。
実際、$hs\,A+ht\,B=0$ は $A=t\ell$, $B=-s\ell$ を与え、
$A,B$ は独立なので第二空間全体が $\ell k[s,t]_1$ となる。

$\sigma_0\varphi^{-1}(q)=hh'$ となる非零元 $q\in W$ を取る。
任意の $z\in U=k[s,t]_2$ について
\[
 z\cdot\sigma_0\varphi^{-1}(q)=hh'z
 \in W\cdot\sigma_0\varphi^{-1}(W)
\]
である。
第一の $R_{12}$ は写像
$U\otimes U'\to k[s,t]_4$, $z\otimes z'\mapsto z\sigma_0\varphi^{-1}(z')$
の全 kernel なので、$B$ において $zq\in(Q_W)_2$ となる。
従って $U\subset\operatorname{Frac}Q_W$ である。

第二の $R_{21}$ において、$W$ の二つの像は
$\varphi^{-1}(W)$ と $\sigma_0(W)$ であり、両方とも共通一次因子を持つ pencil である。
同じ議論を、写像
$U'\otimes U\to k[s,t]_4$, $z'\otimes z\mapsto\varphi^{-1}(z')\sigma_0(z)$
へ適用すると、適切な非零 $q'\in W$ に対し
$z'q'\in(Q_W)_2$ ($z'\in U'$) を得る。
よって $U'\subset\operatorname{Frac}Q_W$ でもある。
これらの分数は全て次数一である。

以上の分数を $B$ の商斜体の中で取ると、$B$ は $Q_W$ の斉次商環の
次数付き部分環になる。
量子平面および Jordan 平面の斉次商環は
$k(\mathbb P^1)[T,T^{-1};\psi]$、$\psi\in\operatorname{Aut}\mathbb P^1$
の形である。
補題~\ref{lem:w-rational-growth} より $\dim_kB_n=O(n)$ となり、
仮定した $(n+1)^2$ と矛盾する。
従って $W$ は基点を持たない。

$W$ の互いに素な二元を $w_1,w_2$ とする。
$r$ が与える $w_1A+w_2B=0$ に一意分解性を使うと、
$A=\lambda w_2$, $B=-\lambda w_1$ となる。
従って
\[
 \sigma_0\varphi^{-1}(W)=W.
\]
特に $\varphi^{-1}(W)$ も基点を持たず、
第二の $R_{21}$ に同じ議論を適用して
$\sigma_0(W)=\varphi^{-1}(W)$ を得る。
そこで
\[
 T_1=\varphi\sigma_0^{-1}:U\longrightarrow U',\qquad
 T_2=\sigma_0^{-1}\varphi^{-1}:U'\longrightarrow U
\]
はいずれも $W$ を保つ。
$K_1=(\varphi\otimes\varphi)K_0$ と書けば
$R_{12}=(1\otimes T_1)K_0$、$R_{21}=(1\otimes T_2)K_1$ である。
共通の関係は
\[
 kr=(1\otimes T_1)(\bigwedge^2W)
    =(1\otimes T_2)(\bigwedge^2W)
\]
なので $T_1|_W$ と $T_2|_W$ は比例する。
$N$ を定数倍すると $\varphi$ が同じ倍率で変わり、
$T_1,T_2$ は互いに逆の倍率で変わる。
基礎体上で平方根を取ることにより、両制限を一致させられる。

$\tau|_U=T_1$, $\tau|_{U'}=T_2$ として $V$ 上へ貼り合わせる。
これは $U/W$ と $U'/W$ を交換する可逆線形写像である。
$\sigma_0$ が $C_0$ の自己同型であることを使うと
\begin{align*}
 (\tau\otimes\tau)R_{12}
  &=(\varphi\sigma_0^{-1}\otimes\sigma_0^{-2})K_0
    =(\varphi\otimes\sigma_0^{-1})K_0=R_{21},\\
 (\tau\otimes\tau)R_{21}
  &=(\sigma_0^{-1}\otimes\varphi\sigma_0^{-2})K_0
    =(1\otimes\varphi\sigma_0^{-1})K_0=R_{12}.
\end{align*}
従って $\tau$ は $B$ の次数付き自己同型である。

$B^\tau$ の二次関係は $(1\otimes\tau^{-1})K$ なので、
\[
 (1\otimes\tau^{-1})R_{12}=K_0,\qquad
 (1\otimes\tau^{-1})R_{21}=K_1
\]
となる。よって二つの支持の部分環は同時に可換 conic になり、$W$ は全環で中央となる。
定理~\ref{thm:w-simultaneous-untwist} の証明末尾の平方完成の議論により、
非退化な $f_1,f_2$ を持つ中央二重平方表示を得る。
Zhang 同値は構造対象を保つ。
\end{proof}

\begin{theorem}[非 cubic 部分の完全分類]
\label{thm:noncubic-classification}
$\HH$ のうち pointed cubic 表示を持たない対象の pointed exact 同値類は、
$\PP^1$ 上の二つの被約な二点集合の非順序対を、
$\PGL_2$ の同時作用で割った集合と一対一に対応する。
具体的には、二対が一致する一点、一点を共有する一点、
および四点が異なる場合の
\[
 \lambda\in k\setminus\{0,1\},\qquad \lambda\sim\lambda^{-1}
\]
からなる。各対象は、二点集合を零点とする非退化二次形式 $f_1,f_2$ を用いた
中央二重平方環の pointed 圏として実現される。
non-standard な適合表示を持つものは、二対が一致する一点だけである。
\end{theorem}
\begin{proof}
定理~\ref{thm:cubic-coverage} により、cubic 表示を持たないことは
二つの MCM 型の合成階数がともに三であることと同値である。
その二型に対する $p=\Omega(1)$ は恒等か交換かのいずれかである。
恒等の場合は定理~\ref{thm:w-simultaneous-untwist}、交換の場合は
命題~\ref{prop:standard-mutual-zero-matrices}、
系~\ref{cor:standard-supports-distinct}、
定理~\ref{thm:standard-distinct-untwist} により、中央二重平方環へ帰着する。
これらの Zhang 同値は構造対象を保つ。
逆の構成、全ての $\HH$ 条件と cubic 被覆不能性、
二つの根対による必要十分な同値判定は
定理~\ref{thm:ore-anticommuting-central} とその直後の系による。
non-standard な適合表示についての最後の主張は
定理~\ref{thm:w-compatible-singleton} である。
\end{proof}

% !TEX root = ../nc-quadric-surface.tex
\section{重み付き四周期表示の偶数環}
\label{sec:weighted-inverse}

以下、$D$ を局所有限で正に次数付けされた連結 $\mathbb Z$-algebra とし、
$P_j=e_jD$ と書く。$D$ は左右 noetherian とし、右単純加群 $S_j$ は
\[
 0\longrightarrow P_{j+4}\longrightarrow
 P_{j+3}^2\oplus P_{j+2}^{d_j}\longrightarrow
 P_{j+1}^2\oplus P_{j+2}^{d_j}\longrightarrow
 P_j\longrightarrow S_j\longrightarrow0
 \tag{*}\label{eq:weighted-inverse-rank-resolution}
\]
という極小分解を持つとする。ここで $d_j=0$（$j$ 奇数）、
$d_j\in\{0,1\}$（$j$ 偶数）、$d_{j+2}=d_j$ とする。
この節ではこの完全性を仮定して逆向きに何が従うかを調べる。

\begin{proposition}[周期偶数環の Koszul 性]
\label{prop:weighted-even-koszul}
$D$ の偶数部分を二添字の周期移動で次数付けして得られる環を $B$ とする。
すると $B$ は左右 noetherian、一次生成、二次 Koszul であり、
$H_B(z)=(1+z)/(1-z)^3$ である。
\end{proposition}

\begin{proof}
極小分解の第一項から、$D$ は隣接する二つの生成元と
偶数添字から始まる距離二の生成元で生成される。
偶数添字間の道は長さ一か二で次の偶数添字を通るから、
$B$ は一次生成である。偶数部分の斉次右イデアルを
$P_j$（$j$ 偶数）の部分加群として $D$ へ拡張し、
偶数成分に制限すると元のイデアルに戻る。
実際、拡張で最終添字が偶数になる項に使われる積は
両端とも偶数の成分だけである。よって斉次右イデアルの昇鎖は停止する。
左側も同様であり、一般のイデアルには最高次の初期項を用いる。

$h_{jk}=\dim D_{jk}$ と書く。極小分解の各成分で Euler 標数を取り、
同じ添字に現れる距離二の項を消去すると、
\[
 h_{jk}-2h_{j+1,k}+2h_{j+3,k}-h_{j+4,k}=\delta_{jk}.
\]
$h_{jk}=0$（$j>k$）から順に解けば、$h_{jk}$ は $k-j$ のみの関数となり、
その母関数は $(1-2z+2z^3-z^4)^{-1}$ である。
偶数差 $2n$ の係数は $(n+1)^2$ になる。

奇数 $i$ に対し、$N_i=P_i|_{\mathrm{even}}$ と置く。
式 \eqref{eq:weighted-inverse-rank-resolution} を偶数成分に制限すると
$S_i$ は零になるため、
\[
 0\to N_{i+4}\to (P_{i+3}|_{\mathrm{even}})^2
   \to(P_{i+1}|_{\mathrm{even}})^2\to N_i\to0
\]
を得る。$i$ を四ずつ増やしてこの列をつなぐと、
$N_i$ の自由分解の各項は、最初の生成次数から一つずつ増加する次数で生成される。
特に $N_1$ の左端の生成次数は $1$、$N_3$ では $2$ である
（偶数添字 $2j$ を次数 $j$ と数える）。

次に $j=0$ の式 \eqref{eq:weighted-inverse-rank-resolution} を
偶数成分に制限し、$k$ の分解を得る。
複体の第 $p=0,1,2,3$ 項は順に
$B$、$N_1^2\oplus B(-1)^{d_0}$、
$N_3^2\oplus B(-1)^{d_0}$、$B(-2)$ である。
この複体に $-\otimes_B^{\mathbf L} k$ を適用した
スペクトル系列の第一頁では、各第 $p$ 項の
$\operatorname{Tor}_q^B(-,k)$ は、上で得た線形分解により
内部次数が $p+q$ 以下にある。
従って $\operatorname{Tor}_n^B(k,k)$ は内部次数 $n$ 以下にある。
一方、$B$ は連結一次生成なので、極小分解から同じ Tor は内部次数 $n$ 以上にある。
よって $B$ は Koszul であり、特に二次である。
\end{proof}

\begin{proposition}[残る許容条件への還元]
\label{prop:weighted-inverse-reduction}
命題~\ref{prop:weighted-even-koszul} の仮定に加え、
$D$ を四周期、次元三・parameter 四の AS 正則環とする。
偶数部分の次数付けを与える周期移動の二乗は、
$D$ の逆 Nakayama 同型を偶数部分に制限したものと
対角型自己同型を除いて一致するとする。
このとき、次の条件が成立すれば $(\qgr D,\pi e_0D)\in\HH$ である。
\begin{enumerate}
\item $B$ は domain である。
\item ある正次数の非零斉次元 $f\in B$ に対して
$\pi(B/fB)$ は $\qgr B$ で有限長でない。
\item 定理~\ref{thm:finite-ambient-test} の有限階数条件を満たす
超平面がある。
\end{enumerate}
これらは、この周期移動を適合根とする逆構成に必要でもある。
\end{proposition}

\begin{proof}
偶数部分への制限により $\qgr D\simeq\qgr B$ となる。
ここでは制限は exact であり、左随伴は偶数頂点の表現可能加群を
対応する $P_{2j}$ に送る加群の延長である。
この随伴の単位は偶数成分で同型であり、
余単位の核と余核は偶数成分が零となる。
偶数成分が零の有限生成加群は、奇数始点の任意の正距離の道が
最初に偶数頂点を通るため、有限次元である。
また有限生成加群の偶数部分も有限生成であり、
奇数頂点に置かれた生成元については、次の偶数頂点への二本の矢でその制限を生成できる。
このためこれらの操作は tails において互いに逆である。
添字零の表現可能加群は $\pi B$ に移る。

\cite[Theorems 4.19, 6.8 and 6.11, Proposition 4.11]{MN2025} から、
$D$ 側の標識は exceptional で、四添字の Nakayama 移動が
反 Serre 関手となる。仮定した対角型の差は加群の各成分を定数倍すれば
自然同型になるので、偶数部分のシフトはその平方根である。
$\qgr D$ の豊富な表現可能加群列は、奇数項を次の偶数項から
生成できるため、偶数部分列でも豊富である。
表現可能加群間の Hom が $D$ の成分と一致することから、
この根の切断環は全次数で $B$ となる。
従って条件 (3) が所要の正則正規被覆を与える。

条件 (1) を仮定する。条件 (3) の被覆からすべての有限生成加群の
Hilbert 多項式が得られ、次数が高々一の対象は $\cC_1$ に入る
（命題~\ref{prop:sklyanin-h2} と同じ長さの議論）。
$f$ は正規である必要はなく、
$f^jB/f^{j+1}B\simeq(B/fB)(-j\deg f)$ が domain 性から従う。
条件 (2) のとき、この列は $\cC/\cC_0$ でも真に下降するから
$\pi B\notin\cC_1$ である。
$\pi B$ の任意の非零部分対象が $\pi B$ のシフトを含むことと、
それによる余商の Hilbert 多項式の次数が高々一であることから、
rank one torsion-free 性が従う。

逆に、適合根を得たなら $B$ の domain 性と被覆の有限条件は既に示した。
もしすべての非零斉次 $f$ に対して $\pi(B/fB)$ が有限長なら、
$\pi B$ の任意の非零部分対象による商は有限長となる。
$\cC/\cC_0$ の部分対象を持ち上げると、
そこにおける $\pi B$ の像は零または単純となり、
$\pi B\in\cC_1$ となって矛盾する。
\end{proof}

この還元では $D$ の四周期最小分解の完全性と noetherian 性を仮定した。
残る分類課題は、前節の距離二・距離三の関係からこれらを保証する条件、
および周期偶数環の domain 性と条件 (2) の具体的な判定である。

% !TeX root = ../nc-quadric-classification.tex

\section{整な楕円四つ組からの構成}
\label{sec:elliptic-inverse}

$D$ を Van den Bergh の三次三次元正則 $\mathbb Z$-algebra とする。
$P_i=e_iD$ を右加群、$S_i$ をその単純商と書き、
$V_i=D_{i,i+1}$、$R_i$ を三次関係空間、
$W_i=k w_i=(V_i\otimes R_{i+1})\cap(R_i\otimes V_{i+3})$ とする。
\cite[arXiv v3, \S 5.2]{VdB2011} に従い、$w_i$ の非零の代表を選び、
\[
 w_i=\sum_j v_j\otimes r_j,\qquad
 w_{i+1}=\sum_j r_j\otimes\theta_i(v_j)
\]
で次数 $+4$ の周期移動 $\theta$ を定める。各 $w_i$ のスカラー変更により、
$\theta$ は $D_{ij}$ に $c_i/c_j$ を掛ける対角型の自己同型まで定まる。

\begin{lemma}[四つの添字の移動と Serre 関手]
\label{lem:cubic-serre-period}
$\theta$ は、Mori--Nyman の Nakayama 同型
$\nu:D\to D(-4)$ の逆と上記の対角型自己同型を除いて一致する。
従って $\qgr D$ 上で $(4)_\theta\simeq S^{-1}[2]$ である。
\end{lemma}

\begin{proof}
右単純加群の極小分解を $F^i$ と書く。その項は
\[
 F^i_0=P_i,\quad F^i_1=V_i\otimes P_{i+1},\quad
 F^i_2=R_i\otimes P_{i+3},\quad F^i_3=W_i\otimes P_{i+4}
\]
であり、微分はテンソルの包含と積で与えられる
\cite[arXiv v3, Theorem 5.5.7]{VdB2011}。
特に $\Ext_D^q(S_i,P_j)$ は $(q,j)=(3,i+4)$ 以外で零である。
局所コホモロジー $R^3\tau(D)_{j,i}$ は
$\Ext^3_D(P_i/(P_i)_{\ge N},P_j)$ の直極限である。
これは torsion 部分の第 $i$ 成分を計算する
$\varinjlim_N\Hom_D(P_i/(P_i)_{\ge N},-)$ を導来すればよい。
右からの作用は始域における射での引き戻し、左からの作用は終域における射との合成である。

$M_i=P_i/(P_i)_{\ge i+2}$ とする。
$j=i+4,i+5$ について $\Ext^3_D(M_i,P_j)\to R^3\tau(D)_{j,i}$ は同型である。
さらに切り詰め次数を上げたときの核の単純因子は $S_l$ ($l\ge i+2$) なので、
上記の消滅と長完全列を用いて確かめられる。
短完全列 $0\to V_i\otimes S_{i+1}\to M_i\to S_i\to0$ に対し、
その拡大類を表す $F^i_1\to V_i\otimes S_{i+1}$ は自然な商写像である。
この写像を $V_i\otimes F^{i+1}$ に持ち上げると、
第一段階は $V_i\otimes P_{i+1}$ の恒等写像、第二段階は
$R_i\subset V_i\otimes V_{i+1}\otimes V_{i+2}$ の最後の因子を
$P_{i+3}$ への左乗法に用いる写像、第三段階は
$W_i\subset V_i\otimes R_{i+1}$ である。
これらが微分と可換であることは、それぞれテンソルの最後の因子を
二通りに取り出して同じ積になることで従う。

従って $F^i$ と $V_i\otimes F^{i+1}$ を項ごとに直和した
$M_i$ の分解では、第三微分の交差項は符号を除いて
$W_i\subset V_i\otimes R_{i+1}$ である。
$\Hom_D(-,P_{i+5})$ を取ると、
第三項は $V_i^*\otimes W_{i+1}^*$ のみで、第二項は零である。
一方、$\Hom_D(-,P_{i+4})$ を取った第三項は
\[
 W_i^*\oplus(V_i^*\otimes W_{i+1}^*\otimes V_{i+4})
\]
であり、第二項 $V_i^*\otimes R_{i+1}^*$ の像で割る。
その第二成分への写像は $w_{i+1}\in R_{i+1}\otimes V_{i+4}$ の縮約なので
同型である。第一成分への写像は $w_i\in V_i\otimes R_{i+1}$ の縮約である。
よって、$a\in V_i$ を用いる $S_{i+1}\to M_i$ での引き戻しと
$\theta_i(a)$ を用いる $P_{i+5}\to P_{i+4}$ との合成は、
双対基底 $w_i^*,w_{i+1}^*$ を使えば、いずれも
$V_i^*$ 成分を $a$ で評価する写像である。交差項の符号は基底変更で吸収できる。

より具体的には $R^3\tau(D)$ の成分ごとの $k$-双対を $L$ とすると、
適切に基底を取り直した $m_i\in L_{i,i+4}$ について
\[
 a m_{i+1}=m_i\theta_i(a)\qquad(a\in V_i)
\]
が成り立つ。右 AS 正則性より $L$ の第 $j$ 列は左加群 $D e_{j-4}$ である
\cite[arXiv v1, Theorem 4.8 and Remark 4.10]{MN2025}。
$D$ は $V_i$ で生成されるため、上の関係が双加群構造を決め、
$L\simeq D(0,-4)_{\theta^{-1}}$ となる。
従って \cite[arXiv v1, Definition 4.15 and Theorem 6.8]{MN2025} の
Nakayama 同型と Serre 関手の式が得られる。
$D_{ij}$ に $c_i/c_j$ を掛ける対角型自己同型による引き戻しは、
第 $j$ 成分を $c_j$ 倍して恒等と自然同型になる。
\end{proof}

\begin{lemma}[整な楕円四つ組の二周期の平方]
\label{lem:elliptic-period-square}
幾何的データが整な算術種数一の曲線 $E$ 上の許容四つ組の場合、
\cite[arXiv v3, Proposition 5.6.1]{VdB2011} の平行移動から生じる
次数 $+2$ の周期移動 $\beta$ は、対角型自己同型を除いて
$\beta^2=\theta$ を満たす。
\end{lemma}

\begin{proof}
滑らかな場合は $E$ に原点を選び、因子和を楕円曲線の群法則で取る。
特異な場合、正規化は種数零で、全 $\delta$ 不変量が一である。
従って正規化は $\PP^1$ で、特異点は一つである。
二枝なら完備局所環は正規化環のうち二つの剰余値が等しい部分環である。
実際、剰余値の一致だけで余次元一となる。
一枝なら一様化変数を $u$ として完備局所環は $k+u^2 k[[u]]$ である。
一次の項が非零の元を含めば、形式的逆写像により正規化環全体を含むことになり、
そのような元を含まない部分環は余次元一である。

二枝の場合は正規化上の二点を $0,\infty$、一枝の場合はその点を $\infty$ とする。
滑らかな部分にそれぞれ乗法群 $k^\times$ と加法群 $k$ の群構造を入れ、
以下では両方を加法記号で書く。この群の平行移動は正規化上の
拡大写像または平行移動で実現され、上記の剰余値または一次項の条件を保つので
$E$ の自己同型となる。滑らかな部分に台を持つ次数零の因子
$\sum_j m_j (p_j)$ が主因子となる条件は、群での和 $\sum_j m_j p_j=0$ である。
実際、正規化での有理関数 $\prod_j(t-p_j)^{m_j}$ は $\infty$ で $1$ となる。
二枝の場合は $0$ での値も $1$ であること、一枝の場合は $1/t$ の一次項が
零であることが特異点で可逆となる条件であり、いずれも上記の群の等式である。
従って滑らかな部分の因子和によって、これらの因子から得た同次数の
直線束の同型を判定できる。

各 $L_i$ を与える線型系の一般のファイバーは、
特異点を避ける相異なる二点からなる。実際、$L_i$ は二重被覆を与え、
標数零なので生成する関数体の拡大は分離的である。
この二点の和を $a_i$ とすると、同じファイバーの相手は
$\iota_i(p)=a_i-p$ と書ける。楕円 helix の関係から
$a_{i+2}-a_i$ は全ての $i$ で一定である。
$2v=a_0-a_2$ を満たす点 $v$ を選べる。群が楕円曲線、
$k^\times$ または $k$ のいずれでも、二倍写像は全射だからである。
この $v$ による平行移動で $L_i$ を引き戻すと $L_{i+2}$ となり、
切断への作用が $\beta$ を与える。

$\theta_i:V_i\to V_{i+4}$ の作用を評価写像で比較する。
一般の $p\in E$ に対し $p'=\iota_{i+2}(p)$ と置き、
$V_j=H^0(E,L_j)$ における $q$ での評価の核を
$k\ell_j(q)$ と書く。
$w_i$ を $L_{i+2}$ の $p$ における評価で縮約する。
$p$ と $p'$ は同じ射影評価を与える。
さらに $L_{i+1}$ の $p$ における評価で縮約すると、
$w_i$ が二つの関係空間に属することから
$\ell_i(p)\otimes\ell_{i+3}(p)$ の非零スカラー倍を得る。
$p'$ でも同様である。非零性は、隣接する二因子の縮約が非零という
幾何的五つ組の条件による \cite[arXiv v3, Definition 5.3.1]{VdB2011}。
$L_{i+1}$ の $p,p'$ での評価は独立なので、
$L_{i+2}$ のみを評価したテンソルは、適当な $\lambda,\gamma\ne0$ に対し
\[
 \lambda\ell_i(p)\otimes\ell_{i+1}(p')\otimes\ell_{i+3}(p)
 +\gamma\ell_i(p')\otimes\ell_{i+1}(p)\otimes\ell_{i+3}(p')
\]
と書ける。$L_{i+3}$ の $p$ における評価で縮約すると、
$L_{i+3}\not\simeq L_{i+2}$ より第二項だけが残る。
$w_{i+2}$ は $w_i$ の二回の回転により与えられ、
回転で末尾へ移る二因子は $\theta_i,\theta_{i+1}$ で移す。
従って $w_{i+2}$ を各 $L_{i+2},L_{i+3},L_{i+4}$ の $p$ における評価で
縮約して零になることから、
$\theta_i(\ell_i(p'))$ は $p$ で消える。ゆえに $\theta_i$ による
$L_{i+4}$ の $p$ での射影評価の引き戻しは $L_i$ の $p'$ での評価である。
$L_i$ に関して $p'$ は
$\iota_i(p')=p+a_i-a_{i+2}=p+2v$ と同じ射影評価を与える。
これは $\beta^2$ の作用である。
全ての $V_i$ で二つの写像はスカラー倍で一致し、
$V_i$ の生成性と添字直線上でのスカラーの逐次選択により
その差は対角型自己同型である。
\end{proof}

\begin{theorem}[整な楕円四つ組からの対象]
\label{thm:integral-elliptic-inverse}
標数零の代数閉体上で、$U=(E,L_0,L_1,L_2)$ を整な算術種数一の曲線上の
Van den Bergh の許容四つ組、$D=D(U)$ とする。
このとき $(\qgr D,\pi P_0)\in\HH$ である。
平方根は VdB の平行移動に由来する偶数添字の周期で与えられ、
二次 AS 正則被覆はその周期で次数付けした標準 hull から得られる。
\end{theorem}

\begin{proof}
$\beta$ を補題~\ref{lem:elliptic-period-square} の周期移動とする。
\cite[arXiv v3, Theorem 6.1, Lemmas 6.4--6.5 and Proposition 6.6]{VdB2011} の
標準 hull はこの周期移動を持ち上げる。
偶数部分を $\beta$ で次数付けした環 $B$ は、
型 $(1,4,6,4,1)$ の四次元二次 AS 正則環 $A$ の
正則正規二次元 $g$ による商である。ここで被覆の標準分解は線型なので
$A$ は Koszul である。左右の AS 条件については
\cite[arXiv v1, Proposition 4.4 and Theorem 4.5]{MN2025} を参照できる。
右 noetherian 性は同じ hull 定理による。
$D$ の反対環は添字を反転すると三次正則型であるため
\cite[arXiv v1, Theorem 4.5]{MN2025}、同じ VdB の分類と
\cite[arXiv v3, Corollary 5.5.9 and Lemma 3.5]{VdB2011} から
$D$ と $B$ は左 noetherian でもある。
$g$ によるフィルトレーションの次数環は $B$ 上の skew polynomial 環なので
$A$ にも左 noetherian 性が昇る。実際、斉次左イデアルについては固定した
内部次数でこのフィルトレーションは有限であり、
初期項の生成元から元のイデアルの生成元を持ち上げられる。
一般の左イデアルについては内部次数による最高次成分を使う。

偶数部分への制限は、指定対象を保つ $\qgr D\simeq\qgr B$ を与える
\cite[arXiv v3, Lemmas 3.4--3.5]{VdB2011}。
$\pi P_{-j}$ の列は ample であり
\cite[arXiv v1, Theorem 6.11]{MN2025}、奇数添字の $\pi P_i$ は
次の偶数添字の $\pi P_{i+1}$ の有限直和の商になる。
これは $D$ の隣接成分による生成性から従う。
従って偶数部分列も ample である。シフト $(1)$ は
$\qgr D$ 側の $(2)_\beta$ であり、
補題~\ref{lem:cubic-serre-period} と補題~\ref{lem:elliptic-period-square}
からその平方は $S^{-1}[2]$ である。
\cite[arXiv v1, Proposition 4.11]{MN2025} により
$\Hom(\pi P_j,\pi P_i)=D_{ij}$ であり、
$\pi P_0$ は exceptional である。この Hom の同一視は合成を保つため、
$\pi P_0$ の切断環は全非負次数で $B$ と同型となる。
実際、次数 $n$ の切断の空間 $D_{-2n,0}$ を $\beta^n$ で $D_{0,2n}$ に移すと、
次数 $m,n$ の積は $a\beta^m(b)$ となり、周期による $B$ の積に一致する。
$\qgr D$ の Ext 有限性は \cite[arXiv v3, Corollary 5.5.9]{VdB2011} による。

最後に rank 条件を確認する。
楕円 helix の切断 $\mathbb Z$-algebra $D_{\mathrm{geom}}$ への商の核は、
次数四の正則な一元列 $(h_i)$ で左右から生成される
\cite[arXiv v3, Theorem 5.5.6]{VdB2011}。
$\beta$ は幾何データ上の平行移動から来るためこの核を保つ。
従って偶数部分の次数付き商では、列 $(h_{2i})$ が一つの
正則正規元 $h\in B_2$ に降りる。
$B/(h)$ は $E$ の次数四の直線束 $L_0\otimes L_1$ と
平行移動による捩れ斉次座標環である。
従ってこれは domain であり、
$h$ による分離フィルトレーションの初期項から $B$ も domain である。
また $\qgr(B/(h))\simeq\coh E$ であり、
その構造対象のシフトは可逆層に対応する
\cite[arXiv v3, Corollary 5.5.9 and Lemma 3.5]{VdB2011}。
$h$ の冪が与える $\pi B$ の下降列の隣接商は、従って有限長でない。
命題~\ref{prop:sklyanin-h2} の証明と同じ Hilbert 多項式の議論を適用すると、
$\pi B$ は Gabriel--Krull の $\cC_1$ に属さず、
任意の非零部分対象は $\pi B$ のシフトを含み、その余商は $\cC_1$ に属する。
ここで使うのは $B$ の Hilbert 級数、domain 性と、
有限生成 $B$ 加群が $A$ 上で有限自由分解を持つことである。
よって $\pi B$ は rank one torsion-free であり、全ての条件が成立する。
\end{proof}

% !TEX root = ../nc-quadric-surface.tex
\section{Sklyanin pencil の一次近傍と同値判定}
\label{sec:pencil-deformation}

命題~\ref{prop:sklyanin-antiserre-center} の無限位数の範囲で、
同一の四次元 Sklyanin 環 $A$ の pencil を比較する。
$P=A^{(2)}$、$Z\subset P_1$ を元の中心二次 pencil、
$D=P/(Z)$ を楕円商と置く。直線 $\ell\subset Z$ に対して $R_\ell=P/(\ell)$ であり、
滑らかな対象に対応する場合にはこれが反 Serre 環である。
$H_\ell=Z/\ell$ の像は $R_\ell$ の中心一次部分である。

\begin{proposition}[中心方向の一次変形]
\label{prop:pencil-tangent}
$U=D_1$、$K=\ker(U\otimes_k U\to D_2)$ とする。
一次補空間を選ぶと、次数二の積により線形写像
\[
 \Delta:Z^*\longrightarrow
 \operatorname{Hom}_k(K,U)/\operatorname{im}\partial,\qquad
 (\partial f)(r)=(f\otimes1+1\otimes f)(r),\quad f\in U^*
\]
が得られ、右辺の商において補空間の選択に依存しない。
$\ell$ を消す非零汎関数 $\lambda$ に対して
$\Delta(\lambda)$ は $R_\ell/(H_\ell^2)$ の変形を与える。
この商の同型は $D$ の次数付き自己同型と $H_\ell$ の基底の変更の下で
$\Delta$ の値を保つ。
\end{proposition}

\begin{proof}
$g,h$ を $Z$ の任意の基底とする。$g$ は $P$ 上で正則であり、
$h$ は $P/(g)$ 上で正則である。これは Sklyanin の中心正則列から従う。
実際、適合する滑らかな一切断 $g$ については domain 性が保証し、
任意の基底については、中心正則列の Koszul 複体が基底変更で同型であることを使える。
後者から任意の非零 $g$ に対する商上の正則性が従う
（$P$ 上の $g$ 自身の正則性は domain 性による）。
この正則列により
\[
 (Z)/(Z)^2 \simeq Z\otimes_k D
\]
となる。例えば $g x+h y\in (Z)^2$ なら右辺の二次式を
$g,h$ の係数に吸収でき、残る関係は $g x'+h y'=0$ となる。
これを $(g)$ で割れば $y'\in(g)$ なので $x'\in(h)$ となる。

$U$ の線形持上げ $U\to P_1$ を選ぶ。
$r\in K$ の積を $P/(Z)^2$ で取ると $Z\otimes U$ に入る。
$\lambda$ で $Z$ を評価して $U$ への写像を得る。
持上げの差 $U\to Z$ は積に $\partial f$ を加え、二つの差の積は零となる。
従って商での類は内在的である。
さらに $(Z)^2+(\ell)$ で割ると $R_\ell/(H_\ell^2)$ となり、
その一次変形が評価した写像で与えられる。
商の同型は $H_\ell$ の同定と $U$ の持上げを移すので、
基底変更と $\partial$ の範囲で作用する。
\end{proof}

\begin{lemma}[持上げ変更の単射性]
\label{lem:pencil-lift-injective}
命題~\ref{prop:sklyanin-antiserre-center} の無限位数の範囲では、
$\partial:U^*\to\operatorname{Hom}_k(K,U)$ は常に単射である。
\end{lemma}
\begin{proof}
$D$ は二次環である。実際、二次環 $A$ の二次 Veronese $P$ の二次性は
命題~\ref{prop:antiserre-quadratic} と同じ関係の挿入による議論から従い、
一次空間 $Z$ で割っても二次性は保たれる。
$\partial f=0$ となる非零 $f\in U^*$ が存在すると仮定する。
各 $r\in K$ に対して
\[
 0=f((\partial f)(r))=2(f\otimes f)(r)
\]
である。標数零と二次性から、$u\mapsto f(u)z$ は次数付き全射
$D\twoheadrightarrow k[z]$ を定める。

この全射による右 $D$-加群を $M=k[z]$ と書く。
非零次数付き部分加群は $z^n k[z]$ なので、$\pi M$ は $\qgr D$ の単純対象である。
また $z$ 倍写像 $M\to M(1)$ の余核は有限次元であり、
$\pi M\simeq(\pi M)(1)$ となる。
命題~\ref{prop:sklyanin-antiserre-center} の楕円商の同定
$\qgr D\simeq\coh E$ の下で、単純対象は閉点 $p\in E$ の skyscraper 層である。
次数シフトは台を $2\tau$ または $-2\tau$ だけ平行移動する。
従って上の同型は $p$ がこの平行移動の固定点であることを意味するが、
$2\tau\ne0$ なので矛盾する。
\end{proof}

\begin{proposition}[固定 pencil の比較を有限化する条件]
\label{prop:pencil-finite-candidates}
$\Delta$ が単射であるとする。
固定 $P$ の滑らかな二切断の pointed exact 同値を判定するには、
有限個の一次線形写像だけを二次関係に照合すればよい。
さらに、各切断と同値な固定 pencil 内の切断は有限個である。
\end{proposition}

\begin{proof}
補題~\ref{lem:pencil-lift-injective} により $\partial$ も単射である。
まず $D$ の次数付き自己同型群を次数スカラーで割ると有限になる。
$D$ は次数八の可逆層を持つ楕円曲線の twisted homogeneous
coordinate ring（またはその反対環）である。
次数付き自己同型は次数零の分数体 $k(E)$ に作用する。
次数一切断の比で与えられる完全線形系の埋込みにより、
この作用は次数八の偏極を保存する $E$ の自己同型である。
詳しくは、基底切断 $u_0$ で割った切断 $u_i/u_0$ を座標比とすると、
代数自己同型が送る比は新基底の座標比であり、埋込みは射影線形変換を介して保たれる。
同じ分数体作用を持つ二自己同型の差は、比を固定するので
$U$ に共通の有理関数を掛ける。
この有理関数は非零有限次元空間を保つため、固有ベクトルに作用させれば
スカラーと分かる。$D$ は一次生成なので、差は次数スカラーである。
次数一の斉次可逆元による分数体上の作用は二倍平行移動であるため、
分数体上の自己同型はこれと可換である。
平行移動を除いた楕円群の自己同型が無限位数の移動点を固定すれば恒等になる。
実際、その自己同型から恒等を引いた非零準同型の核は有限である。
従って残る作用は次数八の偏極を保つ平行移動であり、$E[8]$ に属する。
特に候補は高々64個である。

比較する二切断の $H_\ell$ に基底を選ぶ。
$D$ 上の自己同型の代表 $\varphi$ ごとに
命題~\ref{prop:pencil-tangent} の必要条件を調べる。
二つの変形類は非零なので、それらの比較に必要な
$H_\ell$ 上の倍率は存在するなら一意である。
代表写像の差が $\partial f$ であることを確認すれば、
$\partial$ の単射性により $f$ も一意に定まる。
この $f$ は次数一における持上げのずれである。
従って商 $R_\ell/(H_\ell^2)$ 上の候補から、
$R_\ell$ の次数一の線形写像が一意に決まる。
次数スカラーは常に全環の自己同型なので、$\varphi$ は各剰余類で一つ取ればよい。
最後にその線形写像が全二次関係を保つか検査する。
命題~\ref{prop:antiserre-quadratic} と定理~\ref{thm:antiserre-invariant}
により、この最後の条件は十分でもある。

固定した切断の非零変形類の射影化に $\varphi$ を作用させると、
$\Delta$ の単射性により各 $\varphi$ につき一致する切断は高々一つである。
\end{proof}

\begin{proposition}[通常の Sklyanin パラメータでの開条件]
\label{prop:pencil-tangent-generic}
通常の Sklyanin のパラメータ平面のチャート
$\gamma=-(\alpha+\beta)/(1+\alpha\beta)$ において、
$\partial$ と $\Delta$ がともに単射となる条件は非空 Zariski 開条件を与える。
無限位数の範囲では $\partial$ の単射性は補題~\ref{lem:pencil-lift-injective}
により自動的であり、有限比較に残る仮定は $\Delta$ の単射性だけである。
従ってこの開条件を満たし、かつ命題~\ref{prop:sklyanin-antiserre-center} の仮定を
満たす固定 pencil について、滑らかな切断の同値判定は命題~\ref{prop:pencil-finite-candidates} の有限手順で行える。
\end{proposition}

\begin{proof}
$\dim U=8$、$\dim D_2=16$ なので $K$ は48次元である。
一次持上げと関係空間の基底を替えても、
$\partial$ の階数と、$\partial$ の像に中心二方向の列を加えた行列の階数は不変である。
二つの単射性はこれらがそれぞれ $8,10$ であることと同値であり、
定ランクの座標チャートでは極大小行列式の非消滅で書ける。

非空性を有理数上で確認した。
$(\alpha,\beta,\gamma)=(2,3,-5/7)$ において、
$\sum a_i x_i^2$ の四生成元との交換子を次数三の関係空間で割って
二次元の中心 pencil の基底を求める。
次数四では、二次関係を四文字の三箇所へ挿入した空間で割る。
得られる35次元空間に、中心の三つの二乗積と
二つの中心元と八つの $U$ の持上げとの積を並べると、階数は19となる。
さらにその像で割ると16次元になり、
$U\otimes U$ からの積の核が48次元と求まる。
この核に対し、中心方向の積の係数と $\partial$ の列を並べると
有理数上の階数がそれぞれ $8,10$ となる。
計算の入力と行操作は
\texttt{computations/verify\_pencil\_tangent.py} に収めた。
この計算は開条件の非空性の証明に用い、当該有理パラメータの
平行移動の無限位数性は用いていない。
\end{proof}



\section{同値類の分類後に残る問題}
\label{sec:open-problems}

定理~\ref{thm:complete-h2-classification} は、体上の全対象を分類し、
構成するデータと必要十分な同値判定を与える。
以下では、この結論に含めない構造を区別する。

\begin{problem}[幾何データ上での同値判定の簡約]
全許容四つ組について成立する有限判定を、曲線・直線束の変換による
より小さい正規形や幾何的な対応で記述せよ。
特殊な自己同型を持つ点と全ての退化型を含め、既存の tensor の GIT 商と比較せよ。
この簡約を待たずに、主定理の有限判定は全同値類を区別する。
\end{problem}

\begin{problem}[平方根と表示の忘却]
固定 $(\cC,\OO)$ 上の適合する平方根を分類し、その自己同値群の作用を求めよ。
さらに固定平方根に対する ambient 表示を分類し、
これらの二段階の忘却が既存 tensor モジュライとどう関係するかを記述せよ。
可換点の四つの平方根や固定切断環の非一意な ambient を、
全非可換点での有限性・一般次数の結論に外挿してはならない。
\end{problem}

\begin{remark}[全体の Deligne--Mumford 性]
可換対象 $(\coh Q,\mathcal O_Q)$ における pointed 自己同値群は
\[
 \operatorname{Aut}(Q)=(\operatorname{PGL}_2\times\operatorname{PGL}_2)\rtimes S_2
\]
である。実際、\cite[Theorem 3.1]{BondalOrlov2001} による幾何自己同型、
線束、cohomological shift の記述で、exact 性と $\mathcal O_Q$ の保存により
最後の二つの因子が消える。幾何自己同型は族でも作用するので、
この自己同値を通常の慣性として保持する全体のモジュライは、
可換点で六次元の慣性を持つ。従って目標とする代数スタックは全体では
Deligne--Mumford スタックにはならない。
\end{remark}

\begin{problem}[境界の相対表示と framed 局所構造]
定理~\cite[\texttt{thm:global-h2-smooth-finite-type}]{QuadricModuli2026} により $\HHstack$ の有限型代数性は成立し、
系~\cite[\texttt{cor:global-xas-smooth-scheme}]{QuadricModuli2026} により framed 実現空間は smooth スキームとなる。
有限 foundation の ruling 商 $\mathfrak Q_{\mathrm{fnd}}$ との相対比較を具体化し、
境界で確立した相対商表示（定理~\cite[\texttt{thm:noncubic-intrinsic-presentation-equivalence}]{QuadricModuli2026}）を踏まえて、
全体の相対 groupoid を比較せよ。
全体の既約性と被約境界の滑らかさ・余次元二は、
命題~\cite[\texttt{prop:global-h2-irreducible}]{QuadricModuli2026} と定理~\cite[\texttt{thm:noncubic-boundary-smooth-codim-two}]{QuadricModuli2026} で成立した。
既存の smooth proper dg 特殊化族そのものの相対 $\HHstack$ 認識は、各族について別に確認する。
この内在的な既約性の証明は、実際の相対族の連続性を使い、
命題~\cite[\texttt{prop:h2-component-density}]{QuadricModuli2026} の誘導位相との同一視を必要としない。

また $X_{\mathrm{as}}\to X_3$ の局所閉性と具体的な認識方程式を調べよ。
一致根対では命題~\cite[\texttt{prop:global-h2-not-ambient-open}]{QuadricModuli2026} が開性を否定する。
余分な一次方向は系~\cite[\texttt{cor:artin-polarization-retraction}]{QuadricModuli2026} の
$\operatorname{HH}^1$ に対応する偏極変形である。
全スキーム上の商による延長と、任意 Noetherian 基底上で同じ Abel 公理を課す関手との一致も、
極限保存と canonical 導来圏の比較を伴う別の課題である。
\end{problem}

\bibliographystyle{alpha}
\bibliography{references}
\end{document}
